\chapter{Spectral Algebraic Geometry: Theories, Contexts, and Comprehension}
\label{chap:spectral-algebraic-geometry}

\section{Orientation: theories, models, contexts, and comprehension}
\label{sec:spectral-ag-purpose}

Geometry and logic are both organized by change of context. In geometry,
objects and coefficients are transported along maps of bases, local data are
reconstructed by descent, and relative algebra can determine new spaces over a
given base. In categorical semantics, the corresponding operations appear as
substitution and context extension. The question of this chapter is whether
these operations can be organized together in spectral algebraic geometry:
can geometric change of base, stable coefficients, locality, and context
extension be made into one coherent indexed geometry? We use the words
\textbf{context}, \textbf{substitution}, and \textbf{comprehension} for these
roles. The terminology is drawn from indexed categories in categorical logic
and semantics, while the constructions themselves are spectral-geometric;
compare \cite{JacobsCategoricalLogic}.

The obstruction is that the geometric and linear layers have different
homotopical character. Affine coordinates are connective commutative ring
spectra, whereas the natural coefficient categories are stable module
categories. Mapping spectra, shifts, cofibres, cotangent complexes, and other
linearized constructions require the full stable layer. The chapter therefore
keeps connective geometry and fully stable coefficient theory distinct while
making their changes of context coherent and their local data reconstructible.

The core construction has three stages. First, connective commutative ring
spectra are obtained as models of the commutative operadic theory and then read
contravariantly as affine spectral contexts. In symbols,
\[
    \Comm^\otimes
    \longrightarrow
    \Sp_{\geq0}^\otimes
    \longrightarrow
    \Alg_{\Comm^\otimes}(\Sp_{\geq0}^\otimes)
    \simeq
    \CAlg(\Sp)^{\mathrm{cn}}
    \longrightarrow
    \Aff_{\Sp}
    :=
    \bigl(\CAlg(\Sp)^{\mathrm{cn}}\bigr)^{\op}.
\]
Here \(\Comm^\otimes\) is the commutative \(\infty\)-operad and
\(\Sp_{\geq0}^\otimes\) is the symmetric-monoidal universe in which its
coordinate models are interpreted. The word \textbf{theory} is used in this
operadic sense: this is the symmetric-monoidal analogue of interpreting an
algebraic theory in a semantic universe, not an identification of an
\(\infty\)-operad with a finite-product Lawvere theory. Thus the first stage
is the pattern \textbf{theory \(\to\) models \(\to\) affines}, specialized to
connective spectral algebra; compare
\cite{LurieHigherAlgebra,LurieSAG,KhanBraveNewMotivicI}.

Second, each affine context is equipped with a stable coefficient fibre. If
\(X=\Spec(A)\), then
\[
    \QCoh_{\mathrm{aff}}(X)
    :=
    \Mod_A,
\]
and a context morphism
\[
    f:\Spec(B)\longrightarrow\Spec(A)
\]
represented by \(A\to B\) acts on coefficients by extension of scalars,
\[
    f^*
    =
    -\otimes_AB
    :
    \Mod_A
    \longrightarrow
    \Mod_B.
\]
We read this pullback geometrically as substitution. The module construction
makes these substitutions functorial and strong symmetric monoidal, with
restriction of scalars as right adjoint. The affine coefficient system is then
extended to generalized contexts and, after Zariski descent, restricts to
\[
    \QCoh
    :
    \bigl(\operatorname{Sch}^{\mathrm{sp}}\bigr)^{\op}
    \longrightarrow
    \CAlg(\Pr^L_{\mathrm{st}}),
    \qquad
    X\longmapsto\QCoh(X).
\]
The detailed Kan-extension and adjoint formalism will be introduced only when
it becomes mathematically active. At this stage the point is that stable
coefficients vary coherently with geometric context.

Third, locality selects the geometric contexts and relative algebra returns
coefficient data to geometry. Spectral localizations generate the Zariski
topology, quasi-coherent coefficients satisfy descent, and affine contexts
glue to spectral schemes. Over a spectral scheme \(X\), relative spectrum
then identifies connective quasi-coherent commutative algebras with affine
morphisms over \(X\):
\[
    \operatorname{QCAlg}(X)^{\mathrm{cn},\op}
    \xrightarrow{\ \simeq\ }
    \bigl(\operatorname{Sch}^{\mathrm{sp}}_{/X}\bigr)_{\mathrm{aff}}.
\]
Its universal property identifies maps into the corresponding affine extension
with algebra terms after substitution. This is the comprehension mechanism
constructed in the chapter: suitable algebraic data over a context are
represented by a geometric extension of that context, compatibly with change
of base.

These three stages form the structural spine of the chapter:
\[
    \boxed{
    \text{theory}
    \to
    \text{models}
    \to
    \text{contexts}
    \to
    \text{substitution}
    \to
    \text{descent}
    \to
    \text{comprehension}.}
\]
The distinction between connective coordinates and stable coefficients remains
active throughout. Coordinate algebras lie in
\(\CAlg(\Sp)^{\mathrm{cn}}\), while their coefficient categories retain the
full stable structure required by linear and infinitesimal constructions.
Once this indexed geometry has been established, the later sections develop
its consequences: closed subcontexts and gluing, linear and section
representations, distinguished affine contexts, and the recovery of the
ordinary \(0\)-truncated sector. The purpose of the intervening theory is to
make all of these constructions interact without conflating the connective
and stable levels.

The foundational spectral-algebraic and spectral-geometric results used in
this construction are drawn from
\cite{LurieHigherAlgebra,LurieSAG,LurieDAGIV,KhanBraveNewMotivicI}. The
following conventions fix the technical background once, so that the main
argument can subsequently refer to it without reopening coherence, size, or
variance issues.
\begin{convention}
\label{conv:spectral-infinity-categorical}
All categories are $\infty$-categories unless explicitly stated otherwise.
Mapping spaces are denoted by
\[
    \Map_{\mathcal C}(X,Y).
\]
Mapping spectra in a stable $\infty$-category are denoted by
\[
    \map_{\mathcal C}(X,Y).
\]
Equivalences are denoted by
\[
    \simeq.
\]
\end{convention}

\begin{convention}[Higher coherence]
\label{conv:spectral-higher-coherence}
Whenever a structure is produced by a functor, an adjunction, coCartesian
transport, an operadic construction, a Kan extension, a limit, or a descent
totalization, we use the coherent higher structure supplied by that
construction. No independent coherence conditions are imposed. After the
source of the coherent structure has been specified, its higher homotopies are
suppressed from the notation.
\end{convention}

\begin{convention}[Constructed structure and proof witnesses]
\label{conv:spectral-constructed-structure}
Choices of representatives, atlases, trivializations, cell expressions, and
presentations made inside proofs are not retained as additional structure
unless explicitly stated. Likewise, comparison, exchange, base-change,
projection, and coherence morphisms are those constructed from the displayed
functors, adjunctions, Kan extensions, or descent limits. Their invertibility
is a theorem to be proved or cited where it is used. Terminology belonging to
an external structured-geometric presentation is not transported into the
native category unless it is explicitly defined here.
\end{convention}

\begin{convention}[Size]
\label{conv:size-spectral-ag}
Fix Grothendieck universes
\[
    \mathbb U
    \subset
    \mathbb V.
\]
Spectra, ring spectra, modules, and geometric objects are taken to have
\(\mathbb U\)-small underlying data. The \(\infty\)-categories collecting all
such objects are regarded at the \(\mathbb V\)-level, and presheaf, sheaf,
functor, and presentable-category constructions are formed in that enlarged
ambient universe. Enlargement functors are suppressed when they induce the
canonical fully faithful identification of the construction in question.
\end{convention}

\begin{convention}[Monoidal structures on categories]
\label{conv:spectral-monoidal-structures-on-categories}
The notation
\[
    \CAlg(\Cat_\infty)
\]
uses the Cartesian symmetric monoidal structure on $\Cat_\infty$.
The notation
\[
    \CAlg(\Pr^L_{\mathrm{st}})
\]
uses the tensor product of presentable stable $\infty$-categories. Its
morphisms are colimit-preserving strong symmetric monoidal functors. We write
\[
    \Pr^R_{\mathrm{st}}
\]
for presentable stable $\infty$-categories with accessible right-adjoint
functors. Passage to right adjoints gives the canonical equivalence
\[
\begin{tikzcd}[column sep=huge]
    (\Pr^L_{\mathrm{st}})^{\op}
        \arrow[r, "{\operatorname{RAdj}}", "\simeq"']
    &
    \Pr^R_{\mathrm{st}}.
\end{tikzcd}
\]
Its functoriality supplies the coherent composition of right adjoints used
below.
\end{convention}

\section{Spectral semantic universes}
\label{sec:spectra-postnikov-structure}

The construction begins by separating the two spectral roles identified in
the orientation. The full category \(\Sp\) supplies stable coefficients and
linear constructions, whereas \(\Sp_{\geq0}\) supplies the
symmetric-monoidal universe for connective coordinates. This section records
only the structural features of spectra needed to keep those roles distinct
and to compare them through the Postnikov \(t\)-structure.

\subsection{The stable symmetric monoidal category of spectra}

\begin{proposition}[Spectral semantic category]
\label{prop:spectra-presentably-stable-monoidal}
The $\infty$-category of spectra admits a canonical structure
\[
    \Sp^\otimes
    \in
    \CAlg(\Pr^L_{\mathrm{st}}).
\]
Its tensor product is the smash product
\[
    -\otimes-
    \colon
    \Sp\times\Sp
    \longrightarrow
    \Sp,
\]
its tensor unit is the sphere spectrum $\mathbb S$, and its tensor product
preserves small colimits separately in each variable.
\end{proposition}

\begin{proof}
\leavevmode
\begin{enumerate}[label=\textup{(\arabic*)}]
    \item \textbf{Stabilize pointed spaces.}
    The stabilization of pointed spaces is a presentable stable
    $\infty$-category.

    \item \textbf{Stabilize the smash product.}
    The smash product on pointed spaces stabilizes to a
    symmetric monoidal product on spectra.

    \item \textbf{Identify the tensor unit.}
    The stabilization universal property identifies its tensor unit with the sphere spectrum.

    \item \textbf{Preserve colimits separately.}
    Separate preservation of colimits follows from the presentably symmetric monoidal
    construction; see \cite{LurieHigherAlgebra}.
\end{enumerate}
\end{proof}

\begin{proposition}[Closedness of spectra]
\label{prop:spectra-closedness}
For every spectrum $X$, the functor
\[
    -\otimes X
    \colon
    \Sp
    \longrightarrow
    \Sp
\]
admits a right adjoint
\[
    \HHom_{\Sp}(X,-)
    \colon
    \Sp
    \longrightarrow
    \Sp.
\]
Thus there are natural equivalences
\[
    \Map_{\Sp}(Z\otimes X,Y)
    \simeq
    \Map_{\Sp}
    \bigl(
        Z,
        \HHom_{\Sp}(X,Y)
    \bigr).
\]
\end{proposition}

\begin{proof}
\leavevmode
\begin{enumerate}[label=\textup{(\arabic*)}]
    \item \textbf{Preserve small colimits.}
    The functor $-\otimes X$ preserves small colimits by
    \Ref{prop:spectra-presentably-stable-monoidal}.

    \item \textbf{Apply the adjoint functor theorem.}
    Since $\Sp$ is presentable, the adjoint functor theorem supplies the right adjoint.
\end{enumerate}
\end{proof}

The closed structure will later be used through its adjunction, which we fix
in the following form:
\[
\begin{tikzcd}[column sep=huge]
    \Sp
        \arrow[r, shift left=1.1ex, "-\otimes X"]
    &
    \Sp
        \arrow[l, swap, shift left=1.1ex, "{\HHom_{\Sp}(X,-)}"'].
\end{tikzcd}
\]

\subsection{The canonical \texorpdfstring{$t$}{t}-structure}

\begin{definition}[Postnikov $t$-structure]
\label{def:postnikov-t-structure}
For every integer $n$, let
\[
    \Sp_{\geq n}
\]
be the full subcategory of spectra $X$ for which
\[
    \pi_i(X)=0
    \qquad
    (i<n),
\]
and let
\[
    \Sp_{\leq n}
\]
be the full subcategory of spectra $X$ for which
\[
    \pi_i(X)=0
    \qquad
    (i>n).
\]
The pair
\[
    (\Sp_{\geq0},\Sp_{\leq0})
\]
is the \textbf{canonical Postnikov $t$-structure} on $\Sp$. Its truncation
functors are denoted
\[
    \tau_{\geq n}^{\Sp}
    \colon
    \Sp
    \longrightarrow
    \Sp_{\geq n}
\]
and
\[
    \tau_{\leq n}^{\Sp}
    \colon
    \Sp
    \longrightarrow
    \Sp_{\leq n}.
\]
\end{definition}

\begin{proposition}[Postnikov truncation triangle]
\label{prop:postnikov-truncation-triangle}
For every spectrum $X$ and every integer $n$, there is a natural fibre and
cofibre sequence
\[
    \tau_{\geq n+1}^{\Sp}X
    \longrightarrow
    X
    \longrightarrow
    \tau_{\leq n}^{\Sp}X.
\]
The functors $\tau_{\geq n}^{\Sp}$ and $\tau_{\leq n}^{\Sp}$ are respectively right
and left adjoint to the inclusions of the corresponding full
subcategories.
\end{proposition}

\begin{proof}
\leavevmode
\begin{enumerate}[label=\textup{(\arabic*)}]
    \item \textbf{Identify the truncation sequence.}
    This is the defining truncation sequence of the Postnikov $t$-structure.

    \item \textbf{Use stability.}
    The equality of fibre and cofibre sequences follows from stability.
\end{enumerate}
\end{proof}

\begin{proposition}[Completeness and detection]
\label{prop:postnikov-completeness-detection}
The Postnikov $t$-structure is accessible, compatible with filtered
colimits, left complete, and right complete. In particular, there are
natural equivalences
\[
    X
    \simeq
    \varprojlim\nolimits_n\tau_{\leq n}^{\Sp}X
\]
and
\[
    X
    \simeq
    \varinjlim_n\tau_{\geq -n}^{\Sp}X.
\]
A morphism of spectra is an equivalence if and only if it induces an
isomorphism on every homotopy group.
\end{proposition}

\begin{proof}
\leavevmode
\begin{enumerate}[label=\textup{(\arabic*)}]
    \item \textbf{Invoke the Postnikov completeness theorem.}
    These are the established completeness and accessibility properties of the
    Postnikov $t$-structure; see \cite{LurieHigherAlgebra}.
\end{enumerate}
\end{proof}

\begin{definition}[Discrete spectrum]
\label{def:discrete-spectrum}
A spectrum $X$ is \textbf{discrete} if
\[
    X\in\Sp_{\geq0}\cap\Sp_{\leq0}.
\]
The heart of the Postnikov $t$-structure is denoted
\[
    \Sp^\heartsuit.
\]
\end{definition}

\begin{proposition}[Eilenberg--Mac Lane equivalence]
\label{prop:eilenberg-mac-lane-heart}
The Eilenberg--Mac Lane functor induces an equivalence
\[
    H_{\mathrm{Ab}}
    \colon
    \operatorname{Ab}
    \xrightarrow{\simeq}
    \Sp^\heartsuit.
\]
Under this equivalence, the inverse functor is $\pi_0$.
\end{proposition}

\begin{proof}
\leavevmode
\begin{enumerate}[label=\textup{(\arabic*)}]
    \item \textbf{Recover a discrete spectrum from degree zero.}
    A discrete spectrum is determined by its degree-zero homotopy group.

    \item \textbf{Realize every abelian group.}
    Every abelian group is realized by an Eilenberg--Mac Lane spectrum.
\end{enumerate}
\end{proof}

\subsection{Compatibility with the smash product}

\begin{proposition}[Connectivity of the smash product]
\label{prop:smash-product-connectivity}
If
\[
    X\in\Sp_{\geq m}
    \qquad\text{and}\qquad
    Y\in\Sp_{\geq n},
\]
then
\[
    X\otimes Y
    \in
    \Sp_{\geq m+n}.
\]
In particular, $\Sp_{\geq0}$ is closed under the smash product and contains
the tensor unit $\mathbb S$.
\end{proposition}

\begin{proof}
\leavevmode
\begin{enumerate}[label=\textup{(\arabic*)}]
    \item \textbf{Reduce to the connective case.}
    After suspension, it is enough to consider $m=n=0$.

    \item \textbf{Choose connective generators.}
    Connective spectra are
    generated under colimits and extensions by suspension spectra of pointed
    spaces in nonnegative degrees.

    \item \textbf{Tensor the generators.}
    The smash product preserves colimits
    separately and carries the corresponding generators to connective
    spectra.

    \item \textbf{Invoke monoidal compatibility.}
    Equivalently, the compatibility follows from the monoidal
    compatibility of the Postnikov $t$-structure; see
    \cite{LurieHigherAlgebra}.
\end{enumerate}
\end{proof}

\begin{corollary}[Right $t$-exactness of connective tensoring]
\label{cor:connective-tensor-right-t-exact}
For every connective spectrum $X$, the functor
\[
    -\otimes X
    \colon
    \Sp
    \longrightarrow
    \Sp
\]
is right $t$-exact.
\end{corollary}

\begin{proof}
\leavevmode
\begin{enumerate}[label=\textup{(\arabic*)}]
    \item \textbf{Apply connectivity of the smash product.}
    If $Y$ is connective, then $Y\otimes X$ is connective by
    \Ref{prop:smash-product-connectivity}.
\end{enumerate}
\end{proof}

\subsection{The connective coordinate universe}

\begin{proposition}[Connective spectral semantic universe]
\label{prop:connective-spectral-semantic-universe}
The full subcategory
\[
    \Sp_{\geq0}
    \subseteq
    \Sp
\]
inherits a presentably symmetric monoidal structure
\(\Sp_{\geq0}^\otimes\). Its tensor product is the restriction of the
smash product, its tensor unit is \(\mathbb S\), and the inclusion
\[
    \Sp_{\geq0}^\otimes
    \hookrightarrow
    \Sp^\otimes
\]
is strong symmetric monoidal. Tensoring in \(\Sp_{\geq0}\) preserves small
colimits separately in each variable.
\end{proposition}

\begin{proof}
\leavevmode
\begin{enumerate}[label=\textup{(\arabic*)}]
    \item \textbf{Closure under the monoidal product.}
    The tensor unit is connective, and
    \Ref{prop:smash-product-connectivity} shows that the smash product of
    connective spectra is connective. Hence the monoidal structure on
    \(\Sp\) restricts to \(\Sp_{\geq0}\).

    \item \textbf{Presentability.}
    The connective aisle of the accessible Postnikov \(t\)-structure is a
    presentable full subcategory closed under small colimits.

    \item \textbf{Compute connective colimits.}
    Those
    colimits are computed in \(\Sp\).

    \item \textbf{Preserve colimits under tensoring.}
    Since the smash product on \(\Sp\)
    preserves colimits separately, its restriction has the same property on
    connective spectra. See \cite{LurieHigherAlgebra}.
\end{enumerate}
\end{proof}

The preceding results isolate the two universes needed later:
\(\Sp_{\geq0}^\otimes\) for connective coordinates and \(\Sp^\otimes\) for
stable coefficients. The next section uses the connective universe to form
the coordinate models.

\section{The commutative \texorpdfstring{$\infty$}{infinity}-operadic theory and its spectral models}
\label{sec:commutative-operadic-spectral-algebra}

With the two spectral roles separated, the next step is to construct the
coordinate models. Interpreting the commutative \(\infty\)-operad in
\(\Sp^\otimes\) gives commutative ring spectra; restricting the semantic
universe to \(\Sp_{\geq0}^\otimes\) selects the connective models that will
become affine contexts. The results below record exactly the free,
reflective, discrete, and Postnikov structures used in that passage.

Let
\[
    \Comm^\otimes
    \longrightarrow
    N(\mathrm{Fin}_*)
\]
denote the commutative $\infty$-operad.

\begin{definition}[Commutative ring spectrum]
\label{def:commutative-ring-spectrum}
A \textbf{commutative ring spectrum} is an algebra over the commutative
$\infty$-operad in spectra:
\[
    A
    \in
    \Alg_{\Comm^\otimes}(\Sp^\otimes).
\]
The resulting $\infty$-category is denoted
\[
    \CAlg(\Sp)
    :=
    \Alg_{\Comm^\otimes}(\Sp^\otimes).
\]
Its objects are also called $\EE_\infty$-rings.
\end{definition}

\begin{remark}[Symmetric monoidal envelope]
\label{rem:spectral-commutative-envelope}
The operadic definition above has an equivalent symmetric-monoidal
presentation. Let
\[
    \operatorname{Env}(\Comm)^\otimes
\]
denote the symmetric monoidal envelope of the commutative
$\infty$-operad. Its universal property gives a natural equivalence
\[
    \CAlg(\Sp)
    \simeq
    \Fun^\otimes
    \bigl(
        \operatorname{Env}(\Comm)^\otimes,
        \Sp^\otimes
    \bigr).
\]
Thus a commutative ring spectrum may be read either as an operadic algebra or
as a symmetric-monoidal interpretation of the envelope; in either form, the
commutative operations and their higher coherences are retained. See
\cite[Proposition~2.2.4.9]{LurieHigherAlgebra}.
\end{remark}

For later comparison, the two presentations are related by the canonical
equivalence
\[
\begin{tikzcd}[column sep=huge]
    \Alg_{\Comm^\otimes}(\Sp^\otimes)
        \arrow[r, "\simeq"]
    &
    \Fun^\otimes
    \bigl(
        \operatorname{Env}(\Comm)^\otimes,
        \Sp^\otimes
    \bigr).
\end{tikzcd}
\]

\begin{remark}[Theory, semantic universe, and relative affine geometry]
\label{rem:theory-semantic-universe-relative-geometry}
The preceding equivalence isolates the semantic role of the commutative
\(\infty\)-operad. A commutative ring spectrum is an interpretation of the
commutative operadic theory in a symmetric-monoidal spectral universe. In
the connective coordinate universe, the theory-to-model passage is
\[
    \Comm^\otimes
    \longrightarrow
    \Sp_{\geq0}^\otimes
    \longrightarrow
    \Alg_{\Comm^\otimes}(\Sp_{\geq0}^\otimes)
    \simeq
    \CAlg(\Sp)^{\mathrm{cn}}.
\]
The variance reversal of \Ref{def:affine-spectral-scheme} then reads these
models as affine contexts. This theory--models--affines pattern is the
absolute instance of the relative-algebraic-geometric construction that will
reappear inside \(\QCoh(X)\) under relative spectrum. See
\cite{LurieHigherAlgebra,LurieSAG,KhanBraveNewMotivicI}.
\end{remark}

\begin{proposition}[Presentability of commutative ring spectra]
\label{prop:calg-spectra-presentable}
The $\infty$-category
\[
    \CAlg(\Sp)
\]
is presentable. Its forgetful functor
\[
    \operatorname{oblv}_{\CAlg}
    \colon
    \CAlg(\Sp)
    \longrightarrow
    \Sp
\]
is accessible, creates limits and sifted colimits, and admits a left
adjoint.
\end{proposition}

\begin{proof}
\leavevmode
\begin{enumerate}[label=\textup{(\arabic*)}]
    \item \textbf{Apply presentability for operadic algebras.}
    This is the presentability theorem for algebras over an accessible
    $\infty$-operad in a presentably symmetric monoidal $\infty$-category
    whose tensor product preserves small colimits separately; see
    \cite{LurieHigherAlgebra}.
\end{enumerate}
\end{proof}

\subsection{Free commutative ring spectra}

\begin{definition}[Free commutative algebra]
\label{def:free-commutative-spectral-algebra}
The left adjoint to the forgetful functor is denoted
\[
    \operatorname{Sym}
    \colon
    \Sp
    \longrightarrow
    \CAlg(\Sp).
\]
It is the \textbf{free commutative algebra functor} associated to the
commutative $\infty$-operad.
\end{definition}

The free--forgetful adjunction will be used repeatedly:
\[
\begin{tikzcd}[column sep=large]
    \Sp
        \arrow[r, shift left=1.1ex, "\operatorname{Sym}"]
    &
    \CAlg(\Sp)
        \arrow[l, swap, shift left=1.1ex, "\operatorname{oblv}_{\CAlg}"'].
\end{tikzcd}
\]

\begin{proposition}[Underlying free algebra]
\label{prop:underlying-free-spectral-algebra}
For every spectrum $M$, there is a natural equivalence of underlying
spectra
\[
    \operatorname{oblv}_{\CAlg}\operatorname{Sym}(M)
    \simeq
    \bigoplus_{n\geq0}
    \left(M^{\otimes n}\right)_{h\Sigma_n},
\]
where the summand for $n=0$ is the sphere spectrum.
\end{proposition}

\begin{proof}
\leavevmode
\begin{enumerate}[label=\textup{(\arabic*)}]
    \item \textbf{Apply the free-algebra formula.}
    This is the explicit free-algebra formula for the commutative
    $\infty$-operad in a presentably symmetric monoidal $\infty$-category; see
    \cite{LurieHigherAlgebra}.
\end{enumerate}
\end{proof}

\begin{corollary}[Connectivity of free algebras]
\label{cor:free-algebra-connective}
If $M$ is connective, then $\operatorname{Sym}(M)$ is connective.
\end{corollary}

\begin{proof}
\leavevmode
\begin{enumerate}[label=\textup{(\arabic*)}]
    \item \textbf{Check tensor powers.}
    Every tensor power $M^{\otimes n}$ is connective by
    \Ref{prop:smash-product-connectivity}.

    \item \textbf{Check the remaining colimits.}
    Homotopy orbits and small coproducts
    of connective spectra are connective.

    \item \textbf{Apply the free-algebra formula.}
    Apply
    \Ref{prop:underlying-free-spectral-algebra}.
\end{enumerate}
\end{proof}

\subsection{Connective commutative ring spectra}

\begin{definition}[Connective $\EE_\infty$-ring]
\label{def:connective-e-infinity-ring}
A commutative ring spectrum $A$ is \textbf{connective} if its underlying
spectrum is connective. The full subcategory of connective commutative
ring spectra is denoted
\[
    \CAlg(\Sp)^{\mathrm{cn}}
    \subseteq
    \CAlg(\Sp).
\]
\end{definition}

\begin{proposition}[Connective models as algebras in connective spectra]
\label{prop:connective-calg-operadic-universe}
The strong symmetric monoidal inclusion
\[
    \Sp_{\geq0}^\otimes
    \hookrightarrow
    \Sp^\otimes
\]
induces a fully faithful functor
\[
    \Alg_{\Comm^\otimes}(\Sp_{\geq0}^\otimes)
    \longrightarrow
    \Alg_{\Comm^\otimes}(\Sp^\otimes)
    =
    \CAlg(\Sp),
\]
whose essential image is \(\CAlg(\Sp)^{\mathrm{cn}}\). Hence there is a
canonical equivalence
\[
    \Alg_{\Comm^\otimes}(\Sp_{\geq0}^\otimes)
    \xrightarrow{\simeq}
    \CAlg(\Sp)^{\mathrm{cn}}.
\]
\end{proposition}

\begin{proof}
\leavevmode
\begin{enumerate}[label=\textup{(\arabic*)}]
    \item \textbf{Restrict the operadic structure.}
    By \Ref{prop:connective-spectral-semantic-universe}, connective spectra
    form a full symmetric monoidal subcategory of spectra. An algebra object
    in this subcategory is therefore a commutative ring spectrum whose
    underlying spectrum is connective.

    \item \textbf{Identify the essential image.}
    Conversely, if a commutative ring spectrum has connective underlying
    spectrum, all of its operadic structure maps already lie in the full
    subcategory \(\Sp_{\geq0}\). Full faithfulness and the displayed
    equivalence follow. Compare \cite{LurieHigherAlgebra,LurieSAG}.
\end{enumerate}
\end{proof}

\begin{proposition}[Presentability, connective cover, and limits]
\label{prop:connective-calg-presentable}
The \(\infty\)-category
\[
    \CAlg(\Sp)^{\mathrm{cn}}
\]
is presentable. The inclusion
\[
    \iota
    \colon
    \CAlg(\Sp)^{\mathrm{cn}}
    \hookrightarrow
    \CAlg(\Sp)
\]
preserves small colimits and admits a right adjoint
\[
    \tau_{\geq0}^{\CAlg}
    \colon
    \CAlg(\Sp)
    \longrightarrow
    \CAlg(\Sp)^{\mathrm{cn}},
\]
with
\[
    \operatorname{oblv}_{\CAlg}\bigl(\tau_{\geq0}^{\CAlg}A\bigr)
    \simeq
    \tau_{\geq0}^{\Sp}\operatorname{oblv}_{\CAlg}(A).
\]
Thus
\[
\begin{tikzcd}[column sep=large]
    \CAlg(\Sp)^{\mathrm{cn}}
        \arrow[r, shift left=1.1ex, "\iota"]
    &
    \CAlg(\Sp)
        \arrow[l, swap, shift left=1.1ex, "\tau_{\geq0}^{\CAlg}"']
\end{tikzcd}
\qquad
\iota\dashv\tau_{\geq0}^{\CAlg}.
\]
For every small diagram
\[
    D:K\longrightarrow\CAlg(\Sp)^{\mathrm{cn}},
\]
its limit is
\[
    \varprojlim\nolimits_K^{\CAlg(\Sp)^{\mathrm{cn}}}D
    \simeq
    \tau_{\geq0}^{\CAlg}
    \left(
        \varprojlim\nolimits_K^{\CAlg(\Sp)}\iota D
    \right).
\]
The inclusion \(\iota\) fails to preserve finite limits in general.
\end{proposition}

\begin{proof}
\leavevmode
\begin{enumerate}[label=\textup{(\arabic*)}]
    \item \textbf{Connective cover.}
    The connective-cover theorem for structured ring spectra identifies
    \(\CAlg(\Sp)^{\mathrm{cn}}\) as a presentable full subcategory closed under
    small colimits, and identifies the right adjoint of its inclusion with
    connective cover on underlying spectra; see
    \cite{LurieHigherAlgebra,LurieSAG}. Hence
    \[
        \iota\dashv\tau_{\geq0}^{\CAlg}
        \qquad\text{and}\qquad
        \operatorname{oblv}_{\CAlg}\tau_{\geq0}^{\CAlg}\simeq\tau_{\geq0}^{\Sp}\operatorname{oblv}_{\CAlg}.
    \]

    \item \textbf{Limit formula.}
    Put
    \[
        L:=\varprojlim\nolimits_K^{\CAlg(\Sp)}\iota D.
    \]
    For every connective commutative ring spectrum \(C\), adjunction and the
    universal property of \(L\) give
    \[
    \begin{aligned}
        \Map_{\CAlg(\Sp)^{\mathrm{cn}}}
        \bigl(C,\tau_{\geq0}^{\CAlg}L\bigr)
        &\simeq
        \Map_{\CAlg(\Sp)}(\iota C,L)
        \\
        &\simeq
        \varprojlim\nolimits_{k\in K}
        \Map_{\CAlg(\Sp)}(\iota C,\iota D(k))
        \\
        &\simeq
        \varprojlim\nolimits_{k\in K}
        \Map_{\CAlg(\Sp)^{\mathrm{cn}}}(C,D(k)).
    \end{aligned}
    \]
    Yoneda identifies \(\tau_{\geq0}^{\CAlg}L\) with the limit of \(D\).

    \item \textbf{Failure of finite-limit preservation.}
    Consider
    \[
        P
        :=
        H\mathbb Z
        \times_{H(\mathbb Z[x])}
        H\mathbb Z,
    \]
    where both maps are induced by inclusion of constant polynomials. The
    forgetful functor creates limits, so the underlying spectra fit into
    \[
    \begin{tikzcd}[column sep=large]
    P
        \arrow[r]
    &
    H\mathbb Z\oplus H\mathbb Z
        \arrow[r, "{\mathrm{diff}}"]
    &
    H(\mathbb Z[x]).
    \end{tikzcd}
    \]
    The long exact homotopy sequence gives
    \[
        \pi_{-1}(P)
        \simeq
        \operatorname{coker}
        \bigl(
            \mathbb Z\oplus\mathbb Z\to\mathbb Z[x]
        \bigr)
        \simeq
        \mathbb Z[x]/\mathbb Z,
    \]
    so \(P\) is nonconnective. The connective pullback is therefore
    \(\tau_{\geq0}^{\CAlg}P\).
\end{enumerate}
\end{proof}

\begin{corollary}[Connective pushouts]
\label{cor:connective-calg-pushouts}
If
\[
    B
    \longleftarrow
    A
    \longrightarrow
    C
\]
is a span in $\CAlg(\Sp)^{\mathrm{cn}}$, then the pushout
\[
    B\otimes_A C
\]
computed in $\CAlg(\Sp)$ is connective.
\end{corollary}

\begin{proof}
\leavevmode
\begin{enumerate}[label=\textup{(\arabic*)}]
    \item \textbf{Identify the pushout as a colimit.}
    It is a colimit of connective commutative ring spectra.

    \item \textbf{Apply closure under colimits.}
    Apply
    \Ref{prop:connective-calg-presentable}.
\end{enumerate}
\end{proof}

\subsection{The \texorpdfstring{$0$}{0}-truncated algebraic sector}

\begin{proposition}[Discrete commutative ring spectra]
\label{prop:discrete-commutative-ring-spectra}
Operadic Postnikov truncation induces an equivalence
\[
    \CAlg(\Sp)^{\mathrm{cn}}_{\leq0}
    \xrightarrow{\ \pi_0\ }
    \operatorname{CRing}
\]
whose inverse is the Eilenberg--Mac Lane construction
\[
    H_{\mathrm{CRing}}
    \colon
    \operatorname{CRing}
    \hookrightarrow
    \CAlg(\Sp)^{\mathrm{cn}}.
\]
Consequently, \(H_{\mathrm{CRing}}\) is fully faithful and its essential image consists of the
commutative ring spectra with discrete underlying spectrum.
\end{proposition}

\begin{proof}
\leavevmode
\begin{enumerate}[label=\textup{(\arabic*)}]
    \item \textbf{Operadic truncation.}
    The Postnikov truncation functors on spectra lift to commutative algebra
    objects, and the full subcategory of connective \(0\)-truncated
    commutative ring spectra is equivalent to ordinary commutative rings; see
    \cite{LurieHigherAlgebra,LurieSAG}.

    \item \textbf{Identification of the inverse.}
    Under this equivalence, degree-zero homotopy is the functor to ordinary
    commutative rings, while the inverse sends \(R\) to \(HR\). The asserted
    full faithfulness and description of the essential image follow.
\end{enumerate}
\end{proof}

\begin{proposition}[Degree-zero reflection adjunction]
\label{prop:pi-zero-h-adjunction}
The functor
\[
    \pi_0
    \colon
    \CAlg(\Sp)^{\mathrm{cn}}
    \longrightarrow
    \operatorname{CRing}
\]
is left adjoint to the Eilenberg--Mac Lane functor. Thus
\[
\begin{tikzcd}[column sep=large]
    \CAlg(\Sp)^{\mathrm{cn}}
        \arrow[r, shift left=1.1ex, "\pi_0"]
    &
    \operatorname{CRing}
        \arrow[l, swap, shift left=1.1ex, "H_{\mathrm{CRing}}"']
\end{tikzcd}
\qquad
\pi_0\dashv H_{\mathrm{CRing}},
\]
and there are natural equivalences
\[
    \Map_{\operatorname{CRing}}(\pi_0(A),R)
    \simeq
    \Map_{\CAlg(\Sp)}(A,HR).
\]
The adjunction unit is the natural morphism
\[
    A\longrightarrow H\pi_0(A).
\]
\end{proposition}

\begin{proof}
\leavevmode
\begin{enumerate}[label=\textup{(\arabic*)}]
    \item \textbf{Reflection to discrete rings.}
    By \Ref{prop:discrete-commutative-ring-spectra}, ordinary commutative
    rings identify with the full subcategory of connective \(0\)-truncated
    commutative ring spectra.

    \item \textbf{Mapping-space computation.}
    Operadic \(0\)-truncation is left adjoint to the inclusion of this full
    subcategory. Under the equivalence with \(\operatorname{CRing}\), the
    reflector is \(\pi_0\) and the inclusion is \(H_{\mathrm{CRing}}\). This gives the
    displayed adjunction and its unit; see \cite{LurieHigherAlgebra}.
\end{enumerate}
\end{proof}

\subsection{Operadic Postnikov stages}

\begin{proposition}[Operadic Postnikov truncations]
\label{prop:operadic-postnikov-truncations}
For every integer $n\geq0$, the truncation functor on spectra lifts to a
functor
\[
    \tau_{\leq n}^{\CAlg}
    \colon
    \CAlg(\Sp)^{\mathrm{cn}}
    \longrightarrow
    \CAlg(\Sp)^{\mathrm{cn}}
\]
whose underlying spectrum is the ordinary Postnikov truncation. For every
connective commutative ring spectrum $A$, there is a natural equivalence
\[
    A
    \simeq
    \varprojlim\nolimits_n\tau_{\leq n}^{\CAlg}A
\]
in $\CAlg(\Sp)$.
\end{proposition}

\begin{proof}
\leavevmode
\begin{enumerate}[label=\textup{(\arabic*)}]
    \item \textbf{Use reflectivity of truncated algebras.}
    The full subcategory of $n$-truncated objects in commutative ring spectra
    is reflective.

    \item \textbf{Identify the reflector on connective objects.}
    The forgetful functor identifies the reflector on
    connective objects with the spectral Postnikov truncation.

    \item \textbf{Transfer Postnikov convergence.}
    Since the
    forgetful functor creates limits, Postnikov convergence in spectra from
    \Ref{prop:postnikov-completeness-detection} implies convergence in
    commutative ring spectra. See \cite{LurieHigherAlgebra,LurieDAGIV}.
\end{enumerate}
\end{proof}

\begin{remark}[Postnikov notation on algebra objects]
\label{rem:postnikov-not-t-structure-on-algebras}
The notation
\[
    \CAlg(\Sp)^{\mathrm{cn}},
    \qquad
    \CAlg(\Sp)^{\mathrm{cn}}_{\leq n},
    \qquad
    \tau_{\leq n}^{\CAlg}
\]
is inherited from the Postnikov structure on underlying spectra: connectivity
and truncation are detected after forgetting the commutative algebra
structure. No $t$-structure on the nonstable $\infty$-category
$\CAlg(\Sp)$ is being asserted. The same convention will apply to
connective quasi-coherent commutative algebra objects, whose connectivity is
measured in the underlying stable coefficient category.
\end{remark}

The coordinate category \(\CAlg(\Sp)^{\mathrm{cn}}\) is now equipped with
the free, reflective, discrete, and Postnikov structures needed later. Its
next role is geometric rather than algebraic: variance reversal turns these
models into affine spectral contexts.

\section{Affine spectral contexts}
\label{sec:affine-spectral-schemes}

The connective coordinate models now become geometry. Reversing variance
turns a connective commutative ring spectrum into an affine spectral context;
algebraic pushouts consequently become geometric fibre products, while
\(\pi_0\) provides the topological support on which Zariski locality will be
built. This section records that affine geometric interface before stable
coefficients are attached to it.

\begin{definition}[Affine spectral scheme]
\label{def:affine-spectral-scheme}
The $\infty$-category of \textbf{affine spectral schemes} is
\[
    \Aff_{\Sp}
    :=
    \bigl(\CAlg(\Sp)^{\mathrm{cn}}\bigr)^{\op}.
\]
For
\[
    A\in\CAlg(\Sp)^{\mathrm{cn}},
\]
the corresponding affine spectral scheme is denoted
\[
    \Spec(A).
\]
\end{definition}

\begin{proposition}[Mapping spaces of affine spectral schemes]
\label{prop:affine-spectral-mapping-spaces}
For connective commutative ring spectra $A$ and $B$, there is a natural
equivalence
\[
    \Map_{\Aff_{\Sp}}
    \bigl(
        \Spec(B),
        \Spec(A)
    \bigr)
    \simeq
    \Map_{\CAlg(\Sp)}(A,B).
\]
\end{proposition}

\begin{proof}
\leavevmode
\begin{enumerate}[label=\textup{(\arabic*)}]
    \item \textbf{Unfold the opposite category.}
    The mapping-space equivalence is the definition of the opposite $\infty$-category.
\end{enumerate}
\end{proof}

Thus an affine context morphism reverses the direction of its coordinate map:
\[
\begin{tikzcd}[column sep=large]
    A
        \arrow[r, "\varphi"]
    &
    B
        \arrow[r, phantom, "\Longleftrightarrow" description]
    &
    \Spec(B)
        \arrow[r, "{\Spec(\varphi)}"]
    &
    \Spec(A).
\end{tikzcd}
\]

\subsection{Finite limits of affine spectral schemes}

\begin{proposition}[Affine spectral fibre products]
\label{prop:affine-spectral-fibre-products}
For a diagram of affine spectral schemes
\[
    \Spec(B)
    \longrightarrow
    \Spec(A)
    \longleftarrow
    \Spec(C),
\]
there is a canonical equivalence
\[
    \Spec(B)
    \times_{\Spec(A)}
    \Spec(C)
    \simeq
    \Spec(B\otimes_A C).
\]
The commutative ring spectrum $B\otimes_A C$ is connective.
\end{proposition}

\begin{proof}
\leavevmode
\begin{enumerate}[label=\textup{(\arabic*)}]
    \item \textbf{Reverse the pushout.}
    The fibre product in the opposite category is represented by the pushout
    of commutative ring spectra.

    \item \textbf{Check connectivity.}
    Connectivity follows from
    \Ref{cor:connective-calg-pushouts}.
\end{enumerate}
\end{proof}

This variance reversal is summarized by the correspondence between algebraic
pushouts and affine fibre products:
\[
\begin{tikzcd}[column sep=large, row sep=large]
    A \arrow[r] \arrow[d] & B \arrow[d] \\
    C \arrow[r] & B\otimes_AC
\end{tikzcd}
\qquad\longleftrightarrow\qquad
\begin{tikzcd}[column sep=large, row sep=large]
    \Spec(B\otimes_AC) \arrow[r] \arrow[d]
        & \Spec(B) \arrow[d] \\
    \Spec(C) \arrow[r] & \Spec(A).
\end{tikzcd}
\]

\begin{corollary}[Finite-limit closure]
\label{cor:affine-spectral-finite-limits}
The $\infty$-category $\Aff_{\Sp}$ admits finite limits.
\end{corollary}

\begin{proof}
\leavevmode
\begin{enumerate}[label=\textup{(\arabic*)}]
    \item \textbf{Identify the terminal object.}
    The terminal affine spectral scheme is $\Spec(\mathbb S)$.

    \item \textbf{Construct binary fibre products.}
    Binary
    fibre products exist by \Ref{prop:affine-spectral-fibre-products}.
\end{enumerate}
\end{proof}

\subsection{Underlying topological support}

\begin{definition}[Underlying topological space]
\label{def:underlying-topological-space-affine-spectral}
The underlying topological space of $X=\Spec(A)$ is
\[
    |X|
    :=
    \left|\Spec\pi_0(A)\right|.
\]
\end{definition}

We have therefore fixed the affine geometric base: connective coordinate
algebras determine affine contexts, their pushouts determine fibre products,
and \(\lvert\Spec\pi_0(A)\rvert\) supplies the support for locality. The next
step is to attach the stable coefficient fibre carried by each such context.

\section{Stable coefficients in affine context}
\label{sec:module-spectra-relative-algebras}

The stable linear information needed later lives in \(\Mod_A\) for
\(X=\Spec(A)\), and an affine context morphism acts on it by extension of
scalars. This section
constructs the affine coefficient system: substitution is extension of
scalars, restriction of scalars is its right adjoint, and relative
commutative algebras already encode affine extensions over a fixed base.

Fix a connective commutative ring spectrum
\[
    A\in\CAlg(\Sp)^{\mathrm{cn}}.
\]

\begin{definition}[Module spectrum]
\label{def:module-spectrum-over-a}
The $\infty$-category of $A$-module spectra is denoted
\[
    \Mod_A(\Sp)
\]
or, when no ambiguity is possible, simply
\[
    \Mod_A.
\]
\end{definition}

\begin{theorem}[Spectral module category]
\label{thm:spectral-module-category}
The category
\[
    \Mod_A
\]
is a presentably stable symmetric monoidal $\infty$-category. Its tensor
product is the relative tensor product
\[
    -\otimes_A-,
\]
its tensor unit is the regular module $A$, and its tensor product preserves
small colimits separately in each variable.
\end{theorem}

\begin{proof}
\leavevmode
\begin{enumerate}[label=\textup{(\arabic*)}]
    \item \textbf{Apply the operadic module construction.}
    This is the operadic module construction for a commutative algebra object
    in a presentably symmetric monoidal stable $\infty$-category; see
    \cite[Theorem~4.5.2.1]{LurieHigherAlgebra}.

    \item \textbf{Create limits and colimits in spectra.}
    Limits and colimits are created
    in spectra.

    \item \textbf{Deduce stability.}
    Stability follows from stability of $\Sp$.
\end{enumerate}
\end{proof}

\begin{construction}[Relative tensor product]
\label{con:spectral-relative-tensor-product}
For a right $A$-module $M$ and a left $A$-module $N$, the relative tensor
product is represented by the geometric realization of the two-sided bar
construction
\[
    M\otimes_A N
    \simeq
    \left|
        [n]
        \longmapsto
        M\otimes A^{\otimes n}\otimes N
    \right|.
\]
Since $A$ is commutative, its left and right module theories are
canonically identified.
\end{construction}

\begin{proposition}[Closedness of module spectra]
\label{prop:module-spectra-closedness}
For every $A$-module $M$, the functor
\[
    -\otimes_A M
    \colon
    \Mod_A
    \longrightarrow
    \Mod_A
\]
admits a right adjoint
\[
    \HHom_A(M,-).
\]
\end{proposition}

\begin{proof}
\leavevmode
\begin{enumerate}[label=\textup{(\arabic*)}]
    \item \textbf{Preserve small colimits.}
    The functor preserves small colimits by
    \Ref{thm:spectral-module-category}.

    \item \textbf{Apply the adjoint functor theorem.}
    Presentability of $\Mod_A$ and the
    adjoint functor theorem supply the right adjoint.
\end{enumerate}
\end{proof}

\subsection{Extension and restriction of scalars}

\begin{definition}[Extension and restriction of scalars]
\label{def:spectral-extension-restriction-scalars}
Let
\[
    \varphi
    \colon
    A
    \longrightarrow
    B
\]
be a morphism of commutative ring spectra. The extension-of-scalars
functor is
\[
    \operatorname{Ext}_{\varphi}
    :=
    -\otimes_A B
    \colon
    \Mod_A
    \longrightarrow
    \Mod_B.
\]
Its right adjoint is the restriction-of-scalars functor
\[
    \operatorname{Res}_{\varphi}
    \colon
    \Mod_B
    \longrightarrow
    \Mod_A.
\]
\end{definition}

The substitution/right-adjoint pair is therefore
\[
\begin{tikzcd}[column sep=large]
    \Mod_A
        \arrow[r, shift left=1.1ex, "\operatorname{Ext}_{\varphi}"]
    &
    \Mod_B
        \arrow[l, swap, shift left=1.1ex, "\operatorname{Res}_{\varphi}"'].
\end{tikzcd}
\]

For composable coordinate maps \(A\to B\to C\), coherent composition of
substitution has the form
\[
\begin{tikzcd}[column sep=large, row sep=large]
    \Mod_A
        \arrow[r, "-\otimes_AB"]
        \arrow[dr, "-\otimes_AC"']
    &
    \Mod_B
        \arrow[d, "-\otimes_BC"]
    \\
    &
    \Mod_C.
\end{tikzcd}
\]

\begin{proposition}[Monoidality of spectral base change]
\label{prop:spectral-base-change-monoidal}
The functor
\[
    \operatorname{Ext}_{\varphi}
    =
    -\otimes_A B
\]
is colimit-preserving and strong symmetric monoidal. In particular, there
is a natural equivalence
\[
    (M\otimes_A N)\otimes_A B
    \simeq
    (M\otimes_A B)
    \otimes_B
    (N\otimes_A B).
\]
\end{proposition}

\begin{proof}
\leavevmode
\begin{enumerate}[label=\textup{(\arabic*)}]
    \item \textbf{Identify extension of scalars.}
    The coCartesian operadic module construction identifies the transition
    functor associated to $\varphi$ with extension of scalars.

    \item \textbf{Supply the monoidal structure.}
    The same construction equips it with its
    strong symmetric monoidal structure; see
    \cite[Theorem~4.5.3.1 and Remark~4.5.3.2]{LurieHigherAlgebra}.
\end{enumerate}
\end{proof}

Strong symmetric monoidality is the compatibility needed to transport the
coefficient tensor product through substitution:
\[
\begin{tikzcd}[column sep=large, row sep=large]
    \Mod_A\times\Mod_A
        \arrow[r, "-\otimes_A-"]
        \arrow[d, "\operatorname{Ext}_{\varphi}\times \operatorname{Ext}_{\varphi}"']
    &
    \Mod_A
        \arrow[d, "\operatorname{Ext}_{\varphi}"]
    \\
    \Mod_B\times\Mod_B
        \arrow[r, "-\otimes_B-"']
    &
    \Mod_B.
\end{tikzcd}
\]

\begin{proposition}[Functoriality of module categories]
\label{prop:spectral-module-functoriality}
The assignments
\[
    A
    \longmapsto
    \Mod_A
\]
and
\[
    (A\to B)
    \longmapsto
    (-\otimes_A B)
\]
define a functor
\[
    \mathbf{Mod}_{\Sp}
    \colon
    \CAlg(\Sp)
    \longrightarrow
    \CAlg(\Pr^L_{\mathrm{st}}).
\]
For composable maps $A\to B\to C$, there is a canonical coherent
equivalence
\[
    (-\otimes_A B)\otimes_B C
    \simeq
    -\otimes_A C.
\]
\end{proposition}

\begin{proof}
\leavevmode
\begin{enumerate}[label=\textup{(\arabic*)}]
    \item \textbf{Straighten the module family.}
    Straightening the coCartesian family of module categories over
    $\CAlg(\Sp)$ gives the asserted functor.

    \item \textbf{Use coCartesian composition.}
    The coCartesian composition law
    supplies the coherent comparison for composable algebra maps; see
    \cite{LurieHigherAlgebra}.
\end{enumerate}
\end{proof}

\subsection{The affine coefficient system}

\begin{definition}[Affine quasi-coherent modules]
\label{def:affine-spectral-qcoh}
For an affine spectral scheme $\Spec(A)$, define
\[
    \QCoh_{\mathrm{aff}}(\Spec(A))
    :=
    \Mod_A.
\]
For a morphism
\[
    \Spec(B)
    \longrightarrow
    \Spec(A),
\]
the pullback functor is extension of scalars
\[
    -\otimes_A B
    \colon
    \Mod_A
    \longrightarrow
    \Mod_B.
\]
\end{definition}

\begin{proposition}[Affine quasi-coherent coefficient system]
\label{prop:affine-qcoh-coefficient-system}
The assignments of \Ref{def:affine-spectral-qcoh} define a functor
\[
    \QCoh_{\mathrm{aff}}
    \colon
    (\Aff_{\Sp})^{\op}
    \longrightarrow
    \CAlg(\Pr^L_{\mathrm{st}}).
\]
\end{proposition}

\begin{proof}
\leavevmode
\begin{enumerate}[label=\textup{(\arabic*)}]
    \item \textbf{Restrict the module functor.}
    This is the restriction of the functor
    $\mathbf{Mod}_{\Sp}$ from
    \Ref{prop:spectral-module-functoriality} to connective commutative ring
    spectra.
\end{enumerate}
\end{proof}


\begin{remark}[Affine indexed doctrine]
\label{rem:affine-indexed-spectral-doctrine}
The variance reversal from coordinate algebras to affine geometry also
reverses the way the module functor is read. The covariant algebraic functor
of \Ref{prop:spectral-module-functoriality} becomes
\[
    \QCoh_{\mathrm{aff}}
    :
    (\Aff_{\Sp})^{\op}
    \longrightarrow
    \CAlg(\Pr^L_{\mathrm{st}}).
\]
Accordingly, extension of scalars is substitution along an affine context
morphism, and restriction of scalars is its right adjoint. This is the
complete affine form of the indexed coefficient doctrine; the generalized
and global forms will be obtained from it by extension and descent.
\end{remark}

\subsection{Relative commutative algebra objects}

\begin{proposition}[Relative commutative spectral algebras]
\label{prop:relative-commutative-spectral-algebras}
There is a natural equivalence
\[
    \CAlg(\Mod_A)
    \simeq
    \CAlg(\Sp)_{A/}.
\]
Under this equivalence, a commutative algebra object $B$ in $A$-modules is
identified with its structural morphism
\[
    A
    \longrightarrow
    B.
\]
\end{proposition}

\begin{proof}
\leavevmode
\begin{enumerate}[label=\textup{(\arabic*)}]
    \item \textbf{Specialize the operadic equivalence.}
    This is the commutative specialization of the operadic equivalence between
    algebras in a module category and algebras under the base algebra; see
    \cite[Corollary~3.4.1.7]{LurieHigherAlgebra}.
\end{enumerate}
\end{proof}

\begin{proposition}[Connective relative algebras]
\label{prop:connective-relative-algebras}
The equivalence of \Ref{prop:relative-commutative-spectral-algebras}
restricts to an equivalence
\[
    \CAlg(\Mod_A)^{\mathrm{cn}}
    \simeq
    \CAlg(\Sp)^{\mathrm{cn}}_{A/},
\]
where connectivity on the left is detected on the underlying $A$-module.
Consequently,
\[
    \Aff_{\Sp}/\Spec(A)
    \simeq
    \bigl(
        \CAlg(\Mod_A)^{\mathrm{cn}}
    \bigr)^{\op}.
\]
\end{proposition}

\begin{proof}
\leavevmode
\begin{enumerate}[label=\textup{(\arabic*)}]
    \item \textbf{Preserve connective objects.}
    The equivalence preserves underlying spectra and therefore preserves
    connective objects.

    \item \textbf{Pass to opposites.}
    Passing to opposite categories gives the affine
    statement.
\end{enumerate}
\end{proof}

\subsection{The induced module \texorpdfstring{$t$}{t}-structure}

\begin{definition}[Module Postnikov $t$-structure]
\label{def:module-postnikov-t-structure}
Let
\[
    (\Mod_{A})_{\geq0}
\]
be the full subcategory of $A$-modules whose underlying spectra are
connective, and let
\[
    (\Mod_{A})_{\leq0}
\]
be the full subcategory whose underlying spectra are coconnective.
\end{definition}

\begin{proposition}[Induced module $t$-structure]
\label{prop:induced-module-t-structure}
The pair
\[
    ((\Mod_{A})_{\geq0},(\Mod_{A})_{\leq0})
\]
is an accessible, left complete, and right complete $t$-structure on
$\Mod_A$. The forgetful functor
\[
    \operatorname{oblv}_{A}
    \colon
    \Mod_A
    \longrightarrow
    \Sp
\]
is $t$-exact and creates Postnikov truncations.
\end{proposition}

\begin{proof}
\leavevmode
\begin{enumerate}[label=\textup{(\arabic*)}]
    \item \textbf{Create limits and colimits in spectra.}
    Limits and colimits of $A$-modules are created in spectra.

    \item \textbf{Induce the truncation module structure.}
    The module
    structure on a truncation is induced by the connective multiplication of
    $A$.

    \item \textbf{Detect the $t$-structure axioms.}
    The orthogonality and truncation axioms are detected after
    forgetting to spectra.

    \item \textbf{Transfer completeness.}
    Completeness follows from completeness of the
    Postnikov $t$-structure and creation of limits and colimits. See
    \cite{LurieHigherAlgebra,LurieSAG}.
\end{enumerate}
\end{proof}

\begin{proposition}[Heart of the module category]
\label{prop:heart-module-category}
There is a natural equivalence of abelian categories
\[
    \Mod_A^\heartsuit
    \simeq
    \operatorname{Mod}_{\pi_0(A)}.
\]
\end{proposition}

\begin{proof}
\leavevmode
\begin{enumerate}[label=\textup{(\arabic*)}]
    \item \textbf{Recover the underlying discrete module.}
    A discrete $A$-module is determined by its degree-zero homotopy group.

    \item \textbf{Factor the action through degree zero.}
    The
    $A$-action factors uniquely through the discrete commutative ring spectrum
    $H\pi_0(A)$.

    \item \textbf{Identify the ordinary module structure.}
    It is therefore equivalent to an ordinary $\pi_0(A)$-module
    structure.
\end{enumerate}
\end{proof}

\begin{proposition}[Exactness and right $t$-exactness of base change]
\label{prop:base-change-exact-right-t-exact}
Let $\varphi:A\to B$ be a morphism of connective commutative ring spectra.
Then:
\begin{enumerate}[label=(\roman*)]
    \item the extension-of-scalars functor
    \[
        \operatorname{Ext}_{\varphi}
        \colon
        \Mod_A
        \longrightarrow
        \Mod_B
    \]
    is exact as a functor of stable $\infty$-categories;

    \item the functor $\operatorname{Ext}_{\varphi}$ is right $t$-exact;

    \item the restriction-of-scalars functor $\operatorname{Res}_{\varphi}$ is $t$-exact.
\end{enumerate}
\end{proposition}

\begin{proof}
\leavevmode
\begin{enumerate}[label=\textup{(\arabic*)}]
    \item \textbf{Stable exactness.}
    The functor \(\operatorname{Ext}_{\varphi}=-\otimes_AB\) preserves all colimits. A
    colimit-preserving functor between stable \(\infty\)-categories preserves
    finite colimits, and finite colimits agree with finite limits. Hence
    \(\operatorname{Ext}_{\varphi}\) is exact.

    \item \textbf{Right $t$-exactness.}
    The \(A\)-module \(B\) is connective. The relative tensor product of two
    connective \(A\)-modules is connective, so \(\operatorname{Ext}_{\varphi}\) preserves connective
    objects.

    \item \textbf{Restriction of scalars.}
    Restriction of scalars preserves the underlying spectrum of a module.
    Therefore it preserves both connective and coconnective objects and is
    \(t\)-exact.
\end{enumerate}
\end{proof}

\begin{remark}[Stable exactness and geometric flatness]
\label{rem:stable-exactness-not-flatness}
Stable exactness and geometric flatness measure different properties.
Extension of scalars is exact between stable \(\infty\)-categories for every
morphism of commutative ring spectra. Flatness adds compatibility with the
Postnikov \(t\)-structures, as formalized in
\Ref{def:flat-module-spectrum}.
\end{remark}

The affine coefficient system is now in place: every affine context carries a
stable symmetric-monoidal fibre, substitution is coherent extension of
scalars, and relative commutative algebras already describe affine extensions
over that context. The remaining local algebraic issue is to distinguish
geometrically well-behaved substitutions from the exactness that is automatic
in every stable fibre.

\section{Substitution, geometric exactness, and finiteness}
\label{sec:spectral-flatness-finiteness}

Geometric regularity is detected by compatibility with the Postnikov
\(t\)-structures: flatness will mean that substitution is \(t\)-exact.
Extension of scalars remains exact for every morphism of commutative ring
spectra. This section develops
that condition together with the finiteness, perfectness, and projectivity
notions required later for Zariski locality.

\subsection{Flat module spectra}

\begin{definition}[Flat module spectrum]
\label{def:flat-module-spectrum}
Let $A$ be a connective commutative ring spectrum. A connective
$A$-module $M$ is \textbf{flat} if the functor
\[
    -\otimes_A M
    \colon
    \Mod_A
    \longrightarrow
    \Mod_A
\]
is $t$-exact.
\end{definition}

For a connective module, tensoring is already right \(t\)-exact. Thus the
additional content of flatness is left \(t\)-exactness: tensoring with
\(M\) must also preserve coconnective modules.

\begin{theorem}[Spectral Lazard theorem]
\label{thm:spectral-lazard}
Let \(A\) be a connective commutative ring spectrum and let
\(M\in(\Mod_{A})_{\geq0}\). The following conditions are equivalent.
\begin{enumerate}[label=(\roman*)]
    \item The module \(M\) is flat.

    \item The \(\pi_0(A)\)-module \(\pi_0(M)\) is flat and, for every integer
    \(n\), the canonical morphism
    \[
        \pi_n(A)\otimes_{\pi_0(A)}\pi_0(M)
        \longrightarrow
        \pi_n(M)
    \]
    is an isomorphism.

    \item For every \(A\)-module \(N\) and every integer \(n\), the canonical
    morphism
    \[
        \pi_n(N)\otimes_{\pi_0(A)}\pi_0(M)
        \longrightarrow
        \pi_n(N\otimes_A M)
    \]
    is an isomorphism.

    \item The module \(M\) is a filtered colimit of finite free
    \(A\)-modules.
\end{enumerate}
\end{theorem}

\begin{proof}
\leavevmode
\begin{enumerate}[label=\textup{(\arabic*)}]
    \item \textbf{Invoke the spectral flatness criterion.}
    For a connective base ring, the spectral flatness criterion and Lazard
    theorem identify the homotopy-group condition in \textup{(ii)}, left
    \(t\)-exactness of \(-\otimes_A M\), and expression of \(M\) as a
    filtered colimit of finite free modules; see
    \cite[Definition~7.2.2.10 and Theorem~7.2.2.15]{LurieHigherAlgebra}.

    \item \textbf{Use right $t$-exactness.}
    Since \(M\) is connective, tensoring with \(M\) is already right
    \(t\)-exact.

    \item \textbf{Identify three equivalent conditions.}
    Thus the cited theorem gives
    \[
        \textup{(i)}\Longleftrightarrow\textup{(ii)}
        \Longleftrightarrow\textup{(iv)}.
    \]

    \item \textbf{Choose a filtered free presentation.}
    Assume \textup{(iv)} and choose
    \[
        M\simeq\varinjlim_\alpha F_\alpha
    \]
    with each \(F_\alpha\) finite free.

    \item \textbf{Compute homotopy groups after tensoring.}
    For every \(A\)-module \(N\),
    \[
    \begin{aligned}
        \pi_n(N\otimes_A M)
        &\simeq
        \varinjlim_\alpha\pi_n(N\otimes_A F_\alpha)
        \\
        &\simeq
        \varinjlim_\alpha
        \bigl(\pi_n(N)\otimes_{\pi_0(A)}\pi_0(F_\alpha)\bigr)
        \\
        &\simeq
        \pi_n(N)\otimes_{\pi_0(A)}\pi_0(M).
    \end{aligned}
    \]

    \item \textbf{Pass through filtered colimits.}
    Here tensor products and homotopy groups commute with filtered colimits,
    and the formula for a finite free module is immediate. This proves
    \textup{(iv)}\(\Rightarrow\)\textup{(iii)}.

    \item \textbf{Recover the coefficient formula.}
    Assume \textup{(iii)}. Taking \(N=A\) gives the homotopy-group formula
    in \textup{(ii)}. It remains to prove that \(\pi_0(M)\) is flat over
    \(\pi_0(A)\).

    \item \textbf{Choose a short exact sequence.}
    Let
    \[
        0\longrightarrow P'
        \longrightarrow P
        \longrightarrow P''
        \longrightarrow0
    \]
    be a short exact sequence of ordinary \(\pi_0(A)\)-modules, regarded as
    discrete \(A\)-modules. It determines a fibre sequence in \(\Mod_A\).

    \item \textbf{Tensor the fibre sequence.}
    Since tensoring with \(M\) is exact as a functor of stable
    \(\infty\)-categories, there is a fibre sequence
    \[
        P'\otimes_AM
        \longrightarrow
        P\otimes_AM
        \longrightarrow
        P''\otimes_AM.
    \]

    \item \textbf{Apply condition \textup{(iii)}.}
    Condition \textup{(iii)} shows that each of these three modules is
    discrete and identifies its degree-zero homotopy group with the ordinary
    tensor product by \(\pi_0(M)\).

    \item \textbf{Take degree-zero homotopy.}
    Taking \(\pi_0\) therefore gives a short
    exact sequence
    \[
        0\longrightarrow
        P'\otimes_{\pi_0(A)}\pi_0(M)
        \longrightarrow
        P\otimes_{\pi_0(A)}\pi_0(M)
        \longrightarrow
        P''\otimes_{\pi_0(A)}\pi_0(M)
        \longrightarrow0.
    \]

    \item \textbf{Deduce ordinary flatness.}
    Thus \(-\otimes_{\pi_0(A)}\pi_0(M)\) is exact, so \(\pi_0(M)\) is flat.

    \item \textbf{Complete the equivalence proof.}
    This proves \textup{(iii)}\(\Rightarrow\)\textup{(ii)}, and all four
    conditions are equivalent.
\end{enumerate}
\end{proof}

\begin{corollary}[Homotopy groups of a flat module]
\label{cor:flat-modules-detected-homotopy}
If $M$ is a flat $A$-module, then its graded homotopy groups are obtained
from $\pi_0(M)$ by base change:
\[
    \pi_n(M)
    \simeq
    \pi_n(A)
    \otimes_{\pi_0(A)}
    \pi_0(M)
\]
for every integer $n$.
\end{corollary}

\begin{proof}
\leavevmode
\begin{enumerate}[label=\textup{(\arabic*)}]
    \item \textbf{Apply the Lazard criterion.}
    Apply condition \textup{(ii)} of \Ref{thm:spectral-lazard}.
\end{enumerate}
\end{proof}

\begin{proposition}[Permanence of flat modules]
\label{prop:flat-module-permanence}
Let $A$ be a connective commutative ring spectrum.
\begin{enumerate}[label=(\roman*)]
    \item Finite free and finite projective $A$-modules are flat.

    \item Filtered colimits of flat $A$-modules are flat.

    \item If $M$ and $N$ are flat $A$-modules, then
    \[
        M\otimes_A N
    \]
    is flat over $A$.

    \item Let $A\to B$ be a flat morphism of connective commutative ring
    spectra. If $N$ is flat as a $B$-module, then $N$, regarded by
    restriction of scalars as an $A$-module, is flat over $A$.

    \item Let $A\to B$ be any morphism of connective commutative ring
    spectra. If $M$ is flat over $A$, then
    \[
        M\otimes_A B
    \]
    is flat over $B$.
\end{enumerate}
\end{proposition}

\begin{proof}
\leavevmode
\begin{enumerate}[label=\textup{(\arabic*)}]
    \item \textbf{Check finite free modules.}
    For a finite free module \(F\simeq A^{\oplus r}\), tensoring with \(F\)
    is a finite direct sum of the identity functor and is therefore
    \(t\)-exact.

    \item \textbf{Pass to retracts.}
    If \(M\) is a retract of a flat module \(F\), then
    \(-\otimes_AM\) is a retract of \(-\otimes_AF\). The connective and
    coconnective subcategories are closed under retracts, so finite projective
    modules are flat.

    \item \textbf{Commute tensoring with the filtered colimit.}
    For a filtered diagram \(\{M_\alpha\}\) of flat modules and every
    \(K\in\Mod_A\),
    \[
        K\otimes_A\varinjlim_\alpha M_\alpha
        \simeq
        \varinjlim_\alpha(K\otimes_AM_\alpha).
    \]

    \item \textbf{Preserve both halves of the $t$-structure.}
    Filtered colimits preserve both halves of the Postnikov \(t\)-structure,
    since homotopy groups commute with filtered colimits.

    \item \textbf{Conclude flatness of the colimit.}
    Hence the colimit
    module is flat.

    \item \textbf{Rewrite tensoring by the tensor product.}
    If \(M\) and \(N\) are flat, associativity gives
    \[
        -\otimes_A(M\otimes_AN)
        \simeq
        \bigl(-\otimes_AM\bigr)\otimes_AN.
    \]

    \item \textbf{Use $t$-exact composition.}
    The right side is a composite of \(t\)-exact functors, proving
    \textup{(iii)}.

    \item \textbf{Rewrite tensoring after restriction.}
    Let \(A\to B\) be flat and let \(N\) be flat over \(B\). For every
    \(A\)-module \(K\),
    \[
        K\otimes_AN
        \simeq
        (K\otimes_AB)\otimes_BN
    \]
    as underlying \(A\)-modules.

    \item \textbf{Compose the $t$-exact functors.}
    Extension to \(B\), tensoring with \(N\),
    and restriction back to \(A\) are all \(t\)-exact. Their composite is
    therefore \(t\)-exact.

    \item \textbf{Rewrite the base-changed tensor product.}
    Let \(A\to B\) be arbitrary and let \(M\) be flat over \(A\). For a
    \(B\)-module \(K\),
    \[
        K\otimes_B(M\otimes_AB)
        \simeq
        K\otimes_AM,
    \]
    where \(K\) on the right is restricted to \(A\).

    \item \textbf{Use restriction and flat tensoring.}
    Restriction preserves
    the underlying spectrum and hence the Postnikov \(t\)-structure, while
    tensoring with \(M\) is \(t\)-exact.

    \item \textbf{Conclude flatness after base change.}
    Thus \(M\otimes_AB\) is flat over
    \(B\).
\end{enumerate}
\end{proof}

\subsection{Flat morphisms of connective commutative ring spectra}

\begin{definition}[Flat spectral algebra morphism]
\label{def:flat-spectral-algebra-morphism}
A morphism
\[
    \varphi
    \colon
    A
    \longrightarrow
    B
\]
in $\CAlg(\Sp)^{\mathrm{cn}}$ is \textbf{flat} if $B$, regarded as an
$A$-module, is flat.
\end{definition}

\begin{proposition}[Homotopy groups of a flat algebra]
\label{prop:homotopy-groups-flat-algebra}
If $A\to B$ is flat, then $\pi_0(B)$ is flat over $\pi_0(A)$ and there are
natural isomorphisms
\[
    \pi_n(B)
    \simeq
    \pi_n(A)
    \otimes_{\pi_0(A)}
    \pi_0(B)
\]
for every integer $n$.
\end{proposition}

\begin{proof}
\leavevmode
\begin{enumerate}[label=\textup{(\arabic*)}]
    \item \textbf{Apply the Lazard criterion to the algebra.}
    Apply condition (ii) of \Ref{thm:spectral-lazard} to the $A$-module $B$.
\end{enumerate}
\end{proof}

\begin{proposition}[Permanence of flat algebra morphisms]
\label{prop:flat-algebra-permanence}
Flat morphisms in \(\CAlg(\Sp)^{\mathrm{cn}}\) are stable under equivalence,
composition, base change, and filtered colimits in the arrow category.
\end{proposition}

\begin{proof}
\leavevmode
\begin{enumerate}[label=\textup{(\arabic*)}]
    \item \textbf{Handle equivalences.}
    Stability under equivalence follows from the definition.

    \item \textbf{Factor extension of scalars under composition.}
    For flat maps
    \(A\to B\to C\), the extension-of-scalars functor factors as
    \[
    \begin{tikzcd}[column sep=large]
        \Mod_A
            \arrow[r, "-\otimes_A B"]
        &
        \Mod_B
            \arrow[r, "-\otimes_B C"]
        &
        \Mod_C.
    \end{tikzcd}
    \]

    \item \textbf{Conclude composition stability.}
    Both factors are \(t\)-exact, so their composite is \(t\)-exact.

    \item \textbf{Form the base-change square.}
    Given \(A\to A'\), form
    \[
    \begin{tikzcd}[column sep=large, row sep=large]
        A \arrow[r] \arrow[d] & B \arrow[d] \\
        A' \arrow[r] & B':=B\otimes_AA'.
    \end{tikzcd}
    \]

    \item \textbf{Rewrite the base-changed tensor product.}
    For every \(A'\)-module \(M\), associativity gives
    \[
        M\otimes_{A'}B'
        \simeq
        M\otimes_AB,
    \]
    where the right side uses restriction of scalars along \(A\to A'\).

    \item \textbf{Use $t$-exactness.}
    Restriction is \(t\)-exact and tensoring with \(B\) is \(t\)-exact;
    hence \(A'\to B'\) is flat.

    \item \textbf{Form the filtered colimits.}
    Let \(\{A_\alpha\to B_\alpha\}_\alpha\) be filtered and flat, and put
    \[
        A:=\varinjlim_\alpha A_\alpha,
        \qquad
        B:=\varinjlim_\alpha B_\alpha.
    \]

    \item \textbf{Base change each target to the colimit base.}
    Define
    \[
        B'_\alpha
        :=
        A\otimes_{A_\alpha}B_\alpha.
    \]

    \item \textbf{Preserve flatness under base change.}
    Base-change stability makes every \(A\to B'_\alpha\) flat.

    \item \textbf{Identify the filtered colimit in the undercategory.}
    The canonical
    morphisms \(B'_\alpha\to B\) exhibit \(B\) as the filtered colimit of the
    \(B'_\alpha\) in \(\CAlg(\Sp)_{A/}\): mapping into an \(A\)-algebra
    converts both sides into the same limit of mapping spaces.

    \item \textbf{Identify the underlying module colimit.}
    The underlying \(A\)-module of \(B\) is therefore a filtered colimit of
    flat \(A\)-modules.

    \item \textbf{Apply permanence of flat modules.}
    By \Ref{prop:flat-module-permanence}, it is flat, so
    \(A\to B\) is flat.
\end{enumerate}
\end{proof}

\begin{proposition}[Degree-zero homotopy of connective pushouts]
\label{prop:classical-truncation-affine-base-change}
Let
\[
\begin{tikzcd}[column sep=large, row sep=large]
    A
        \arrow[r]
        \arrow[d]
    &
    B
        \arrow[d]
    \\
    A'
        \arrow[r]
    &
    B':=B\otimes_AA'
\end{tikzcd}
\]
be a pushout square in connective commutative ring spectra. Then the
canonical morphism
\[
    \pi_0(B)
    \otimes_{\pi_0(A)}
    \pi_0(A')
    \longrightarrow
    \pi_0(B')
\]
is an isomorphism.
\end{proposition}

\begin{proof}
\leavevmode
\begin{enumerate}[label=\textup{(\arabic*)}]
    \item \textbf{Identify degree zero as a left adjoint.}
    The adjunction of \Ref{prop:pi-zero-h-adjunction} shows that
    \[
        \pi_0
        :
        \CAlg(\Sp)^{\mathrm{cn}}
        \longrightarrow
        \operatorname{CRing}
    \]
    is a left adjoint.

    \item \textbf{Preserve pushouts.}
    It therefore preserves all small colimits, in
    particular pushouts.

    \item \textbf{Compute degree zero of the pushout.}
    Hence
    \[
        \pi_0(B')
        \simeq
        \pi_0(B)
        \otimes_{\pi_0(A)}
        \pi_0(A').
    \]
\end{enumerate}
\end{proof}

\subsection{Local finite presentation and finite cell algebras}

\begin{definition}[Local finite presentation]
\label{def:finitely-presented-spectral-algebra}
A morphism
\[
    A\longrightarrow B
\]
of connective commutative ring spectra is \textbf{locally of finite
presentation} if \(B\) is compact in
\[
    \CAlg(\Sp)_{A/}.
\]
For an affine morphism
\[
    \Spec(B)\longrightarrow\Spec(A),
\]
we also say \textbf{of finite presentation}; affineness supplies the global
quasi-compactness condition. This terminology agrees with the spectral
geometric convention of \cite{LurieSAG,KhanBraveNewMotivicI}.
\end{definition}

\begin{proposition}[Compactness may be checked connectively]
\label{prop:compactness-connective-undercategory}
Let \(A\) and \(B\) be connective commutative ring spectra and regard \(B\)
as an \(A\)-algebra. Then \(B\) is compact in \(\CAlg(\Sp)_{A/}\) if and
only if it is compact in \(\CAlg(\Sp)^{\mathrm{cn}}_{A/}\).
\end{proposition}

\begin{proof}
\leavevmode
\begin{enumerate}[label=\textup{(\arabic*)}]
    \item \textbf{Construct the connective reflection adjunction.}
    Since \(A\) is connective, the connective-cover adjunction of
    \Ref{prop:connective-calg-presentable} induces an adjunction
    \[
    \begin{tikzcd}[column sep=large]
        \CAlg(\Sp)^{\mathrm{cn}}_{A/}
            \arrow[r, shift left=1.1ex, "\iota_A"]
        &
        \CAlg(\Sp)_{A/}
            \arrow[l, swap, shift left=1.1ex, "\tau^{\CAlg}_{\geq0,A}"'].
    \end{tikzcd}
    \]

    \item \textbf{Identify its right adjoint.}
    The right adjoint is obtained by applying connective cover to the target
    algebra; the structural map from \(A\) factors uniquely through that
    connective cover by adjunction.

    \item \textbf{Commute spectral connective cover with filtered colimits.}
    The functor \(\tau_{\geq0}^{\Sp}\) on spectra commutes with filtered colimits,
    and the forgetful functor from commutative ring spectra creates sifted
    colimits.

    \item \textbf{Transfer filtered-colimit preservation to the undercategory.}
    Hence \(\tau^{\CAlg}_{\geq0,A}\) commutes with filtered
    colimits.

    \item \textbf{Compute the mapping space against a filtered colimit.}
    If \(B\) is compact in the connective undercategory and
    \(\{C_\alpha\}\) is filtered in \(\CAlg(\Sp)_{A/}\), then
    \[
    \begin{aligned}
        \Map_{\CAlg(\Sp)_{A/}}
        \left(B,\varinjlim_\alpha C_\alpha\right)
        &\simeq
        \Map_{\CAlg(\Sp)^{\mathrm{cn}}_{A/}}
        \left(
            B,
            \tau^{\CAlg}_{\geq0,A}\varinjlim_\alpha C_\alpha
        \right)
        \\
        &\simeq
        \Map_{\CAlg(\Sp)^{\mathrm{cn}}_{A/}}
        \left(
            B,
            \varinjlim_\alpha\tau^{\CAlg}_{\geq0,A}C_\alpha
        \right)
        \\
        &\simeq
        \varinjlim_\alpha
        \Map_{\CAlg(\Sp)_{A/}}(B,C_\alpha).
    \end{aligned}
    \]

    \item \textbf{Conclude compactness in the unrestricted undercategory.}
    Thus \(B\) is compact in the unrestricted undercategory.

    \item \textbf{Prove the converse by restriction.}
    The converse
    follows by restricting the same compactness condition to filtered
    diagrams of connective \(A\)-algebras.
\end{enumerate}
\end{proof}

\begin{definition}[Finite cell algebra]
\label{def:finite-cell-spectral-algebra}
Write
\[
    \operatorname{Sym}_A
    \colon
    \Mod_A
    \longrightarrow
    \CAlg(\Mod_A)
    \simeq
    \CAlg(\Sp)_{A/}
\]
for the free commutative \(A\)-algebra functor. A connective commutative
\(A\)-algebra is a \textbf{finite cell algebra} if it belongs to the smallest
full subcategory of \(\CAlg(\Sp)^{\mathrm{cn}}_{A/}\) which contains
\[
    \operatorname{Sym}_A(A^{\oplus n}),
    \qquad n\geq0,
\]
and is closed under finite colimits. Equivalently, it is obtained by a finite
iteration of finite colimits starting from finitely generated free
commutative \(A\)-algebras. No separate presentation datum is part of the
definition.
\end{definition}

\begin{proposition}[Degree zero of relative free commutative algebras]
\label{prop:pi-zero-relative-free-algebra}
Let \(A\) be a connective commutative ring spectrum and let
\(M\in(\Mod_A)_{\geq0}\). Then the canonical map of ordinary commutative
\(\pi_0(A)\)-algebras is an isomorphism
\[
    \operatorname{Sym}_{\pi_0(A)}\bigl(\pi_0(M)\bigr)
    \xrightarrow{\ \cong\ }
    \pi_0\operatorname{Sym}_A(M).
\]
In particular,
\[
    \pi_0\operatorname{Sym}_A(A^{\oplus n})
    \simeq
    \pi_0(A)[t_1,\ldots,t_n].
\]
\end{proposition}

\begin{proof}
\leavevmode
\begin{enumerate}[label=\textup{(\arabic*)}]
    \item \textbf{Compute degree zero of a relative tensor product.}
    For connective \(A\)-modules \(M_1\) and \(M_2\), the bar construction
    of \Ref{con:spectral-relative-tensor-product} is a simplicial connective
    spectrum. On connective spectra, \(\pi_0\) preserves colimits, and the
    monoidal structure induced on the heart is the ordinary tensor product.
    Hence geometric realization of the bar construction gives
    \[
        \pi_0(M_1\otimes_AM_2)
        \simeq
        \pi_0(M_1)\otimes_{\pi_0(A)}\pi_0(M_2).
    \]

    \item \textbf{Iterate the tensor-product formula.}
    Iterating yields
    \[
        \pi_0(M^{\otimes_A r})
        \simeq
        \pi_0(M)^{\otimes_{\pi_0(A)}r}.
    \]

    \item \textbf{Write the relative free-algebra formula.}
    The relative free-algebra formula gives
    \[
        \operatorname{oblv}_A\operatorname{Sym}_A(M)
        \simeq
        \bigoplus_{r\geq0}
        \bigl(M^{\otimes_A r}\bigr)_{h\Sigma_r}.
    \]

    \item \textbf{Apply degree zero to the colimits.}
    Since every tensor power is connective, \(\pi_0\) carries the homotopy
    orbit colimit to ordinary coinvariants and carries the coproduct to the
    coproduct.

    \item \textbf{Compute the symmetric algebra.}
    Therefore
    \[
    \begin{aligned}
        \pi_0\operatorname{Sym}_A(M)
        &\simeq
        \bigoplus_{r\geq0}
        \left(
            \pi_0(M)^{\otimes_{\pi_0(A)}r}
        \right)_{\Sigma_r}
        \\
        &\simeq
        \operatorname{Sym}_{\pi_0(A)}\bigl(\pi_0(M)\bigr).
    \end{aligned}
    \]

    \item \textbf{Identify the multiplication.}
    The multiplication is the one induced by concatenation of tensor powers,
    hence agrees with ordinary symmetric-algebra multiplication.

    \item \textbf{Specialize to a finite free module.}
    Taking
    \(M=A^{\oplus n}\) gives the polynomial-algebra formula. Compare the
    relative free-algebra construction of \cite{LurieHigherAlgebra}.
\end{enumerate}
\end{proof}

\begin{remark}[Finite-cell and finite-presentation terminology]
\label{rem:finite-cell-terminology}
Two different finiteness notions are being kept distinct. The term
\textbf{finite cell} refers to the explicit finite-colimit closure of finitely
generated free commutative \(A\)-algebras. The phrase \textbf{locally of
finite presentation} refers instead to the compactness condition of
\Ref{def:finitely-presented-spectral-algebra}. In
\cite{LurieHigherAlgebra} the former class is called the subcategory of
algebras of finite presentation, while the latter objects are called locally
of finite presentation. The terminology here is chosen to keep construction
separate from compactness. Neither notion retains a choice of generators,
relations, or presentation data.
\end{remark}

\begin{proposition}[Finite-cell criterion]
\label{prop:finite-presentation-compactness-criterion}
A connective commutative \(A\)-algebra is locally of finite presentation if
and only if it is a retract of a finite cell commutative \(A\)-algebra.
\end{proposition}

\begin{proof}
\leavevmode
\begin{enumerate}[label=\textup{(\arabic*)}]
    \item \textbf{Define the finite free algebra.}
    For every integer \(n\geq0\), put
    \[
        F_n:=\operatorname{Sym}_A(A^{\oplus n}).
    \]

    \item \textbf{Compute mapping spaces against a filtered colimit.}
    Let \(\{B_\alpha\}\) be a filtered diagram in
    \(\CAlg(\Sp)^{\mathrm{cn}}_{A/}\). Since the forgetful functor to
    \(A\)-modules creates filtered colimits, the free-algebra adjunction gives
    \[
    \begin{aligned}
        \Map_{\CAlg(\Sp)^{\mathrm{cn}}_{A/}}
        \left(F_n,\varinjlim_\alpha B_\alpha\right)
        &\simeq
        \Map_{\Mod_A}
        \left(A^{\oplus n},\varinjlim_\alpha B_\alpha\right)
        \\
        &\simeq
        \varinjlim_\alpha
        \Map_{\Mod_A}(A^{\oplus n},B_\alpha)
        \\
        &\simeq
        \varinjlim_\alpha
        \Map_{\CAlg(\Sp)^{\mathrm{cn}}_{A/}}(F_n,B_\alpha).
    \end{aligned}
    \]

    \item \textbf{Conclude compactness of finite free algebras.}
    Thus every finitely generated free commutative \(A\)-algebra is compact.

    \item \textbf{Pass compactness through finite colimits.}
    Finite colimits of compact objects are compact because mapping out of a
    finite colimit is a finite limit of mapping spaces and finite limits of
    spaces commute with filtered colimits.

    \item \textbf{Pass compactness through retracts.}
    Hence every finite cell
    \(A\)-algebra is compact, and so is every retract of a finite cell
    \(A\)-algebra.

    \item \textbf{Invoke the Ind-completion theorem.}
    The finite-colimit closure of the finitely generated free commutative
    \(A\)-algebras has Ind-completion
    \[
        \operatorname{Ind}
        \bigl(\CAlg(\Sp)^{\mathrm{fc}}_{A/}\bigr)
        \simeq
        \CAlg(\Sp)^{\mathrm{cn}}_{A/},
    \]
    where \(\CAlg(\Sp)^{\mathrm{fc}}_{A/}\) denotes the full finite-cell
    subcategory of \Ref{def:finite-cell-spectral-algebra}. This is the
    connective commutative specialization of
    \cite[Definition~7.2.4.26 and Proposition~7.2.4.27]{LurieHigherAlgebra}.

    \item \textbf{Rewrite the theorem as a filtered-colimit presentation.}
    Equivalently, every connective commutative \(A\)-algebra is a filtered
    colimit of finite cell \(A\)-algebras.

    \item \textbf{Choose a finite-cell presentation of the compact algebra.}
    Let \(B\) be compact and choose
    \[
        B\simeq\varinjlim_\alpha F_\alpha
    \]
    with every \(F_\alpha\) finite cell.

    \item \textbf{Apply compactness to the identity mapping space.}
    Compactness gives
    \[
        \Map(B,B)
        \simeq
        \varinjlim_\alpha\Map(B,F_\alpha).
    \]

    \item \textbf{Represent the identity at a finite-cell stage.}
    Hence the point \(\id_B\) is represented at some stage by a morphism
    \(s:B\to F_\alpha\).

    \item \textbf{Construct the retraction.}
    If \(r:F_\alpha\to B\) is the colimit structure
    morphism, then
    \[
        r\circ s\simeq\id_B.
    \]

    \item \textbf{Conclude the retract statement.}
    Thus \(B\) is a retract of the finite cell algebra \(F_\alpha\).

    \item \textbf{Return to the stated definition.}
    By \Ref{prop:compactness-connective-undercategory}, compactness in the
    connective undercategory is equivalent to compactness in
    \(\CAlg(\Sp)_{A/}\), which is the local finite-presentation condition of
    \Ref{def:finitely-presented-spectral-algebra}.
\end{enumerate}
\end{proof}

\begin{proposition}[Permanence of finite presentation]
\label{prop:finite-presentation-permanence}
Morphisms locally of finite presentation are stable under equivalence,
composition, and base change.
\end{proposition}

\begin{proof}
\leavevmode
\begin{enumerate}[label=\textup{(\arabic*)}]
    \item \textbf{Equivalences.}
    Compactness is invariant under equivalence, giving the first permanence
    property.

    \item \textbf{Base change.}
    Base change
    \[
        \CAlg(\Sp)_{A/}
        \longrightarrow
        \CAlg(\Sp)_{A'/},
        \qquad
        B\longmapsto B\otimes_AA',
    \]
    is left adjoint to restriction of scalars. Restriction preserves filtered
    colimits, so the base-change functor preserves compact objects.

    \item \textbf{Composition.}
    Local finite presentation is stable under composition for morphisms of
    connective commutative ring spectra; this is the compactness permanence
    theorem for spectral algebra morphisms. See
    \cite{LurieHigherAlgebra,LurieSAG}. Applying it to
    \(A\to B\to C\) gives local finite presentation of \(A\to C\).
\end{enumerate}
\end{proof}

\subsection{Perfect and projective modules}

\begin{definition}[Perfect module]
\label{def:perfect-module-spectrum}
An $A$-module $M$ is \textbf{perfect} if it belongs to the thick
subcategory of $\Mod_A$ generated by the tensor unit $A$.
\end{definition}

\begin{proposition}[Perfect equals dualizable]
\label{prop:perfect-equals-dualizable}
An $A$-module is perfect if and only if it is dualizable in the symmetric
monoidal $\infty$-category $\Mod_A$.
\end{proposition}

\begin{proof}
\leavevmode
\begin{enumerate}[label=\textup{(\arabic*)}]
    \item \textbf{Place the unit in both classes.}
    The unit $A$ is compact and dualizable.

    \item \textbf{Close dualizable objects thickly.}
    The dualizable objects form a
    thick subcategory.

    \item \textbf{Show a dualizable module is compact.}
    Conversely, a dualizable $A$-module is compact because
    its mapping functor is represented after tensoring with its dual.

    \item \textbf{Identify the compact objects.}
    The
    compact objects of $\Mod_A$ form the thick subcategory generated by $A$.
    See \cite{LurieHigherAlgebra}.
\end{enumerate}
\end{proof}

\begin{definition}[Tor-amplitude]
\label{def:tor-amplitude}
An $A$-module $M$ has \textbf{Tor-amplitude in $[a,b]$} if, for every
discrete $A$-module $N$, the module
\[
    M\otimes_A N
\]
has homotopy groups only in degrees $a$ through $b$.
\end{definition}

\begin{definition}[Finite projective module]
\label{def:finite-projective-module-spectrum}
An $A$-module $M$ is \textbf{finite projective} if it is a retract of a
finite free $A$-module.
\end{definition}

\begin{proposition}[Characterization of finite projectives]
\label{prop:finite-projective-characterization}
For an $A$-module $M$, the following are equivalent.
\begin{enumerate}[label=(\roman*)]
    \item The module $M$ is finite projective.

    \item The module $M$ is perfect and has Tor-amplitude contained in
    $[0,0]$.

    \item The module $M$ is perfect, connective, and flat.
\end{enumerate}
\end{proposition}
\begin{proof}
\leavevmode
\begin{enumerate}[label=\textup{(\arabic*)}]
    \item \textbf{Check finite free modules.}
    A finite free \(A\)-module is perfect, connective, flat, and has
    Tor-amplitude in \([0,0]\).

    \item \textbf{Pass the properties to retracts.}
    These properties are stable under retracts.
    Hence \textup{(i)} implies both \textup{(ii)} and \textup{(iii)}.

    \item \textbf{Pass to the discrete base.}
    Assume \textup{(iii)} and put
    \[
        M_0:=M\otimes_AH\pi_0(A).
    \]

    \item \textbf{Compute the homotopy of the base change.}
    Base change preserves perfectness, and flatness gives
    \[
        \pi_n(M_0)
        \simeq
        \pi_n(H\pi_0(A))\otimes_{\pi_0(A)}\pi_0(M).
    \]

    \item \textbf{Identify \(\pi_0(M)\) as finitely generated projective.}
    Thus \(M_0\) is discrete with \(\pi_0(M_0)\simeq\pi_0(M)\). A discrete
    perfect module over \(H\pi_0(A)\) corresponds to a perfect ordinary
    \(\pi_0(A)\)-module, so \(\pi_0(M)\) is finitely presented. It is also
    flat by \Ref{thm:spectral-lazard}; hence it is finitely generated
    projective.

    \item \textbf{Choose an idempotent.}
    Choose an idempotent
    \[
        e\in M_r(\pi_0(A))
    \]
    with image \(\pi_0(M)\).

    \item \textbf{Split the idempotent.}
    The corresponding idempotent endomorphism of
    \(A^{\oplus r}\) splits in \(\Mod_A\), producing a finite projective
    module \(P\) with
    \[
        \pi_0(P)\simeq\pi_0(M).
    \]

    \item \textbf{Lift the degree-zero isomorphism.}
    Because \(P\) is a retract of a finite free module, the map
    \[
        \pi_0\Map_{\Mod_A}(P,M)
        \longrightarrow
        \Hom_{\pi_0(A)}\bigl(\pi_0(P),\pi_0(M)\bigr)
    \]
    is surjective. Lift a chosen isomorphism on \(\pi_0\) to
    \(u:P\to M\).

    \item \textbf{Compare all homotopy groups.}
    Both modules are flat, so
    \[
        \pi_n(u)
        \simeq
        \pi_n(A)\otimes_{\pi_0(A)}\pi_0(u)
    \]
    is an isomorphism for every \(n\).

    \item \textbf{Conclude finite projectivity.}
    Hence \(u\) is an equivalence and
    \(M\) is finite projective.

    \item \textbf{Pass again to the discrete base.}
    Assume \textup{(ii)} and again put
    \[
        M_0:=M\otimes_AH\pi_0(A).
    \]
    The Tor-amplitude condition makes \(M_0\) discrete. Perfectness over a
    connective ring spectrum makes \(M\) bounded below.

    \item \textbf{Choose the first nonzero homotopy degree.}
    Let
    \[
        I:=\operatorname{fib}\bigl(A\to H\pi_0(A)\bigr),
    \]
    so \(I\in(\Mod_{A})_{\geq1}\). If \(M\neq0\), let \(r\) be the least
    integer with \(\pi_r(M)\neq0\). Then
    \[
        I\otimes_AM\in(\Mod_{A})_{\geq r+1}.
    \]

    \item \textbf{Tensor the defining fibre sequence.}
    Tensoring the fibre sequence
    \[
        I\longrightarrow A\longrightarrow H\pi_0(A)
    \]
    with \(M\) gives
    \[
    \begin{tikzcd}[column sep=large]
        I\otimes_AM \arrow[r] & M \arrow[r] & M_0.
    \end{tikzcd}
    \]

    \item \textbf{Compare the first nonzero homotopy group.}
    Hence \(\pi_r(M)\to\pi_r(M_0)\) is an isomorphism.

    \item \textbf{Deduce connectivity.}
    Since \(M_0\) is
    discrete, \(r=0\). Therefore \(M\) is connective.

    \item \textbf{Define the tensor functor.}
    Let \(F=-\otimes_AM\). Connectivity of \(M\) makes \(F\) right
    \(t\)-exact, while Tor-amplitude \([0,0]\) makes \(F\) carry discrete
    modules to discrete modules.

    \item \textbf{Approximate a coconnective module by bounded truncations.}
    For \(N\in(\Mod_{A})_{\leq0}\), right
    completeness gives
    \[
        N\simeq\varinjlim_k\tau_{\geq-k}^{A}N.
    \]

    \item \textbf{Build each bounded truncation from discrete pieces.}
    Each bounded truncation is obtained by finitely many extensions from the
    shifted discrete modules \(H\pi_j(N)[j]\), \(-k\leq j\leq0\).

    \item \textbf{Apply exactness to the bounded truncations.}
    Exactness
    of \(F\) therefore gives
    \[
        F(\tau_{\geq-k}^{A}N)\in(\Mod_{A})_{\leq0}.
    \]

    \item \textbf{Pass to the filtered colimit.}
    Filtered colimits preserve coconnective objects, so
    \(F(N)\in(\Mod_{A})_{\leq0}\).

    \item \textbf{Conclude $t$-exactness.}
    Thus \(F\) is \(t\)-exact and \(M\) is flat.

    \item \textbf{Return to finite projectivity.}
    Step \textup{(10)} now gives \textup{(ii)}\(\Rightarrow\)\textup{(i)}.
\end{enumerate}
\end{proof}

At the affine level, the geometric conditions needed for locality now have
precise algebraic forms: flatness is Postnikov-exact substitution, local
finite presentation is compactness of the coordinate algebra, and finite
projectivity is perfect connective flatness. The next problem is scope: the coefficient system must be extended beyond
affines before locality can select the global geometric contexts.

\section{Generalized contexts and extension of the coefficient system}
\label{sec:spectral-prestacks-qcoh}

The coefficient system has so far been constructed on affines. The next step
is therefore to enlarge the base before imposing any geometric locality
condition, so that locality can be formulated globally. We pass from affine
contexts to space-valued prestacks and extend \(\QCoh_{\mathrm{aff}}\) by
right Kan extension along opposite Yoneda. The resulting coefficient fibres
exist on arbitrary generalized contexts, while unstraightening records the
same substitution data as Cartesian transport.

\subsection{Spectral prestacks and Yoneda}


\begin{definition}[Spectral prestacks]
\label{def:spectral-prestacks-affine-stage}
The $\infty$-category of space-valued presheaves on affine spectral
schemes is
\[
    \PSh(\Aff_{\Sp})
    :=
    \Fun
    \bigl(
        (\Aff_{\Sp})^{\op},
        \mathcal S
    \bigr).
\]
Its objects are called \textbf{spectral prestacks}.
\end{definition}

\begin{definition}[Functor of points]
\label{def:spectral-functor-of-points}
For $A\in\CAlg(\Sp)^{\mathrm{cn}}$, the functor of points of
$\Spec(A)$ is
\[
    h_A
    \colon
    (\Aff_{\Sp})^{\op}
    \longrightarrow
    \mathcal S,
\]
defined by
\[
    h_A(\Spec(B))
    :=
    \Map_{\CAlg(\Sp)}(A,B).
\]
\end{definition}

\begin{proposition}[Spectral Yoneda embedding]
\label{prop:spectral-yoneda-embedding}
The assignment
\[
    \Spec(A)
    \longmapsto
    h_A
\]
defines a fully faithful functor
\[
    h
    \colon
    \Aff_{\Sp}
    \hookrightarrow
    \PSh(\Aff_{\Sp}).
\]
\end{proposition}

\begin{proof}
\leavevmode
\begin{enumerate}[label=\textup{(\arabic*)}]
    \item \textbf{Apply Yoneda.}
    The $\infty$-categorical Yoneda lemma gives the asserted full faithfulness.
\end{enumerate}
\end{proof}

\subsection{Affine points and right Kan extension}

\begin{definition}[Affine points of a prestack]
\label{def:affine-points-prestack}
For a spectral prestack
\[
    X\in\PSh(\Aff_{\Sp}),
\]
let
\[
    \operatorname{AffPt}(X)
\]
denote the $\infty$-category of pairs
\[
    (\Spec(A),x),
    \qquad
    x\in X(A),
\]
equivalently the category of objects of \(\Aff_{\Sp}\) mapping to \(X\)
through the Yoneda embedding. This notation is used instead of
\(\Aff_{\Sp}/X\), since an arbitrary prestack \(X\) is not an object of
\(\Aff_{\Sp}\).
\end{definition}

\begin{definition}[Quasi-coherent modules on a spectral prestack]
\label{def:qcoh-spectral-prestack}
For a spectral prestack $X$, define
\[
    \QCoh(X)
    :=
    \varprojlim\nolimits_{(\Spec(A)\to X)
    \in\operatorname{AffPt}(X)^{\op}}
    \Mod_A.
\]
The limit is formed in
\[
    \CAlg(\Pr^L_{\mathrm{st}}).
\]
\end{definition}

The affine-point formula makes the generalized coefficient fibre explicit:
\[
\begin{tikzcd}[column sep=large]
    \operatorname{AffPt}(X)^{\op}
        \arrow[r]
    &
    \CAlg(\Pr^L_{\mathrm{st}})
    \\
    (\Spec(A)\to X)
        \arrow[r, mapsto]
    &
    \Mod_A.
\end{tikzcd}
\]

A morphism of generalized contexts \(X\to Y\) acts on affine points by
composition. The induced comparison of indexing diagrams is
\[
\begin{tikzcd}[column sep=huge, row sep=large]
    \operatorname{AffPt}(X)^{\op}
        \arrow[rr]
        \arrow[dd, "{\text{postcomposition}}"']
    &&
    \CAlg(\Pr^L_{\mathrm{st}})
    \\
    \\
    \operatorname{AffPt}(Y)^{\op}
        \arrow[uurr]
    &&
\end{tikzcd}
\]
and the resulting map on limits is the substitution functor
\(\QCoh(Y)\to\QCoh(X)\).

\begin{theorem}[Affine right Kan extension]
\label{thm:qcoh-affine-right-kan-extension}
Let
\[
    h
    \colon
    \Aff_{\Sp}
    \longrightarrow
    \PSh(\Aff_{\Sp})
\]
be the Yoneda embedding. The construction of
\Ref{def:qcoh-spectral-prestack} defines a functor
\[
    \QCoh
    \colon
    \PSh(\Aff_{\Sp})^{\op}
    \longrightarrow
    \CAlg(\Pr^L_{\mathrm{st}})
\]
which is the right Kan extension of
\[
    \QCoh_{\mathrm{aff}}
    \colon
    (\Aff_{\Sp})^{\op}
    \longrightarrow
    \CAlg(\Pr^L_{\mathrm{st}})
\]
along the opposite Yoneda embedding
\[
    h^{\op}
    \colon
    (\Aff_{\Sp})^{\op}
    \longrightarrow
    \PSh(\Aff_{\Sp})^{\op}.
\]

Equivalently, $\QCoh$ is the unique small-limit-preserving functor on
$\PSh(\Aff_{\Sp})^{\op}$ whose restriction along
$h^{\op}$ is $\QCoh_{\mathrm{aff}}$.
\end{theorem}

\begin{proof}
\leavevmode
\begin{enumerate}[label=\textup{(\arabic*)}]
    \item \textbf{Pointwise right Kan extension.}
    For a spectral prestack \(X\), the pointwise formula along
    \(h^{\op}\) is
    \[
        \operatorname{Ran}_{h^{\op}}(\QCoh_{\mathrm{aff}})(X)
        \simeq
        \varprojlim\nolimits_{(\Spec(A)\to X)\in\operatorname{AffPt}(X)^{\op}}\Mod_A.
    \]
    This is exactly \Ref{def:qcoh-spectral-prestack}.

    \item \textbf{Functoriality.}
    A morphism \(X\to Y\) induces
    \[
        \operatorname{AffPt}(X)\longrightarrow\operatorname{AffPt}(Y)
    \]
    by composition. Passing to opposites and using the universal property of
    the two limits produces
    \[
        \QCoh(Y)\longrightarrow\QCoh(X).
    \]

    \item \textbf{Universal property.}
    The Yoneda embedding exhibits \(\PSh(\Aff_{\Sp})\) as the free
    cocompletion of \(\Aff_{\Sp}\); after passage to opposites,
    \(\PSh(\Aff_{\Sp})^{\op}\) is generated under limits by the image of
    \(h^{\op}\). Hence a small-limit-preserving extension is uniquely determined
    by its restriction to representables, and the right Kan extension has
    the asserted universal property.
\end{enumerate}
\end{proof}

\begin{corollary}[Recovery on affine spectral schemes]
\label{cor:qcoh-recovery-affines}
For every connective commutative ring spectrum $A$, there is a canonical
symmetric monoidal equivalence
\[
    \QCoh(\Spec(A))
    \simeq
    \Mod_A.
\]
\end{corollary}

\begin{proof}
\leavevmode
\begin{enumerate}[label=\textup{(\arabic*)}]
    \item \textbf{Identify a terminal affine point.}
    For the representable prestack
    \[
        X=\Spec(A),
    \]
    the identity map
    \[
        \id_X
        \colon
        X
        \longrightarrow
        X
    \]
    is terminal in the affine-point category
    \[
        \operatorname{AffPt}(X).
    \]

    \item \textbf{Pass to the opposite indexing category.}
    It is therefore initial in
    \[
        \operatorname{AffPt}(X)^{\op}.
    \]

    \item \textbf{Evaluate the limit at the initial object.}
    A limit of a diagram indexed by an $\infty$-category with an initial
    object is canonically its value at that initial object. Hence
    \[
        \QCoh(\Spec(A))
        \simeq
        \Mod_A.
    \]

    \item \textbf{Preserve the symmetric monoidal structure.}
    The equivalence is symmetric monoidal because the defining limit is
    formed in
    \[
        \CAlg(\Pr^L_{\mathrm{st}}).
    \]
\end{enumerate}
\end{proof}

\begin{corollary}[Colimits of prestacks]
\label{cor:qcoh-colimits-to-limits}
For every small diagram of spectral prestacks $X_\alpha$, there is a
canonical equivalence
\[
    \QCoh
    \left(
        \varinjlim_\alpha X_\alpha
    \right)
    \simeq
    \varprojlim\nolimits_\alpha\QCoh(X_\alpha).
\]
\end{corollary}

\begin{proof}
\leavevmode
\begin{enumerate}[label=\textup{(\arabic*)}]
    \item \textbf{Apply limit preservation.}
    Apply the limit-preservation assertion of
    \Ref{thm:qcoh-affine-right-kan-extension} to the opposite functor.
\end{enumerate}
\end{proof}

\subsection{The coefficient fibration}

\begin{construction}[Grothendieck construction of quasi-coherent coefficients]
\label{con:qcoh-coefficient-fibration}
Forget the monoidal structure of the right-Kan-extended functor and regard
\[
    \QCoh
    :
    \PSh(\Aff_{\Sp})^{\op}
    \longrightarrow
    \Cat_\infty.
\]
Unstraightening produces a Cartesian fibration
\[
    \pi_{\QCoh}
    :
    \int\QCoh
    \longrightarrow
    \PSh(\Aff_{\Sp}).
\]
The fibre over a generalized context \(X\) is \(\QCoh(X)\). For a
morphism of contexts \(f:Y\to X\) and an object
\(\mathcal F\in\QCoh(X)\), Cartesian transport along \(f\) has domain
\[
    (Y,f^*\mathcal F)
\]
and codomain \((X,\mathcal F)\). The coherence of iterated substitution is
therefore the Cartesian-composition coherence of this fibration, rather than
additional structure imposed on individual pullback functors.
\end{construction}

\[
\begin{tikzcd}[column sep=large, row sep=large]
    (Y,f^*\mathcal F)
        \arrow[r]
        \arrow[d]
    &
    (X,\mathcal F)
        \arrow[d]
    \\
    Y
        \arrow[r, "f"']
    &
    X.
\end{tikzcd}
\]


\begin{remark}[Lawvere-style indexed semantics]
\label{rem:lawvere-style-indexed-semantics}
The Cartesian fibration \(\pi_{\QCoh}\) is the indexed-categorical object
encoded by the language of coefficients in context: its base consists of
contexts, its fibres are coefficient categories, and its Cartesian transport
is substitution. The right adjoints to substitution will be obtained
functorially by passage to right adjoints, after which the canonical mate
transformations are constructed from the resulting adjunctions. We call this
indexed organization a \textbf{Lawvere-style doctrine}. The term
\textbf{hyperdoctrine} is reserved for settings in which the additional
dependent adjoints exist and the relevant exchange transformations have been
proved invertible. Compare \cite{JacobsCategoricalLogic}.
\end{remark}

\subsection{Pullback, direct image, and the structure sheaf}

\begin{proposition}[Quasi-coherent pullback]
\label{prop:qcoh-pullback}
Every morphism of spectral prestacks
\[
    f
    \colon
    X
    \longrightarrow
    Y
\]
induces a colimit-preserving strong symmetric monoidal functor
\[
    f^*
    \colon
    \QCoh(Y)
    \longrightarrow
    \QCoh(X).
\]
For affine spectral schemes it agrees with extension of scalars.
\end{proposition}

\begin{proof}
\leavevmode
\begin{enumerate}[label=\textup{(\arabic*)}]
    \item \textbf{Obtain the transition functor.}
    Functoriality of the right Kan extension gives the transition functor.

    \item \textbf{Read off its monoidal properties.}
    The target of the right Kan extension is
    $\CAlg(\Pr^L_{\mathrm{st}})$, so the transition is colimit-preserving and
    strong symmetric monoidal.

    \item \textbf{Recover the affine case.}
    The affine identification follows from
    \Ref{cor:qcoh-recovery-affines} and the defining affine coefficient
    system.
\end{enumerate}
\end{proof}

\begin{corollary}[Quasi-coherent direct image]
\label{cor:qcoh-direct-image}
The pullback functor of \Ref{prop:qcoh-pullback} admits a right adjoint
\[
    f_*
    \colon
    \QCoh(X)
    \longrightarrow
    \QCoh(Y).
\]
\end{corollary}

\begin{proof}
\leavevmode
\begin{enumerate}[label=\textup{(\arabic*)}]
    \item \textbf{Preserve small colimits.}
    The functor $f^*$ preserves small colimits between presentable
    $\infty$-categories.

    \item \textbf{Apply the adjoint functor theorem.}
    The adjoint functor theorem supplies a right adjoint.
\end{enumerate}
\end{proof}

\begin{proposition}[Coherent direct-image coefficient system]
\label{prop:qcoh-coherent-direct-image}
Forget the symmetric monoidal enhancement of quasi-coherent pullback and write
\[
    \QCoh^*
    :
    \PSh(\Aff_{\Sp})^{\op}
    \longrightarrow
    \Pr^L_{\mathrm{st}}.
\]
Passage to right adjoints defines a functor
\[
    \QCoh_*
    :
    \PSh(\Aff_{\Sp})
    \longrightarrow
    \Pr^R_{\mathrm{st}}
\]
whose value on a morphism
\(f:X\to Y\) is the right adjoint
\[
    f_*:\QCoh(X)\longrightarrow\QCoh(Y)
\]
of \(f^*\). In particular, the identity and composition equivalences
\[
    (\id_X)_*\simeq\id_{\QCoh(X)},
    \qquad
    (g\circ f)_*\simeq g_*f_*
\]
and all their higher associativity coherences are supplied by the functor
\(\QCoh_*\); no coherent family of adjoints is chosen independently.
\end{proposition}

\begin{proof}
\leavevmode
\begin{enumerate}[label=\textup{(\arabic*)}]
    \item \textbf{Pass functorially to right adjoints.}
    Taking opposites gives
    \[
    \begin{tikzcd}[column sep=huge]
        \PSh(\Aff_{\Sp})
            \arrow[r, "{(\QCoh^*)^{\op}}"]
        &
        (\Pr^L_{\mathrm{st}})^{\op}
            \arrow[r, "{\operatorname{RAdj}}", "\simeq"']
        &
        \Pr^R_{\mathrm{st}}.
    \end{tikzcd}
    \]
    Define the composite to be \(\QCoh_*\). The value of
    \(\operatorname{RAdj}\) on a colimit-preserving functor is its right
    adjoint, so the transition attached to \(f:X\to Y\) is exactly the
    right adjoint of
    \[
        f^*:\QCoh(Y)\longrightarrow\QCoh(X).
    \]
    Existence of this right adjoint is
    \Ref{cor:qcoh-direct-image}.

    \item \textbf{Read off higher coherence.}
    For composable morphisms
    \(X\xrightarrow{f}Y\xrightarrow{g}Z\), the two adjunction directions may
    be displayed simultaneously as
    \[
    \begin{tikzcd}[column sep=huge, row sep=large]
        \QCoh(Z)
            \arrow[r, "g^*", shift left=1.1ex]
        &
        \QCoh(Y)
            \arrow[l, swap, "g_*", shift left=1.1ex]
            \arrow[r, "f^*", shift left=1.1ex]
        &
        \QCoh(X)
            \arrow[l, swap, "f_*", shift left=1.1ex].
    \end{tikzcd}
    \]
    Functoriality of \(\QCoh_*\) supplies
    \((g\circ f)_*\simeq g_*f_*\), while its higher functoriality supplies
    the unit and associativity homotopies. This is the coherent structure
    used below. Compare the standard passage-to-adjoints formalism in
    \cite{LurieHigherAlgebra}.
\end{enumerate}
\end{proof}

\begin{proposition}[Canonical composition, base-change, and projection morphisms]
\label{prop:qcoh-canonical-mates}
The functorial right adjoints of \Ref{prop:qcoh-coherent-direct-image} carry the following
canonical mate transformations.
\begin{enumerate}[label=(\roman*)]
    \item For composable morphisms
    \[
        X\xrightarrow{f}Y\xrightarrow{g}Z,
    \]
    there is a canonical coherent equivalence
    \[
        (g\circ f)_*
        \simeq
        g_*f_*.
    \]

    \item For every Cartesian square of spectral prestacks
    \[
    \begin{tikzcd}[column sep=large, row sep=large]
        X' \arrow[r, "g'"] \arrow[d, "f'"']
            & X \arrow[d, "f"] \\
        Y' \arrow[r, "g"']
            & Y,
    \end{tikzcd}
    \]
    there is a canonical base-change morphism
    \[
        \operatorname{BC}_{f,g}:
        g^*f_*
        \longrightarrow
        f'_*g'^*.
    \]

    \item For \(\mathcal E\in\QCoh(Y)\) and
    \(\mathcal F\in\QCoh(X)\), there is a canonical projection morphism
    \[
        \operatorname{PF}_{f}(\mathcal E,\mathcal F):
        \mathcal E\otimes f_*\mathcal F
        \longrightarrow
        f_*(f^*\mathcal E\otimes\mathcal F).
    \]
\end{enumerate}
No equivalence assertion is part of the construction in \textup{(ii)} or
\textup{(iii)}.
\end{proposition}

\begin{proof}
\leavevmode
\begin{enumerate}[label=\textup{(\arabic*)}]
    \item \textbf{Composition from the direct-image functor.}
    The equivalence
    \[
        (g\circ f)_*\simeq g_*f_*
    \]
    and all of its higher associativity coherence are the composition data of
    the functor \(\QCoh_*\) constructed in
    \Ref{prop:qcoh-coherent-direct-image}. Thus no pointwise choice of
    coherent right adjoints is required.

    \item \textbf{Base change as an adjoint mate.}
    The Cartesian square supplies the canonical pullback equivalence
    \[
        g'^*f^*
        \simeq
        f'^*g^*.
    \]
    For \(\mathcal F\in\QCoh(X)\), define
    \(\operatorname{BC}_{f,g}(\mathcal F)\) to be the composite
    \[
    \begin{aligned}
        g^*f_*\mathcal F
        &\longrightarrow
        f'_*f'^*g^*f_*\mathcal F
        \\
        &\simeq
        f'_*g'^*f^*f_*\mathcal F
        \\
        &\longrightarrow
        f'_*g'^*\mathcal F,
    \end{aligned}
    \]
    where the first arrow is the unit of \(f'^*\dashv f'_*\), the middle
    equivalence is the pullback comparison, and the final arrow is induced by
    the counit of \(f^*\dashv f_*\). This is precisely the adjoint mate of
    the pullback equivalence.

    \item \textbf{Projection morphism as an adjoint mate.}
    Define the displayed projection morphism to be adjoint to
    \[
    \begin{aligned}
        f^*(\mathcal E\otimes f_*\mathcal F)
        &\simeq
        f^*\mathcal E\otimes f^*f_*\mathcal F
        \\
        &\longrightarrow
        f^*\mathcal E\otimes\mathcal F,
    \end{aligned}
    \]
    where the final morphism uses the counit. Naturality and all coherence
    follow from the calculus of adjoint mates and the coherent monoidality of
    pullback.
\end{enumerate}
\end{proof}

\begin{remark}[No general dependent left adjoint]
\label{rem:no-general-dependent-left-adjoint-qcoh}
The indexed coefficient doctrine does not supply a left adjoint to every
substitution functor. A concrete obstruction already appears for extension
of scalars
\[
    -\otimes_{H\mathbb Z}H\mathbb Q:
    \Mod_{H\mathbb Z}
    \longrightarrow
    \Mod_{H\mathbb Q}.
\]
This functor does not preserve infinite products. On degree-zero homotopy,
the canonical map
\[
    \left(\prod_{n\geq1}\mathbb Z\right)\otimes_{\mathbb Z}\mathbb Q
    \longrightarrow
    \prod_{n\geq1}\mathbb Q
\]
is not surjective: elements on the left have a common denominator, whereas
\[
    \left(1,\frac12,\frac13,\frac14,\ldots\right)
\]
does not. Thus this substitution functor is not a right adjoint and cannot
itself admit a left adjoint. Any dependent left adjoints considered in a
more specialized geometric setting must therefore come from additional
geometry rather than from the indexed doctrine alone.
\end{remark}

\begin{definition}[Structure sheaf]
\label{def:structure-sheaf-prestack}
The \textbf{structure sheaf} of a spectral prestack $X$ is the tensor unit
\[
    \mathcal O_X
    \in
    \QCoh(X).
\]
\end{definition}

\begin{definition}[Global sections]
\label{def:global-sections-qcoh}
For $\mathcal F\in\QCoh(X)$, define its spectrum of global sections by
\[
    \Gamma(X,\mathcal F)
    :=
    \map_{\QCoh(X)}
    (\mathcal O_X,\mathcal F).
\]
\end{definition}

\begin{proposition}[Pullback of the structure sheaf]
\label{prop:pullback-structure-sheaf}
For every morphism $f:X\to Y$, there is a canonical equivalence
\[
    f^*\mathcal O_Y
    \simeq
    \mathcal O_X.
\]
\end{proposition}

\begin{proof}
\leavevmode
\begin{enumerate}[label=\textup{(\arabic*)}]
    \item \textbf{Preserve the tensor unit.}
    The functor $f^*$ is strong symmetric monoidal and therefore preserves the
    tensor unit.
\end{enumerate}
\end{proof}

\subsection{Connectivity on prestacks}

\begin{definition}[Objectwise connective quasi-coherent module]
\label{def:objectwise-connective-qcoh}
A quasi-coherent module $\mathcal F\in\QCoh(X)$ is
\textbf{objectwise connective} if, for every affine point
\[
    x
    \colon
    \Spec(A)
    \longrightarrow
    X,
\]
its pullback
\[
    x^*\mathcal F
\]
belongs to $(\Mod_{A})_{\geq0}$. Objectwise coconnective modules are defined
similarly.
\end{definition}

The affine coefficient system has now been extended to every spectral
prestack, with coherent substitution encoded either functorially or by the
associated Cartesian fibration. The remaining step toward geometry is
locality: we must identify which generalized contexts can be reconstructed
from affine opens, and show that the coefficient system reconstructs with
them.

\section{Zariski locality and coherent descent}
\label{sec:affine-opens-zariski-descent}

The generalized coefficient system is available before geometry has been
selected. Zariski locality now performs that selection. We first characterize
affine opens and principal spectral localizations, then define covering
families by joint conservativity of substitution. Descendability and the
Amitsur complex convert this local detection criterion into reconstruction of
representables, modules, and commutative algebras.

\begin{convention}[Sheaves and hyperdescent]
\label{conv:spectral-sheaves-hyperdescent}
Unless explicitly stated otherwise,
\[
    \Sh_\tau(\mathcal C)
\]
denotes the $\infty$-category of $\tau$-sheaves of spaces without
hypercompletion. When hyperdescent is required, the hypercomplete
$\infty$-topos is denoted by
\[
    \Sh_\tau(\mathcal C)^\wedge.
\]
\end{convention}

\begin{convention}[Effective epimorphisms and descent data]
\label{conv:spectral-effective-epimorphisms}
In an \(\infty\)-topos, a morphism
\[
    p:U\longrightarrow X
\]
is an \textbf{effective epimorphism} if the augmentation from its \v{C}ech
nerve \(U_\bullet\to X\) exhibits \(X\) as the geometric realization
\[
    X
    \simeq
    \operatorname*{colim}_{[n]\in\Delta^{\op}}U_n.
\]
Whenever descent data for a cosimplicial diagram of \(\infty\)-categories
\(\mathcal C^\bullet\) are mentioned, they mean an object of the totalization
\[
    \operatorname*{Tot}_{[n]\in\Delta}\mathcal C^n.
\]
Thus the compatibility morphisms and all higher cocycle coherences are part of
the displayed limit rather than additional axiomatic data. We also use the
standard local-surjectivity characterization of effective epimorphisms in a
sheaf \(\infty\)-topos: sections lift after passage to a covering family; see
\cite{LurieSAG}.
\end{convention}

\begin{lemma}[Mapping descent along an effective epimorphism]
\label{lem:mapping-descent-effective-epimorphism}
Let \(\mathcal X\) be an \(\infty\)-topos, let
\[
    p:U\longrightarrow X
\]
be an effective epimorphism with \v{C}ech nerve \(U_\bullet\), and let
\[
    Y\longrightarrow X,
    \qquad
    Z\longrightarrow X
\]
be objects of the slice \(\mathcal X_{/X}\). Then there is a natural
equivalence
\[
    \Map_X(Y,Z)
    \simeq
    \operatorname*{Tot}_{[n]\in\Delta}
    \Map_{U_n}
    \bigl(
        Y\times_XU_n,
        Z\times_XU_n
    \bigr).
\]
In particular, morphisms and all homotopies between them may be reconstructed
from their restrictions to an effective epimorphic cover and its iterated
intersections.
\end{lemma}

\begin{proof}
\leavevmode
\begin{enumerate}[label=\textup{(\arabic*)}]
    \item \textbf{Base change the effective epimorphism.}
    Effective epimorphisms are stable under base change in an
    \(\infty\)-topos. Hence
    \[
        Y\times_XU\longrightarrow Y
    \]
    is an effective epimorphism in the slice \(\mathcal X_{/X}\), with
    \v{C}ech nerve
    \[
        [n]\longmapsto Y\times_XU_n.
    \]
    Therefore
    \[
        Y
        \simeq
        \operatorname*{colim}_{[n]\in\Delta^{\op}}
        Y\times_XU_n
    \]
    in \(\mathcal X_{/X}\).

    \item \textbf{Map out of the geometric realization.}
    Mapping from the preceding colimit into \(Z\) gives
    \[
        \Map_X(Y,Z)
        \simeq
        \operatorname*{Tot}_{[n]\in\Delta}
        \Map_X(Y\times_XU_n,Z).
    \]
    Since \(Y\times_XU_n\) lies over \(U_n\), the universal property of the
    slice identifies
    \[
        \Map_X(Y\times_XU_n,Z)
        \simeq
        \Map_{U_n}
        \bigl(
            Y\times_XU_n,
            Z\times_XU_n
        \bigr).
    \]
    Combining the displayed equivalences proves the assertion.
\end{enumerate}
\end{proof}

\subsection{Affine monomorphisms and open immersions}

\begin{proposition}[Affine monomorphism criterion]
\label{prop:affine-monomorphism-criterion}
Let $A\to B$ be a morphism of connective commutative ring spectra. The
following are equivalent.
\begin{enumerate}[label=(\roman*)]
    \item The morphism
    \[
        \Spec(B)
        \longrightarrow
        \Spec(A)
    \]
    is a monomorphism of spectral prestacks.

    \item The multiplication morphism
    \[
        B\otimes_A B
        \longrightarrow
        B
    \]
    is an equivalence.
\end{enumerate}
\end{proposition}

\begin{proof}
\leavevmode
\begin{enumerate}[label=\textup{(\arabic*)}]
    \item \textbf{Use the diagonal criterion.}
    A morphism $X\to Y$ is a monomorphism if and only if its diagonal is an
    equivalence.

    \item \textbf{Identify the affine diagonal.}
    By \Ref{prop:affine-spectral-fibre-products}, the diagonal of
    $\Spec(B)\to\Spec(A)$ is represented by the multiplication map
    $B\otimes_A B\to B$.
\end{enumerate}
\end{proof}

\begin{definition}[Affine open immersion]
\label{def:affine-open-immersion-spectral}
A morphism
\[
    \Spec(B)
    \longrightarrow
    \Spec(A)
\]
of affine spectral schemes is an \textbf{affine open immersion} if the
algebra map $A\to B$ is flat, locally of finite presentation, and satisfies
\[
    B\otimes_A B
    \xrightarrow{\simeq}
    B.
\]
\end{definition}

\begin{proposition}[Permanence of affine open immersions]
\label{prop:affine-open-permanence}
Affine open immersions are stable under equivalence, composition, and base
change.
\end{proposition}

\begin{proof}
\leavevmode
\begin{enumerate}[label=\textup{(\arabic*)}]
    \item \textbf{Preserve flatness and finite presentation.}
    Flatness and finite presentation have the required permanence properties
    by \Ref{prop:flat-algebra-permanence} and
    \Ref{prop:finite-presentation-permanence}.

    \item \textbf{Preserve monomorphisms.}
    Monomorphisms are stable under
    equivalence, composition, and base change.

    \item \textbf{Apply the definition.}
    Apply
    \Ref{def:affine-open-immersion-spectral}.
\end{enumerate}
\end{proof}

\subsection{Principal opens}

\begin{construction}[Spectral localization]
\label{con:spectral-localization}
Let $A$ be connective and let
\[
    f\in\pi_0(A).
\]
The \textbf{localization of $A$ at $f$} is the initial commutative
$A$-algebra
\[
    A[f^{-1}]
\]
in which the image of $f$ is invertible.
\end{construction}

\begin{proposition}[Homotopy groups of a localization]
\label{prop:homotopy-groups-localization}
For every integer $n$, there is a natural isomorphism
\[
    \pi_n(A[f^{-1}])
    \simeq
    \pi_n(A)[f^{-1}].
\]
In particular, $A\to A[f^{-1}]$ is flat.
\end{proposition}

\begin{proof}
\leavevmode
\begin{enumerate}[label=\textup{(\arabic*)}]
    \item \textbf{Underlying module of the localization.}
    Choose a representative of the homotopy class of multiplication by
    \(f\) in \(\Map_{\Mod_A}(A,A)\). The spectral localization theorem
    identifies the underlying \(A\)-module of \(A[f^{-1}]\) with the
    telescope
    \[
    \begin{tikzcd}[column sep=large]
        A \arrow[r, "f"] & A \arrow[r, "f"] & A \arrow[r, "f"] & \cdots.
    \end{tikzcd}
    \]
    Any representative of the same degree-zero class computes the same
    localization object by its universal property, so this proof choice is not
    part of the construction. See \cite{LurieHigherAlgebra}.

    \item \textbf{Homotopy groups.}
    Homotopy groups commute with filtered colimits, so
    \[
    \begin{aligned}
        \pi_n(A[f^{-1}])
        &\simeq
        \varinjlim
        \bigl(
            \pi_n(A)\xrightarrow{f}\pi_n(A)\xrightarrow{f}\cdots
        \bigr)
        \\
        &\simeq
        \pi_n(A)[f^{-1}].
    \end{aligned}
    \]

    \item \textbf{Flatness.}
    In degree zero this gives
    \(\pi_0(A[f^{-1}])\simeq\pi_0(A)[f^{-1}]\), a flat
    \(\pi_0(A)\)-algebra. Together with the displayed homotopy formula,
    \Ref{thm:spectral-lazard} implies that \(A\to A[f^{-1}]\) is flat.
\end{enumerate}
\end{proof}

\begin{proposition}[Principal spectral opens]
\label{prop:principal-spectral-open}
The morphism
\[
    D_A^{\mathrm{sp}}(f)
    :=
    \Spec(A[f^{-1}])
    \longrightarrow
    \Spec(A)
\]
is an affine open immersion.
\end{proposition}

\begin{proof}
\leavevmode
\begin{enumerate}[label=\textup{(\arabic*)}]
    \item \textbf{Flatness.}
    The localization map is flat by
    \Ref{prop:homotopy-groups-localization}.

    \item \textbf{Choose the universal polynomial algebra.}
    Let
    \[
        P:=\operatorname{Sym}_{\mathbb S}(\mathbb S)
    \]
    with universal generator \(t\).

    \item \textbf{Choose a representative of the class.}
    Choose a point
    \(\widetilde f\in\Omega^\infty A\) representing the component
    \(f\in\pi_0(A)\), and use \(t\mapsto\widetilde f\) to regard \(A\)
    as a \(P\)-algebra for this proof.

    \item \textbf{Introduce the two free $P$-algebras.}
    In \(\CAlg(\Mod_P)\), let
    \[
        R:=\operatorname{Sym}_P(P),
        \qquad
        Q:=\operatorname{Sym}_P(P),
    \]
    with generators \(u\) and \(z\), respectively.

    \item \textbf{Form the pushout.}
    Form
    \[
    \begin{tikzcd}[column sep=large, row sep=large]
        Q
            \arrow[r, "{z\mapsto tu-1}"]
            \arrow[d, "{z\mapsto0}"']
        &
        R
            \arrow[d]
        \\
        P
            \arrow[r]
        &
        P[t^{-1}].
    \end{tikzcd}
    \]

    \item \textbf{Describe maps out of the pushout.}
    A map from this pushout to a commutative \(P\)-algebra \(C\) is a choice
    of \(u\in\Omega^\infty C\) together with a path \(tu\simeq1\).

    \item \textbf{Use invertibility of $t$.}
    If the
    image of \(t\) is a unit, multiplication by \(t\) is an equivalence of
    \(C\)-modules, so
    \[
        \Omega^\infty C\xrightarrow{t\cdot-}\Omega^\infty C
    \]
    is an equivalence of spaces and the homotopy fibre over \(1\) is
    contractible.

    \item \textbf{Identify the universal localization.}
    The pushout therefore has the universal property of
    \(P[t^{-1}]\).

    \item \textbf{Identify the finite-cell generators.}
    The algebras \(Q\) and \(R\) are finitely generated free commutative
    \(P\)-algebras, while \(P\) is the free \(P\)-algebra on the zero
    module.

    \item \textbf{Conclude finite-cell structure.}
    The displayed pushout is therefore a finite colimit of the
    generators appearing in \Ref{def:finite-cell-spectral-algebra}.

    \item \textbf{Apply the finite-cell criterion.}
    Hence
    \(P[t^{-1}]\) is a finite cell \(P\)-algebra and is locally of finite
    presentation by \Ref{prop:finite-presentation-compactness-criterion}.

    \item \textbf{Base change the universal localization.}
    For the chosen representative \(\widetilde f\), the base change
    \(A\otimes_PP[t^{-1}]\) has the universal property of an
    \(A\)-algebra in which the degree-zero class \(f\) is invertible.

    \item \textbf{Identify it with \(A[f^{-1}]\).}
    Hence the localization universal property gives a canonical equivalence
    \[
        A[f^{-1}]
        \simeq
        A\otimes_PP[t^{-1}].
    \]

    \item \textbf{Remove dependence on the representative.}
    In particular, the object on the left and the conclusion below are
    independent of the chosen representative of \(f\).

    \item \textbf{Preserve local finite presentation.}
    Local finite presentation is preserved by base change, so
    \(A\to A[f^{-1}]\) is locally of finite presentation.

    \item \textbf{Verify the multiplication equivalence.}
    The multiplication
    map
    \[
        A[f^{-1}]\otimes_AA[f^{-1}]
        \longrightarrow
        A[f^{-1}]
    \]
    is an equivalence by the same universal property.

    \item \textbf{Conclude that the principal localization is open.}
    Thus \(D_A^{\mathrm{sp}}(f)\to\Spec A\)
    is an affine open immersion.
\end{enumerate}
\end{proof}

\begin{proposition}[Principal opens form a basis]
\label{prop:principal-opens-basis}
Every affine open immersion
\[
    \Spec(B)\longrightarrow\Spec(A)
\]
is covered by principal spectral opens \(D_A^{\mathrm{sp}}(f)\) with
\(f\in\pi_0(A)\). More precisely, there is a family of elements
\(\{f_\alpha\}\subseteq\pi_0(A)\) such that every
\(D_A^{\mathrm{sp}}(f_\alpha)\to\Spec(A)\) factors through
\(\Spec(B)\), and after pullback to \(\Spec(B)\) these principal opens are
jointly conservative on \(\Mod_B\).
\end{proposition}

\[
\begin{tikzcd}[column sep=large, row sep=large]
    D_A^{\mathrm{sp}}(f_\alpha f_\beta)
        \arrow[r]
        \arrow[d]
    &
    D_A^{\mathrm{sp}}(f_\beta)
        \arrow[d]
    \\
    D_A^{\mathrm{sp}}(f_\alpha)
        \arrow[r]
    &
    \Spec(A).
\end{tikzcd}
\]


\begin{proof}
\leavevmode
\begin{enumerate}[label=\textup{(\arabic*)}]
    \item \textbf{Pass flatness to degree zero.}
    Let \(A\to B\) represent the affine open immersion. Flatness gives a
    flat ring map \(\pi_0(A)\to\pi_0(B)\).

    \item \textbf{Pass the diagonal equivalence to degree zero.}
    Applying \(\pi_0\) to the
    diagonal equivalence
    \[
        B\otimes_AB\xrightarrow{\simeq}B
    \]
    and using \Ref{prop:classical-truncation-affine-base-change} shows that
    \(\pi_0(A)\to\pi_0(B)\) is a ring epimorphism.

    \item \textbf{Set up finite presentation in degree zero.}
    Local finite presentation of \(A\to B\) also descends to degree zero.
    Let \(\{R_\lambda\}\) be a filtered diagram of
    \(\pi_0(A)\)-algebras.

    \item \textbf{Commute Eilenberg--Mac Lane with the filtered colimit.}
    Filtered colimits of discrete connective
    commutative ring spectra are created on underlying spectra and remain
    discrete, so there is a canonical equivalence
    \[
        H\!\left(\varinjlim_\lambda R_\lambda\right)
        \simeq
        \varinjlim_\lambda HR_\lambda.
    \]

    \item \textbf{Use compactness to compare mapping spaces.}
    Compactness of \(B\) and \Ref{prop:pi-zero-h-adjunction} therefore give
    \[
    \begin{aligned}
        \Map_{\operatorname{CRing}_{\pi_0(A)/}}
        \left(\pi_0(B),\varinjlim_\lambda R_\lambda\right)
        &\simeq
        \Map_{\CAlg(\Sp)_{A/}}
        \left(B,H\!\left(\varinjlim_\lambda R_\lambda\right)\right)
        \\
        &\simeq
        \Map_{\CAlg(\Sp)_{A/}}
        \left(B,\varinjlim_\lambda HR_\lambda\right)
        \\
        &\simeq
        \varinjlim_\lambda
        \Map_{\CAlg(\Sp)_{A/}}(B,HR_\lambda)
        \\
        &\simeq
        \varinjlim_\lambda
        \Map_{\operatorname{CRing}_{\pi_0(A)/}}
        (\pi_0(B),R_\lambda).
    \end{aligned}
    \]

    \item \textbf{Conclude finite presentation in degree zero.}
    Hence \(\pi_0(B)\) is finitely presented over \(\pi_0(A)\).

    \item \textbf{Identify the resulting ordinary open immersion.}
    A flat
    finitely presented ring epimorphism determines an open immersion on prime
    spectra; equivalently, its image is locally principal. See
    \cite{LurieSAG}.

    \item \textbf{Choose a principal neighbourhood in degree zero.}
    Let \(\mathfrak q\in\Spec\pi_0(B)\) and let
    \(\mathfrak p\in\Spec\pi_0(A)\) be its image. Choose
    \(f\in\pi_0(A)\) such that \(\mathfrak p\in D(f)\) and the principal
    open \(D(f)\) is contained in the image of
    \(\Spec\pi_0(B)\to\Spec\pi_0(A)\).

    \item \textbf{Form the spectral pullback.}
    Form
    \[
    \begin{tikzcd}[column sep=large, row sep=large]
        P:=\Spec(B)\times_{\Spec(A)}\Spec(A[f^{-1}])
            \arrow[r]
            \arrow[d, "p"']
        &
        \Spec(B) \arrow[d]
        \\
        \Spec(A[f^{-1}]) \arrow[r]
        &
        \Spec(A).
    \end{tikzcd}
    \]

    \item \textbf{Identify its affine coordinate algebra.}
    The object \(P\) is affine and \(p\) is an affine open immersion by
    base change. Its coordinate algebra is
    \(B\otimes_AA[f^{-1}]\).

    \item \textbf{Compute its degree-zero ring.}
    By
    \Ref{prop:classical-truncation-affine-base-change}, its degree-zero ring
    is
    \[
        \pi_0(B)\otimes_{\pi_0(A)}\pi_0(A)[f^{-1}].
    \]

    \item \textbf{Use the chosen principal neighbourhood.}
    By the choice of the principal neighbourhood, this ring is canonically
    \(\pi_0(A)[f^{-1}]\).

    \item \textbf{Analyze the comparison algebra map.}
    The map
    \[
        A[f^{-1}]
        \longrightarrow
        B\otimes_AA[f^{-1}]
    \]
    is flat and induces an isomorphism on \(\pi_0\).

    \item \textbf{Detect its equivalence on homotopy groups.}
    The flat
    homotopy-group formula therefore makes it an isomorphism on every
    homotopy group, hence an equivalence.

    \item \textbf{Obtain the spectral factorization.}
    Thus
    \[
        P\xrightarrow{\simeq}\Spec(A[f^{-1}]),
    \]
    and the inverse exhibits \(D_A^{\mathrm{sp}}(f)\) as an open spectral
    subscheme factoring through \(\Spec(B)\).

    \item \textbf{Repeat the construction over every prime.}
    Repeat the construction for every prime of \(\pi_0(B)\). After pullback
    to \(\Spec(B)\), the resulting principal localizations cover the prime
    spectrum of \(\pi_0(B)\).

    \item \textbf{Detect vanishing on the principal cover.}
    If a \(B\)-module restricts to zero on all of
    them, every homotopy group localizes to zero on this principal cover and
    hence vanishes.

    \item \textbf{Conclude joint conservativity.}
    Homotopy groups detect equivalences, so the pullback
    functors are jointly conservative on \(\Mod_B\).
\end{enumerate}
\end{proof}

\subsection{The affine spectral Zariski topology}

\begin{definition}[Affine Zariski covering family]
\label{def:affine-zariski-cover}
A small family of affine open immersions
\[
    \{\Spec(B_i)\to\Spec(A)\}_{i\in I}
\]
is an \textbf{affine spectral Zariski covering family} if the substitution
functors
\[
    \{-\otimes_AB_i:\Mod_A\to\Mod_{B_i}\}_{i\in I}
\]
are jointly conservative.
\end{definition}

The definition makes locality part of the coefficient system itself: an
affine family covers exactly when substitution along those context changes
jointly detects coefficients.

\begin{proposition}[The spectral Zariski pretopology]
\label{prop:spectral-zariski-pretopology}
Affine spectral Zariski covering families define a Grothendieck pretopology
on \(\Aff_{\Sp}\).
\end{proposition}

\begin{proof}
\leavevmode
\begin{enumerate}[label=\textup{(\arabic*)}]
    \item \textbf{Identity covers.}
    The identity substitution on \(\Mod_A\) is conservative, so the identity
    family covers.

    \item \textbf{Base change.}
    Let \(\{A\to B_i\}\) be covering and let \(A\to A'\) be any morphism.
    Put \(B_i':=B_i\otimes_AA'\). If \(M\in\Mod_{A'}\) satisfies
    \(M\otimes_{A'}B_i'\simeq0\) for every \(i\), then, after restriction to
    \(A\),
    \[
        M\otimes_AB_i
        \simeq
        M\otimes_{A'}B_i'
        \simeq0.
    \]
    Joint conservativity over \(A\) gives \(M\simeq0\). The base-changed
    family therefore covers. Its members remain affine open immersions by
    \Ref{prop:affine-open-permanence}.

    \item \textbf{Transitivity.}
    Suppose \(\{A\to B_i\}\) is covering and, for every \(i\),
    \(\{B_i\to C_{ij}\}_j\) is covering. If \(M\in\Mod_A\) restricts to
    zero over every \(C_{ij}\), then \(M\otimes_AB_i\simeq0\) for every
    \(i\) by joint conservativity over \(B_i\). Joint conservativity over
    \(A\) then gives \(M\simeq0\). Thus the composite family covers.
\end{enumerate}
\end{proof}

\begin{definition}[Spectral Zariski topology]
\label{def:spectral-zariski-topology}
The \textbf{spectral Zariski topology}
\[
    \tau_{\mathrm{Zar}}
\]
on \(\Aff_{\Sp}\) is the Grothendieck topology generated by the pretopology
of \Ref{prop:spectral-zariski-pretopology}.
\end{definition}

\begin{proposition}[Finite principal refinements]
\label{prop:finite-principal-refinement}
Every affine spectral Zariski covering family of \(\Spec(A)\) admits a finite
refinement by principal spectral opens
\[
    \{D_A^{\mathrm{sp}}(f_j)\to\Spec(A)\}_{j=1}^r
\]
with \(f_1,\ldots,f_r\in\pi_0(A)\) generating the unit ideal. The
refinement may be displayed as
\[
\begin{tikzcd}[column sep=huge]
    \displaystyle\coprod_{j=1}^r D_A^{\mathrm{sp}}(f_j)
        \arrow[r]
    &
    \displaystyle\coprod_{i\in I}\Spec(B_i)
        \arrow[r]
    &
    \Spec(A).
\end{tikzcd}
\]
Consequently, finite principal covers generate the spectral Zariski topology
on affines.
\end{proposition}

\begin{proof}
\leavevmode
\begin{enumerate}[label=\textup{(\arabic*)}]
    \item \textbf{Refine each member by principal opens.}
    Let \(\{\Spec(B_i)\to\Spec(A)\}_{i\in I}\) be a covering family. By
    \Ref{prop:principal-opens-basis}, each \(\Spec(B_i)\) is covered by
    principal spectral opens of \(\Spec(A)\) which factor through
    \(\Spec(B_i)\).

    \item \textbf{Form the combined principal family.}
    Let
    \(\{D_A^{\mathrm{sp}}(f_\alpha)\}_{\alpha\in J}\) be the combined
    family.

    \item \textbf{Assume vanishing on the combined family.}
    If \(M\in\Mod_A\) restricts to zero on every
    \(D_A^{\mathrm{sp}}(f_\alpha)\), then for each \(i\) the module
    \(M\otimes_AB_i\) restricts to zero on a jointly conservative principal
    cover of \(\Spec(B_i)\).

    \item \textbf{Descend vanishing to each original member.}
    Hence \(M\otimes_AB_i\simeq0\) for every
    \(i\).

    \item \textbf{Use the original cover.}
    Since the original family is covering, \(M\simeq0\).

    \item \textbf{Conclude conservativity of the combined family.}
    Thus the
    combined principal family is jointly conservative.

    \item \textbf{Define the coefficient-detecting module.}
    Let \(J_0\subseteq\pi_0(A)\) be the ideal generated by all
    \(f_\alpha\), and regard
    \[
        M:=H\bigl(\pi_0(A)/J_0\bigr)
    \]
    as an \(A\)-module through
    \[
        A\longrightarrow H\pi_0(A)
        \longrightarrow H\bigl(\pi_0(A)/J_0\bigr).
    \]

    \item \textbf{Apply the localization tensor formula.}
    For every \(\alpha\), the map
    \(A\to A[f_\alpha^{-1}]\) is flat by
    \Ref{prop:homotopy-groups-localization}. Applying the tensor formula of
    \Ref{thm:spectral-lazard} to the flat \(A\)-module
    \(A[f_\alpha^{-1}]\) gives
    \[
        \pi_n\bigl(M\otimes_AA[f_\alpha^{-1}]\bigr)
        \simeq
        \pi_n(M)[f_\alpha^{-1}]
        =0
    \]
    for all \(n\), because \(M\) is discrete and \(f_\alpha\) acts by zero
    on \(\pi_0(M)\).

    \item \textbf{Deduce vanishing after every principal localization.}
    Hence
    \[
        M\otimes_AA[f_\alpha^{-1}]\simeq0
    \]
    for every \(\alpha\).

    \item \textbf{Use joint conservativity to identify the unit ideal.}
    Joint conservativity of the combined principal
    family gives \(M\simeq0\), so
    \[
        \pi_0(A)/J_0=0
        \qquad\text{and therefore}\qquad
        J_0=(1).
    \]

    \item \textbf{Extract finitely many generators.}
    By the definition of the generated ideal, the relation \(1\in J_0\)
    involves only finitely many generators: there are
    \(f_1,\ldots,f_r\) among the \(f_\alpha\) and
    \(a_1,\ldots,a_r\in\pi_0(A)\) with
    \[
        1=a_1f_1+\cdots+a_rf_r.
    \]

    \item \textbf{Obtain the finite principal refinement.}
    Thus \((f_1,\ldots,f_r)=(1)\), and the corresponding finite principal
    family is the required refinement.
\end{enumerate}
\end{proof}

\begin{proposition}[Coefficient criterion for Zariski covers]
\label{prop:zariski-cover-joint-conservativity}
For a family of affine open immersions
\[
    \{\Spec(B_i)\to\Spec(A)\}_{i\in I},
\]
the following conditions are equivalent.
\begin{enumerate}[label=(\roman*)]
    \item The family is an affine spectral Zariski covering family.
    \item The pullback functors \(\{-\otimes_AB_i\}_{i\in I}\) are jointly
    conservative on \(\Mod_A\).
    \item The family admits a finite refinement by principal spectral opens
    \(D_A^{\mathrm{sp}}(f_j)\) with \((f_1,\ldots,f_r)=(1)\) in
    \(\pi_0(A)\).
\end{enumerate}
\end{proposition}

\begin{proof}
\leavevmode
\begin{enumerate}[label=\textup{(\arabic*)}]
    \item \textbf{Identify covering with conservative substitution.}
    Conditions \textup{(i)} and \textup{(ii)} are equivalent by
    \Ref{def:affine-zariski-cover}. The implication
    \textup{(i)}\(\Rightarrow\)\textup{(iii)} is
    \Ref{prop:finite-principal-refinement}.

    \item \textbf{Principal unit-ideal families are conservative.}
    Assume \textup{(iii)} and let \(M\in\Mod_A\) restrict to zero over every
    member of the original family. It then restricts to zero over every
    refining principal open, so by
    \Ref{prop:homotopy-groups-localization}
    \[
        \pi_n(M)[f_j^{-1}]=0
    \]
    for every \(n\) and \(j\). Since the \(f_j\) generate the unit ideal,
    an ordinary \(\pi_0(A)\)-module which vanishes after all these
    localizations is zero. Thus every \(\pi_n(M)\) vanishes, and
    \Ref{prop:postnikov-completeness-detection} gives \(M\simeq0\).
    Hence the original family is jointly conservative.
\end{enumerate}
\end{proof}

For a finite covering family, let
\[
    U:=\coprod_i\Spec(B_i)
\]
in \(\PSh(\Aff_{\Sp})\). Its \v{C}ech nerve packages all iterated overlaps;
the first stage is
\[
\begin{tikzcd}[column sep=huge, row sep=large]
    U\times_{\Spec(A)}U
        \arrow[r, shift left=1.2ex]
        \arrow[r, shift right=1.2ex]
    &
    U
        \arrow[r]
    &
    \Spec(A),
\end{tikzcd}
\]
and each affine overlap is represented by the corresponding tensor product
\(B_i\otimes_AB_j\). The next results show that these overlap diagrams detect covers and
reconstruct the objects being covered.

\[
\begin{tikzcd}[column sep=large, row sep=large]
    \Spec(B_i\otimes_A B_j)
        \arrow[r]
        \arrow[d]
    &
    \Spec(B_j)
        \arrow[d]
    \\
    \Spec(B_i)
        \arrow[r]
    &
    \Spec(A).
\end{tikzcd}
\]


\subsection{Descendability and Zariski descent}

\begin{definition}[Thick tensor ideal]
\label{def:spectral-thick-tensor-ideal}
Let $\mathcal C^\otimes$ be a stable symmetric monoidal
$\infty$-category. A full subcategory
\[
    \mathcal I
    \subseteq
    \mathcal C
\]
is a \textbf{thick tensor ideal} if it is closed under finite limits,
finite colimits, retracts, and tensoring with arbitrary objects of
$\mathcal C$.
\end{definition}

\begin{definition}[Descendable morphism]
\label{def:spectral-descendable-morphism}
A morphism
\[
    A
    \longrightarrow
    B
\]
of commutative ring spectra is \textbf{descendable} if $B$ generates
$\Mod_A$ as a thick tensor ideal.
\end{definition}

\begin{proposition}[Permanence of descendability]
\label{prop:spectral-descendability-permanence}
Descendable morphisms are stable under equivalence, base change, and
composition. If $A\to B\to C$ is a factorization and $A\to C$ is
descendable, then $A\to B$ is descendable.
\end{proposition}

\begin{proof}
\leavevmode
\begin{enumerate}[label=\textup{(\arabic*)}]
    \item \textbf{Handle base change.}
    Base change follows because extension of scalars is exact and strong
    symmetric monoidal.

    \item \textbf{Invoke the remaining permanence results.}
    The remaining assertions are the permanence results of
    \cite[Proposition~3.24]{MathewGaloisStableHomotopy}.
\end{enumerate}
\end{proof}

\begin{lemma}[Modules over finite product algebras]
\label{lem:spectral-modules-finite-products}
For a finite family of commutative ring spectra $\{B_i\}_{i\in I}$, there
is a canonical symmetric monoidal equivalence
\[
    \Mod_{\prod_iB_i}
    \simeq
    \prod_i\Mod_{B_i}.
\]
\end{lemma}

\begin{proof}
\leavevmode
\begin{enumerate}[label=\textup{(\arabic*)}]
    \item \textbf{Identify the underlying biproduct.}
    The finite product of commutative ring spectra has underlying finite
    biproduct of spectra.

    \item \textbf{Split modules by idempotents.}
    Its orthogonal idempotents split every module into
    its components.

    \item \textbf{Construct the inverse functors.}
    Extension of scalars along the projections and finite
    direct sum of restrictions of scalars define mutually inverse symmetric
    monoidal functors, exactly as for ordinary product rings.
\end{enumerate}
\end{proof}

\begin{lemma}[Finite products and base change]
\label{lem:finite-products-base-change}
Let \(A\to D\) be a morphism of commutative ring spectra and let
\(\{B_i\}_{i\in I}\) be a finite family of commutative \(A\)-algebras. The
canonical morphism
\[
    D\otimes_A\left(\prod_{i\in I}B_i\right)
    \longrightarrow
    \prod_{i\in I}(D\otimes_AB_i)
\]
is an equivalence. Consequently,
\[
    \left(\prod_iB_i\right)^{\otimes_A(n+1)}
    \simeq
    \prod_{(i_0,\ldots,i_n)}
    B_{i_0}\otimes_A\cdots\otimes_AB_{i_n}.
\]
\end{lemma}

\begin{proof}
\leavevmode
\begin{enumerate}[label=\textup{(\arabic*)}]
    \item \textbf{Underlying modules.}
    Finite products and finite coproducts agree in the stable category
    \(\Mod_A\), so
    \[
        \prod_iB_i\simeq\bigoplus_iB_i
    \]
    as \(A\)-modules.

    \item \textbf{Extension of scalars.}
    Tensoring with \(D\) preserves coproducts, giving
    \[
        D\otimes_A\prod_iB_i
        \simeq
        \bigoplus_i(D\otimes_AB_i)
        \simeq
        \prod_i(D\otimes_AB_i)
    \]
    on underlying \(D\)-modules. The canonical algebra morphism is therefore
    an equivalence. Iteration gives the second formula.
\end{enumerate}
\end{proof}

\begin{proposition}[Finite Zariski covers yield descendable product algebras]
\label{prop:finite-zariski-covers-descendable}
Let
\[
    \{\Spec(B_i)\to\Spec(A)\}_{i\in I}
\]
be a finite affine spectral Zariski covering family and put
\[
    B:=\prod_{i\in I}B_i.
\]
Then \(A\to B\) is descendable.
\end{proposition}

\begin{proof}
\leavevmode
\begin{enumerate}[label=\textup{(\arabic*)}]
    \item \textbf{Reduce to a finite principal cover.}
    By \Ref{prop:finite-principal-refinement}, choose a finite principal
    refinement
    \[
        \{D_A^{\mathrm{sp}}(f_j)\to\Spec(A)\}_{j=1}^r
    \]
    subordinate to the given cover. Put
    \[
        A_j:=A[f_j^{-1}],
        \qquad
        C:=\prod_{j=1}^rA_j.
    \]
    The refinement gives a factorization
    \[
    \begin{tikzcd}[column sep=large]
        A \arrow[r] & B \arrow[r] & C.
    \end{tikzcd}
    \]
    By \Ref{prop:spectral-descendability-permanence}, it is enough to prove
    that \(A\to C\) is descendable.

    \item \textbf{Fibres of the localization maps.}
    For each \(j\), let
    \[
        K_j:=\operatorname{fib}(A\longrightarrow A_j)
        \in\Mod_A.
    \]
    Localization is idempotent, so
    \[
        A_j\otimes_AA_j\simeq A_j.
    \]
    Exactness of extension of scalars therefore gives
    \[
        K_j\otimes_AA_j
        \simeq
        \operatorname{fib}
        \bigl(A_j\to A_j\otimes_AA_j\bigr)
        \simeq0.
    \]

    \item \textbf{The product of the fibres vanishes.}
    Put
    \[
        K:=K_1\otimes_A\cdots\otimes_AK_r.
    \]
    For every \(j\), one tensor factor becomes zero after extension to
    \(A_j\), hence
    \[
        K\otimes_AA_j\simeq0.
    \]
    The principal spectral opens \(D_A^{\mathrm{sp}}(f_j)\) cover \(\Spec(A)\), so the functors
    \(-\otimes_AA_j\) are jointly conservative by
    \Ref{prop:zariski-cover-joint-conservativity}. Consequently
    \[
        K\simeq0.
    \]

    \item \textbf{Place the localizations in the thick tensor ideal.}
    Let \(\mathcal I\) be the thick tensor ideal of \(\Mod_A\) generated by
    \(C\). Every \(A_j\) is a retract of the underlying finite product
    \(C\), so \(A_j\in\mathcal I\).

    \item \textbf{Define the inductive objects and target.}
    Define
    \[
        L_m:=K_1\otimes_A\cdots\otimes_AK_m,
        \qquad
        L_0:=A,
    \]
    and let \(u_m:L_m\to A\) be the tensor product of the maps
    \(K_j\to A\), followed by multiplication. We prove inductively that
    \[
        \operatorname{cofib}(u_m)\in\mathcal I.
    \]

    \item \textbf{Verify the base case.}
    For \(m=1\), the defining fibre sequence gives
    \[
        \operatorname{cofib}(K_1\to A)\simeq A_1\in\mathcal I.
    \]

    \item \textbf{Factor the inductive map.}
    For \(m>1\), factor
    \[
        L_m=L_{m-1}\otimes_AK_m
        \longrightarrow
        L_{m-1}
        \xrightarrow{u_{m-1}}
        A.
    \]

    \item \textbf{Identify the first cofiber.}
    The cofiber of the first map is
    \[
        L_{m-1}\otimes_AA_m,
    \]
    which lies in \(\mathcal I\) because \(A_m\in\mathcal I\) and
    \(\mathcal I\) is a tensor ideal.

    \item \textbf{Apply the composite cofiber sequence.}
    The cofiber sequence associated to a
    composite in a stable \(\infty\)-category then expresses
    \(\operatorname{cofib}(u_m)\) as an extension of
    \(L_{m-1}\otimes_AA_m\) by
    \(\operatorname{cofib}(u_{m-1})\). Hence
    \(\operatorname{cofib}(u_m)\in\mathcal I\).

    \item \textbf{Use vanishing of the total fibre product.}
    For \(m=r\), Step \textup{(3)} gives \(L_r=K\simeq0\). Therefore
    \[
        A
        \simeq
        \operatorname{cofib}(0\to A)
        \simeq
        \operatorname{cofib}(u_r)
        \in
        \mathcal I.
    \]

    \item \textbf{Generate the tensor unit.}
    Thus \(C\) generates the tensor unit and \(A\to C\) is descendable.

    \item \textbf{Return to the original cover.}
    The factorization \(A\to B\to C\) and the intermediate-algebra property
    of descendability imply that \(A\to B\) is descendable; see
    \cite[Proposition~3.24]{MathewGaloisStableHomotopy}.
\end{enumerate}
\end{proof}

\begin{theorem}[Amitsur convergence]
\label{thm:spectral-amitsur-convergence}
If $A\to B$ is descendable, then the augmented Amitsur object
\[
    A
    \longrightarrow
    B
    \triplearrows
    B\otimes_A B
    \cdots
\]
is a limit diagram in $\CAlg(\Sp)$. In particular,
\[
    A
    \simeq
    \operatorname*{Tot}_{[n]\in\Delta}
    B^{\otimes_A(n+1)}.
\]
\end{theorem}

\begin{proof}
\leavevmode
\begin{enumerate}[label=\textup{(\arabic*)}]
    \item \textbf{Regard the cover algebra in the module category.}
    Regard $B$ as a commutative algebra object of $\Mod_A$.

    \item \textbf{Translate descendability.}
    Descendability
    means that $B$ generates the tensor unit.

    \item \textbf{Apply the descendability theorem.}
    The descendability theorem
    identifies the augmented cobar object as a limit diagram in $\Mod_A$; see
    \cite[Proposition~3.20]{MathewGaloisStableHomotopy}.

    \item \textbf{Create the same limit in commutative algebras.}
    The forgetful functor
    from commutative $A$-algebras to $A$-modules creates limits, so the same
    object is a limit in commutative $A$-algebras.

    \item \textbf{Pass to commutative ring spectra.}
    It is therefore a limit in
    $\CAlg(\Sp)$.
\end{enumerate}
\end{proof}

\begin{theorem}[Affine module descent]
\label{thm:spectral-affine-module-descent}
If $A\to B$ is descendable, then extension of scalars induces an
equivalence in $\CAlg(\Cat_\infty)$
\[
    \Mod_A
    \xrightarrow{\simeq}
    \operatorname*{Tot}_{[n]\in\Delta}
    \Mod_{B^{\otimes_A(n+1)}}.
\]
\end{theorem}

\begin{proof}
\leavevmode
\begin{enumerate}[label=\textup{(\arabic*)}]
    \item \textbf{Apply comonadicity.}
    The extension--restriction adjunction is comonadic for a descendable map.

    \item \textbf{Identify the cobar resolution.}
    Its cobar resolution is the displayed Amitsur diagram; see
    \cite[Proposition~3.22]{MathewGaloisStableHomotopy}.

    \item \textbf{Retain the symmetric monoidal structure.}
    Every transition
    functor is strong symmetric monoidal, so the comparison is an equivalence
    of symmetric monoidal $\infty$-categories.
\end{enumerate}
\end{proof}

\begin{theorem}[Affine algebra descent]
\label{thm:spectral-affine-algebra-descent}
If $A\to B$ is descendable, then base change induces an equivalence
\[
    \CAlg(\Sp)_{A/}
    \xrightarrow{\simeq}
    \operatorname*{Tot}_{[n]\in\Delta}
    \CAlg(\Sp)_{B^{\otimes_A(n+1)}/}.
\]
\end{theorem}

\[
\begin{tikzcd}[column sep=large, row sep=large]
    \Mod_A
        \arrow[rr, "\simeq"]
        \arrow[d, "\CAlg"']
    &&
    \displaystyle
    \operatorname*{Tot}_{[n]\in\Delta}
    \Mod_{B^{\otimes_A(n+1)}}
        \arrow[d, "\CAlg"]
    \\
    \CAlg(\Mod_A)
        \arrow[rr, "\simeq"']
    &&
    \displaystyle
    \CAlg\!\left(
        \operatorname*{Tot}_{[n]\in\Delta}
        \Mod_{B^{\otimes_A(n+1)}}
    \right)
        \arrow[r, "\simeq"']
    &
    \displaystyle
    \operatorname*{Tot}_{[n]\in\Delta}
    \CAlg\!\bigl(\Mod_{B^{\otimes_A(n+1)}}\bigr).
\end{tikzcd}
\]


\begin{proof}
\leavevmode
\begin{enumerate}[label=\textup{(\arabic*)}]
    \item \textbf{Apply affine module descent.}
    By \Ref{thm:spectral-affine-module-descent}, extension of scalars gives a
    symmetric monoidal equivalence
    \[
        \Mod_A
        \xrightarrow{\simeq}
        \operatorname*{Tot}_{[n]\in\Delta}
        \Mod_{B^{\otimes_A(n+1)}}.
    \]

    \item \textbf{Use limit preservation of commutative algebra objects.}
    The functor assigning to a symmetric monoidal $\infty$-category its
    $\infty$-category of commutative algebra objects preserves limits.

    \item \textbf{Apply commutative algebra objects.}
    Applying it to the displayed equivalence gives
    \[
        \CAlg(\Mod_A)
        \xrightarrow{\simeq}
        \operatorname*{Tot}_{[n]\in\Delta}
        \CAlg
        \bigl(
            \Mod_{B^{\otimes_A(n+1)}}
        \bigr).
    \]

    \item \textbf{Identify relative commutative algebras.}
    Using
    \Ref{prop:relative-commutative-spectral-algebras} in every cosimplicial
    degree identifies this equivalence with
    \[
        \CAlg(\Sp)_{A/}
        \xrightarrow{\simeq}
        \operatorname*{Tot}_{[n]\in\Delta}
        \CAlg(\Sp)_{B^{\otimes_A(n+1)}/}.
    \]
\end{enumerate}
\end{proof}

\begin{corollary}[Connective algebra descent for finite Zariski covers]
\label{cor:connective-zariski-algebra-descent}
Let
\[
    \{\Spec(B_i)\to\Spec(A)\}_{i\in I}
\]
be a finite affine spectral Zariski covering family and put
\[
    B
    :=
    \prod_{i\in I}B_i.
\]
Then the equivalence of
\Ref{thm:spectral-affine-algebra-descent} restricts to an equivalence
\[
    \CAlg(\Sp)^{\mathrm{cn}}_{A/}
    \xrightarrow{\simeq}
    \operatorname*{Tot}_{[n]\in\Delta}
    \CAlg(\Sp)^{\mathrm{cn}}_
    {B^{\otimes_A(n+1)}/}.
\]
\end{corollary}

\begin{proof}
\leavevmode
\begin{enumerate}[label=\textup{(\arabic*)}]
    \item \textbf{Flat descendable cover algebra.}
    The morphism \(A\to B\) is descendable by
    \Ref{prop:finite-zariski-covers-descendable}. Each \(B_i\) is flat over
    \(A\), and the finite product \(B\) is the finite direct sum of the
    \(B_i\) as an \(A\)-module. Hence \(A\to B\) is flat. Every Amitsur term
    \(B^{\otimes_A(n+1)}\) is therefore connective, and base change along
    \(A\to B\) is \(t\)-exact and conservative.

    \item \textbf{Descend an object of the connective totalization.}
    Let an object of the right-hand totalization be given whose component in
    every cosimplicial degree is connective. By
    \Ref{thm:spectral-affine-algebra-descent}, it descends to a commutative
    \(A\)-algebra \(C\). Its base change \(C\otimes_AB\) is connective, so
    \[
        \tau_{\leq-1}^{A}(C)\otimes_AB
        \simeq
        \tau_{\leq-1}^{B}(C\otimes_AB)
        \simeq0,
    \]
    where the algebras are regarded as their underlying modules.
    Conservativity of \(-\otimes_AB\) gives
    \[
        \tau_{\leq-1}^{A}(C)\simeq0,
    \]
    hence the underlying \(A\)-module of \(C\), and therefore \(C\), is
    connective.

    \item \textbf{Restrict the descent equivalence.}
    Essential surjectivity on connective objects follows from the preceding
    step, while full faithfulness is inherited from the unrestricted algebra
    descent equivalence. This gives the asserted totalization equivalence.
\end{enumerate}
\end{proof}

\begin{theorem}[Subcanonicity of the spectral Zariski topology]
\label{thm:spectral-zariski-subcanonical}
Every representable spectral prestack is a sheaf for
\(\tau_{\mathrm{Zar}}\).
\end{theorem}

\begin{proof}
\leavevmode
\begin{enumerate}[label=\textup{(\arabic*)}]
    \item \textbf{Reduction to finite principal covers.}
    By \Ref{prop:finite-principal-refinement}, finite principal covers generate
    the affine spectral Zariski topology. It therefore suffices to prove
    \v{C}ech descent for a finite affine cover
    \(\{\Spec(B_i)\to\Spec(A)\}\).

    \item \textbf{Passage to one descendable algebra.}
    Put \(B:=\prod_iB_i\). By
    \Ref{prop:finite-zariski-covers-descendable}, \(A\to B\) is descendable,
    so \Ref{thm:spectral-amitsur-convergence} gives
    \[
        A
        \simeq
        \operatorname*{Tot}_{[n]\in\Delta}
        B^{\otimes_A(n+1)}.
    \]

    \item \textbf{Mapping into the Amitsur totalization.}
    For every connective commutative ring spectrum \(C\), mapping out of
    \(C\) converts this limit into
    \[
        \Map_{\CAlg(\Sp)}(C,A)
        \simeq
        \operatorname*{Tot}_{[n]\in\Delta}
        \Map_{\CAlg(\Sp)}
        \bigl(C,B^{\otimes_A(n+1)}\bigr).
    \]
    By \Ref{lem:finite-products-base-change}, the Amitsur terms decompose into
    the iterated intersections of the original covering family. This is the
    \v{C}ech descent condition for the representable \(h_C\).
\end{enumerate}
\end{proof}

\begin{theorem}[Zariski descent for affine quasi-coherent modules]
\label{thm:affine-qcoh-zariski-descent}
Let
\[
    U:=\coprod_i\Spec(B_i)
    \longrightarrow
    \Spec(A)
\]
be a finite affine spectral Zariski cover, and let $U_\bullet$ be its
\v{C}ech nerve. Then the comparison functor is an equivalence
\[
    \Mod_A
    \xrightarrow{\simeq}
    \operatorname*{Tot}_{[n]\in\Delta}
    \QCoh(U_n).
\]
\end{theorem}

\begin{proof}
\leavevmode
\begin{enumerate}[label=\textup{(\arabic*)}]
    \item \textbf{Identify the prestack \v{C}ech terms.}
    Put \(B:=\prod_iB_i\). By
    \Ref{lem:finite-products-base-change}, the degree-\(n\) Amitsur algebra
    decomposes as
    \[
        B^{\otimes_A(n+1)}
        \simeq
        \prod_{(i_0,\ldots,i_n)}
        B_{i_0}\otimes_A\cdots\otimes_AB_{i_n}.
    \]
    The corresponding affine components are exactly the iterated fibre
    products occurring in the prestack \v{C}ech nerve of
    \(U\to\Spec(A)\).

    \item \textbf{Compute quasi-coherent modules on the coproducts.}
    Since \(\QCoh\) carries prestack colimits to limits by
    \Ref{cor:qcoh-colimits-to-limits}, and modules over a finite product split
    by \Ref{lem:spectral-modules-finite-products},
    \[
        \QCoh(U_n)
        \simeq
        \prod_{(i_0,\ldots,i_n)}
        \Mod_{B_{i_0}\otimes_A\cdots\otimes_AB_{i_n}}
        \simeq
        \Mod_{B^{\otimes_A(n+1)}}.
    \]

    \item \textbf{Apply descendable module descent.}
    The map \(A\to B\) is descendable by
    \Ref{prop:finite-zariski-covers-descendable}. Therefore
    \Ref{thm:spectral-affine-module-descent} gives
    \[
        \Mod_A
        \xrightarrow{\simeq}
        \operatorname*{Tot}_{[n]\in\Delta}
        \Mod_{B^{\otimes_A(n+1)}}.
    \]
    Under the identifications of the preceding steps this is the asserted
    \v{C}ech descent equivalence.
\end{enumerate}
\end{proof}

\begin{proposition}[Affine open covers and effective epimorphisms]
\label{prop:affine-open-cover-effective-epi}
Let
\[
    \{U_i=\Spec(B_i)\longrightarrow X=\Spec(A)\}_{i\in I}
\]
be a small family of affine open immersions. The following conditions are
equivalent.
\begin{enumerate}[label=(\roman*)]
    \item The family is an affine spectral Zariski covering family.
    \item The induced morphism of spectral Zariski stacks
    \[
        \coprod_{i\in I}U_i
        \longrightarrow
        X
    \]
    is an effective epimorphism.
\end{enumerate}
For every affine base change \(T\to X\), the covering morphism is carried to
\[
\begin{tikzcd}[column sep=large, row sep=large]
    \displaystyle\coprod_i(T\times_XU_i)
        \arrow[r]
        \arrow[d, two heads]
    &
    \displaystyle\coprod_iU_i
        \arrow[d, two heads]
    \\
    T
        \arrow[r]
    &
    X,
\end{tikzcd}
\]
and the left vertical morphism is again effective epimorphic.
\end{proposition}

\begin{proof}
\leavevmode
\begin{enumerate}[label=\textup{(\arabic*)}]
    \item \textbf{A covering family is locally surjective.}
    Assume \textup{(i)}. For every affine test \(T\to X\), base change of
    the family is again a Zariski covering family by
    \Ref{prop:spectral-zariski-pretopology}. Hence the section of \(X\) over
    \(T\) determined by \(T\to X\) lifts locally to one of the pullbacks
    \(T\times_XU_i\). The local-surjectivity characterization of effective
    epimorphisms from
    \Ref{conv:spectral-effective-epimorphisms} therefore makes the coproduct
    map an effective epimorphism.

    \item \textbf{An effective affine-open family is conservative.}
    Assume \textup{(ii)}. Apply local surjectivity to the identity section
    \(\id_X\in X(X)\). The resulting covering sieve contains, because the
    topology is generated by the pretopology of
    \Ref{prop:spectral-zariski-pretopology}, an affine spectral Zariski
    covering family
    \[
        \{V_j\longrightarrow X\}_{j\in J}
    \]
    such that each \(V_j\to X\) factors through some \(U_i\). If
    \(M\in\Mod_A\) restricts to zero on every \(U_i\), it restricts to zero
    on every \(V_j\). Joint conservativity of the latter covering family
    gives \(M\simeq0\). Thus the original family is jointly conservative
    and is a covering family by \Ref{def:affine-zariski-cover}.
\end{enumerate}
\end{proof}

The affine locality problem is now solved in the form needed for gluing.
Principal spectral opens generate the topology, joint conservativity detects
covers, and the same covering data reconstruct representables, modules, and
connective commutative algebras. The next section uses these reconstruction
results to pass from affine contexts to global spectral schemes.

\section{Global spectral contexts}
\label{sec:spectral-stacks-schemes}

The affine site is now subcanonical, and its covering families reconstruct
both geometric and coefficient data. We can therefore pass from local
contexts to global ones by gluing affine opens. This section defines spectral
schemes as Zariski sheaves with effective affine open atlases, establishes the
closure properties needed later, and finally identifies this functor-of-points
model with the standard structured spectral-scheme model.

\begin{definition}[Spectral Zariski stack]
\label{def:spectral-zariski-stack}
A \textbf{spectral Zariski stack} is a sheaf of spaces on
$\Aff_{\Sp}$ for the spectral Zariski topology. The
$\infty$-category of spectral Zariski stacks is denoted
\[
    \operatorname{Stk}^{\mathrm{sp}}_{\mathrm{Zar}}
    :=
    \Sh_{\tau_{\mathrm{Zar}}}
    (\Aff_{\Sp}).
\]
\end{definition}

\begin{construction}[Spectral Zariski sheafification]
\label{con:spectral-zariski-sheafification}
The inclusion
\[
    \operatorname{Stk}^{\mathrm{sp}}_{\mathrm{Zar}}
    \hookrightarrow
    \PSh(\Aff_{\Sp})
\]
is an accessible left-exact localization. Its left adjoint is denoted
\[
    a^{\mathrm{sp}}_{\mathrm{Zar}}
    \colon
    \PSh(\Aff_{\Sp})
    \longrightarrow
    \operatorname{Stk}^{\mathrm{sp}}_{\mathrm{Zar}}
\]
and is called \textbf{spectral Zariski sheafification}. It preserves all
small colimits and finite limits. See \cite{LurieSAG}.
\end{construction}

Subcanonicity allows affine contexts to be regarded without ambiguity as
objects of the Zariski topos:
\[
\begin{tikzcd}[column sep=large]
    \Aff_{\Sp}
        \arrow[r, hook, "h"]
    &
    \operatorname{Stk}^{\mathrm{sp}}_{\mathrm{Zar}}
        \arrow[r, hook]
    &
    \PSh(\Aff_{\Sp}).
\end{tikzcd}
\]

\subsection{Representable open immersions}

\begin{definition}[Open immersion of spectral stacks]
\label{def:open-immersion-spectral-stacks}
Let
\[
    j
    \colon
    U
    \longrightarrow
    X
\]
be a morphism of spectral Zariski stacks.

Suppose first that $X$ is an affine spectral scheme. Then $j$ is an
\textbf{open immersion} if:
\begin{enumerate}[label=(\roman*)]
    \item $j$ is a monomorphism;

    \item there exists a small family of affine spectral schemes
    $\{U_\alpha\}_{\alpha\in I}$ and morphisms
    \[
        U_\alpha
        \longrightarrow
        U
    \]
    such that
    \[
        \coprod_{\alpha\in I}U_\alpha
        \longrightarrow
        U
    \]
    is an effective epimorphism and every composite
    \[
        U_\alpha
        \longrightarrow
        U
        \longrightarrow
        X
    \]
    is an affine open immersion in the sense of
    \Ref{def:affine-open-immersion-spectral}.
\end{enumerate}

For a general spectral Zariski stack $X$, the morphism $j$ is an open
immersion if, for every affine spectral scheme $T$ and every morphism
$T\to X$, the pullback
\[
    U\times_XT
    \longrightarrow
    T
\]
is an open immersion in the preceding sense.
\end{definition}

\begin{proposition}[Permanence of open immersions]
\label{prop:global-open-immersion-permanence}
Open immersions of spectral Zariski stacks are stable under equivalence,
composition, and base change.
\end{proposition}

\begin{proof}
\leavevmode
\begin{enumerate}[label=\textup{(\arabic*)}]
    \item \textbf{Handle equivalences.}
    Equivalences preserve open immersions.

    \item \textbf{Handle base change.}
    Stability under base change is
    part of the representable definition.

    \item \textbf{Reduce composition to an affine target.}
    For composable open immersions
    \[
        U\xrightarrow{j}V\xrightarrow{k}X,
    \]
    base change along an affine point \(T\to X\) reduces the assertion to
    the case \(X=T\) affine. Choose an affine open atlas
    \(\coprod_\alpha V_\alpha\to V\) whose composites to \(T\) are affine
    open immersions. Pull \(U\) back to each \(V_\alpha\) and choose an
    affine open atlas \(\coprod_\beta U_{\alpha\beta}\to
    U\times_VV_\alpha\).

    \item \textbf{Compose the affine charts.}
    Every
    \(U_{\alpha\beta}\to V_\alpha\to T\) is an affine open immersion by
    \Ref{prop:affine-open-permanence}. The coproduct
    \[
        \coprod_{\alpha,\beta}U_{\alpha\beta}\longrightarrow U
    \]
    is an effective epimorphism, and \(U\to T\) is a monomorphism. Hence
    \(U\to T\) satisfies the definition of an open immersion. This proves
    composition stability.
\end{enumerate}
\end{proof}

\subsection{Spectral schemes}

\begin{definition}[Spectral scheme]
\label{def:spectral-scheme}
A \textbf{spectral scheme} is a spectral Zariski stack $X$ for which there
exists a small family of affine spectral schemes $\{U_i\}_{i\in I}$ and a
morphism
\[
    \coprod_{i\in I}U_i
    \longrightarrow
    X
\]
such that:
\begin{enumerate}[label=(\roman*)]
    \item the displayed morphism is an effective epimorphism;

    \item every component $U_i\to X$ is an open immersion.
\end{enumerate}
Such a family is an \textbf{affine spectral Zariski atlas} of $X$.
\end{definition}

The full subcategory of
\[
    \operatorname{Stk}^{\mathrm{sp}}_{\mathrm{Zar}}
\]
spanned by spectral schemes is denoted
\[
    \operatorname{Sch}^{\mathrm{sp}}.
\]

\begin{proposition}[Affine spectral schemes are spectral schemes]
\label{prop:affines-are-spectral-schemes}
Every affine spectral scheme is a spectral scheme.
\end{proposition}

\begin{proof}
\leavevmode
\begin{enumerate}[label=\textup{(\arabic*)}]
    \item \textbf{Use the identity atlas.}
    Use the singleton atlas given by the identity map.
\end{enumerate}
\end{proof}

\begin{proposition}[Coproducts of spectral schemes]
\label{prop:spectral-schemes-coproducts}
Small coproducts of spectral schemes exist and are spectral schemes.
\end{proposition}

\begin{proof}
\leavevmode
\begin{enumerate}[label=\textup{(\arabic*)}]
    \item \textbf{Form the disjoint union in the Zariski topos.}
    Coproducts exist in the ambient \(\infty\)-topos of spectral Zariski
    stacks. The inclusion of a coproduct summand is an open immersion. This is
    the usual disjoint-union construction for spectral schemes: after affine
    base change, the summand pullbacks are open-and-closed affine pieces,
    equivalently the pieces determined by idempotents in degree zero; see
    \cite{LurieSAG}.

    \item \textbf{Coproduct the affine atlases.}
    Choose affine open atlases \(\coprod_\alpha U_{i\alpha}\to X_i\).
    The composite
    \[
        \coprod_{i,\alpha}U_{i\alpha}
        \longrightarrow
        \coprod_iX_i
    \]
    is an effective epimorphism, and each component is an open immersion by
    composition with the open summand inclusion. Hence the ambient coproduct
    is a spectral scheme and is the coproduct in
    \(\operatorname{Sch}^{\mathrm{sp}}\).
\end{enumerate}
\end{proof}

\begin{lemma}[Open substacks of spectral schemes]
\label{lem:open-substack-spectral-scheme}
Let \(X\) be a spectral scheme and let
\[
    U\longrightarrow X
\]
be an open immersion of spectral Zariski stacks. Then \(U\) is a spectral
scheme.
\end{lemma}

\begin{proof}
\leavevmode
\begin{enumerate}[label=\textup{(\arabic*)}]
    \item \textbf{Pull back an affine atlas.}
    Choose an affine spectral Zariski atlas
    \[
        \coprod_iX_i\longrightarrow X.
    \]
    Put \(U_i:=U\times_XX_i\). Each \(U_i\to X_i\) is an open immersion by
    base change.

    \item \textbf{Refine over each affine chart.}
    By the definition of an open immersion with affine target, every \(U_i\)
    admits an effective epimorphism
    \[
        \coprod_\alpha U_{i\alpha}\longrightarrow U_i
    \]
    by affine spectral schemes whose composites
    \(U_{i\alpha}\to X_i\) are affine open immersions.

    \item \textbf{Assemble the atlas.}
    Effective epimorphisms are stable under base change and composition in an
    \(\infty\)-topos. Hence
    \[
        \coprod_{i,\alpha}U_{i\alpha}\longrightarrow U
    \]
    is an effective epimorphism. Every component is an open immersion by
    composition and base change. Thus it is an affine spectral Zariski atlas
    of \(U\).
\end{enumerate}
\end{proof}

\begin{proposition}[Small unions of open substacks]
\label{prop:unions-open-substacks}
Let \(X\) be a spectral Zariski stack and let
\[
    \{j_\alpha:U_\alpha\to X\}_{\alpha\in I}
\]
be a small family of open immersions. Factor the induced morphism in the
spectral Zariski \(\infty\)-topos as an effective epimorphism followed by a
monomorphism:
\[
\begin{tikzcd}[column sep=huge]
    \displaystyle\coprod_{\alpha\in I}U_\alpha
        \arrow[r, two heads]
    &
    U
        \arrow[r, hook, "j"]
    &
    X.
\end{tikzcd}
\]
Then \(j:U\to X\) is an open immersion. This construction is stable under
arbitrary base change. If \(X\) is a spectral scheme, then \(U\) is a
spectral scheme.
\end{proposition}

\begin{proof}
\leavevmode
\begin{enumerate}[label=\textup{(\arabic*)}]
    \item \textbf{Base change the image factorization.}
    Image factorizations in an \(\infty\)-topos are stable under pullback.
    For every affine point \(T\to X\), the pullback \(U_T:=U\times_XT\) is
    therefore the image in
    \[
    \begin{tikzcd}[column sep=huge]
        \displaystyle\coprod_{\alpha}
        (U_\alpha\times_XT)
            \arrow[r, two heads]
        &
        U_T
            \arrow[r, hook]
        &
        T.
    \end{tikzcd}
    \]
    Each \(U_\alpha\times_XT\to T\) is open by base-change stability of open
    immersions.

    \item \textbf{Produce an affine open atlas of the image.}
    For every \(\alpha\), choose an affine open atlas
    \[
        \coprod_\beta V_{\alpha\beta}
        \longrightarrow
        U_\alpha\times_XT.
    \]
    Every composite \(V_{\alpha\beta}\to T\) is an affine open immersion.
    Since effective epimorphisms are stable under coproduct and composition,
    the induced map
    \[
        \coprod_{\alpha,\beta}V_{\alpha\beta}
        \longrightarrow
        U_T
    \]
    is an effective epimorphism. The morphism \(U_T\to T\) is a monomorphism
    by construction. Hence \(U_T\to T\) is an open immersion in the sense of
    \Ref{def:open-immersion-spectral-stacks}.

    \item \textbf{Globalize the affine-point argument.}
    Since the preceding argument applies to every affine point of \(X\), the
    morphism \(U\to X\) is an open immersion.

    \item \textbf{Record base-change stability.}
    The construction commutes with
    arbitrary base change because it was defined by the universal image
    factorization.

    \item \textbf{Recognize the open as a spectral scheme.}
    If \(X\) is a spectral scheme, then \(U\) is a spectral
    scheme by \Ref{lem:open-substack-spectral-scheme}.
\end{enumerate}
\end{proof}

\begin{theorem}[Fibre products of spectral schemes]
\label{thm:spectral-schemes-fibre-products}
Let
\[
    X\longrightarrow Z\longleftarrow Y
\]
be morphisms of spectral schemes. Then the fibre product
\[
    X\times_ZY
\]
is a spectral scheme.
\end{theorem}

\begin{proof}
\leavevmode
\begin{enumerate}[label=\textup{(\arabic*)}]
    \item \textbf{Choose one affine atlas of the base.}
    Choose an affine spectral Zariski atlas
    \[
        \coprod_kW_k\longrightarrow Z.
    \]
    Put
    \[
        X_k:=X\times_ZW_k,
        \qquad
        Y_k:=Y\times_ZW_k.
    \]
    Since \(W_k\to Z\) is open, the projections \(X_k\to X\) and
    \(Y_k\to Y\) are open. Hence \(X_k\) and \(Y_k\) are spectral schemes
    by \Ref{lem:open-substack-spectral-scheme}.

    \item \textbf{Choose affine atlases over the same base chart.}
    For every \(k\), choose affine open atlases
    \[
        \coprod_\alpha U_{k\alpha}\longrightarrow X_k,
        \qquad
        \coprod_\beta V_{k\beta}\longrightarrow Y_k.
    \]
    By construction, every \(U_{k\alpha}\) and \(V_{k\beta}\) maps to the
    same affine spectral scheme \(W_k\).

    \item \textbf{Form affine fibre products.}
    Put
    \[
        P_{k\alpha\beta}
        :=
        U_{k\alpha}\times_{W_k}V_{k\beta}.
    \]
    Each \(P_{k\alpha\beta}\) is affine by
    \Ref{prop:affine-spectral-fibre-products}. The canonical map
    \[
        P_{k\alpha\beta}
        \longrightarrow
        X\times_ZY
    \]
    is an open immersion, since it is obtained by composing base changes of
    the open immersions \(U_{k\alpha}\to X_k\),
    \(V_{k\beta}\to Y_k\), and \(W_k\to Z\).

    \item \textbf{Verify that the affine charts cover.}
    The pullbacks
    \[
        (X\times_ZY)\times_ZW_k
        \simeq
        X_k\times_{W_k}Y_k
    \]
    cover \(X\times_ZY\). For fixed \(k\), base changing the atlas of
    \(X_k\) to \(X_k\times_{W_k}Y_k\), and then base changing the atlas of
    \(Y_k\), shows that
    \[
        \coprod_{\alpha,\beta}P_{k\alpha\beta}
        \longrightarrow
        X_k\times_{W_k}Y_k
    \]
    is an effective epimorphism. Universality and composition of effective
    epimorphisms therefore give an effective epimorphism
    \[
        \coprod_{k,\alpha,\beta}P_{k\alpha\beta}
        \longrightarrow
        X\times_ZY.
    \]
    Thus these affine open immersions form an affine spectral Zariski atlas.
\end{enumerate}
\end{proof}

\begin{corollary}[Finite limits of spectral schemes]
\label{cor:spectral-schemes-finite-limits}
The $\infty$-category of spectral schemes admits finite limits, computed
in spectral Zariski stacks.
\end{corollary}

\begin{proof}
\leavevmode
\begin{enumerate}[label=\textup{(\arabic*)}]
    \item \textbf{Identify the terminal object.}
    The terminal object is $\Spec(\mathbb S)$.

    \item \textbf{Construct binary fibre products.}
    Binary fibre products exist
    by \Ref{thm:spectral-schemes-fibre-products}.
\end{enumerate}
\end{proof}

\subsection{Comparison with structured spectral schemes}

\begin{theorem}[Recognition of the structured spectral model]
\label{thm:spectral-scheme-model-comparison}
Let
\[
    \operatorname{Sch}^{\mathrm{sp}}_{\mathrm{str}}
\]
denote the \(\infty\)-category of connective \(0\)-localic spectral schemes
in the structured-space sense of spectral algebraic geometry. The standard
functor of points
\[
    \Phi(\mathfrak X)(\Spec A)
    :=
    \Map(\Spec A,\mathfrak X)
\]
induces an equivalence
\[
\begin{tikzcd}[column sep=huge]
    \operatorname{Sch}^{\mathrm{sp}}_{\mathrm{str}}
        \arrow[r, "\Phi", "\simeq"']
    &
    \operatorname{Sch}^{\mathrm{sp}}
        \arrow[r, hook]
    &
    \operatorname{Stk}^{\mathrm{sp}}_{\mathrm{Zar}}.
\end{tikzcd}
\]
Under this equivalence:
\begin{enumerate}[label=(\roman*)]
    \item \(\Phi(\Spec A)\simeq h_A\) for every connective commutative ring
    spectrum \(A\);

    \item affine morphisms, open immersions, and finite fibre products agree
    in the two models.
\end{enumerate}
\end{theorem}

\begin{proof}
\leavevmode
\begin{enumerate}[label=\textup{(\arabic*)}]
    \item \textbf{Verify that the native affine geometry is the standard
    connective spectral Zariski geometry.}
    Affines are the opposites of connective commutative ring spectra by
    \Ref{def:affine-spectral-scheme}. Principal localizations are affine open
    immersions by \Ref{prop:principal-spectral-open}, principal opens form a
    basis by \Ref{prop:principal-opens-basis}, and
    \Ref{thm:spectral-zariski-subcanonical} identifies representables as
    Zariski sheaves. Thus the affine objects, local morphisms, and covering
    families used in the native construction are the standard connective
    spectral Zariski ones.

    \item \textbf{Invoke the structured-space recognition theorem.}
    The comparison theorem between connective \(0\)-localic structured
    spectral schemes and their functors of points identifies precisely the
    Zariski stacks admitting affine open atlases with the essential image of
    the structured model; see
    \cite{LurieSAG,KhanBraveNewMotivicI}. It identifies an affine structured
    spectral scheme with \(h_A\), and representable properties such as
    affineness and open immersion are preserved under the comparison. Since
    functors of points preserve finite limits objectwise, finite fibre products
    agree as well. This proves all assertions.
\end{enumerate}
\end{proof}

\begin{remark}[The structured model is used only for recognition]
\label{rem:no-hidden-admissibility-spectral-zariski}
The comparison theorem identifies the spectral schemes constructed above with
the structured-space model of \cite{LurieSAG} while leaving the data carried by
an object of \(\operatorname{Sch}^{\mathrm{sp}}\) unchanged. Terminology such
as \textbf{admissible} belongs to the external structured presentation. The
open morphisms and covering families used here remain the spectral localizations
and conservative families constructed before the comparison. The structured
model serves to recognize this geometry as the standard connective spectral
one.
\end{remark}

\subsection{Underlying topological support and qcqs geometry}

\begin{definition}[Underlying topological space]
\label{def:underlying-space-spectral-scheme}
Let \(X\) be a spectral scheme and choose an affine spectral Zariski atlas
\[
    \coprod_i\Spec(A_i)\longrightarrow X.
\]
The spaces \(\lvert\Spec\pi_0(A_i)\rvert\) glue along the open subsets
represented by the affine spectral intersections and their principal-open
refinements. The resulting topological space is independent of the chosen
atlas by common refinement and is denoted
\[
    |X|.
\]
For \(X=\Spec(A)\) this agrees with
\Ref{def:underlying-topological-space-affine-spectral}. See
\cite{LurieSAG} for the equivalent structured-space construction.
\end{definition}

\begin{definition}[Quasi-compact spectral scheme]
\label{def:quasi-compact-spectral-scheme}
A spectral scheme $X$ is \textbf{quasi-compact} if it admits a finite
affine open cover. A morphism $f:X\to Y$ is quasi-compact if the pullback
over every affine open of $Y$ is quasi-compact.
\end{definition}

\begin{definition}[Quasi-separated morphisms and spectral schemes]
\label{def:quasi-separated-spectral-scheme}
A morphism
\[
    f
    \colon
    X
    \longrightarrow
    Y
\]
of spectral schemes is \textbf{quasi-separated} if its diagonal
\[
    \Delta_f
    \colon
    X
    \longrightarrow
    X\times_YX
\]
is quasi-compact.

A spectral scheme $X$ is \textbf{quasi-separated} if its structural
morphism
\[
    X
    \longrightarrow
    \Spec(\mathbb S)
\]
is quasi-separated. Equivalently, the absolute diagonal
\[
    X
    \longrightarrow
    X\times X
\]
is quasi-compact.

A spectral scheme is \textbf{qcqs} if it is quasi-compact and
quasi-separated.
\end{definition}

The global geometric base is now available: spectral schemes are obtained by
Zariski gluing from the affine contexts, remain closed under the finite-limit
and open-subobject operations needed below, and agree with the standard
connective structured model. The coefficient system must now be restricted to
this geometric base and shown to admit the corresponding global structures.

\section{The global quasi-coherent coefficient system}
\label{sec:qcoh-spectral-schemes-relative-spec}

For a global spectral context, the coefficient fibre should be recoverable
from the affine fibres on any atlas. The affine descent theorem provides that
reconstruction and thereby turns the right-Kan-extended prestack assignment
into the coefficient system used on spectral schemes. This section computes
\(\QCoh(X)\) from affine charts and globalizes the Postnikov structure,
flatness, affine direct image, and quasi-coherent commutative algebra objects.
These are precisely the ingredients needed before coefficient algebra can be
returned to geometry by relative spectrum.

\subsection{Geometric descent of quasi-coherent coefficients}

\begin{theorem}[Atlas computation of quasi-coherent modules]
\label{thm:qcoh-atlas-computation}
Let $X$ be a spectral scheme and let
\[
    U
    \longrightarrow
    X
\]
be an affine Zariski atlas. Write $U_\bullet$ for its \v{C}ech nerve in
spectral schemes. Then the canonical comparison functor is a symmetric
monoidal equivalence
\[
    \QCoh(X)
    \xrightarrow{\simeq}
    \operatorname*{Tot}_{[n]\in\Delta}
    \QCoh(U_n).
\]
Equivalently, the augmented cosimplicial coefficient diagram
\[
\begin{tikzcd}[column sep=large]
    \QCoh(X)
        \arrow[r]
    &
    \QCoh(U)
        \arrow[r, shift left=1.1ex]
        \arrow[r, shift right=1.1ex]
    &
    \QCoh(U\times_XU)
        \arrow[r, shift left=1.8ex]
        \arrow[r]
        \arrow[r, shift right=1.8ex]
    &
    \cdots
\end{tikzcd}
\]
is a limit diagram in symmetric monoidal presentable stable
\(\infty\)-categories.
\end{theorem}

\begin{proof}
\leavevmode
\begin{enumerate}[label=\textup{(\arabic*)}]
    \item \textbf{Restriction to the \v{C}ech nerve.}
    Restriction of an object of \(\QCoh(X)\) to the atlas and all its iterated
    intersections produces an object of the totalization and therefore the
    comparison functor
    \[
        \QCoh(X)
        \longrightarrow
        \operatorname*{Tot}_{[n]\in\Delta}\QCoh(U_n).
    \]

    \item \textbf{Pull the atlas back to an affine point.}
    Let \(v:V=\Spec(A)\to X\) be an affine point. The pullback
    \(U_V:=U\times_XV\to V\) is an effective epimorphism by open
    subschemes.

    \item \textbf{Refine by affine opens.}
    Refine its components by affine open atlases.

    \item \textbf{Recognize an affine Zariski cover.}
    The resulting
    coproduct of affine opens is still an effective epimorphism over \(V\),
    so \Ref{prop:affine-open-cover-effective-epi} identifies the resulting
    affine family as an affine spectral Zariski covering family of \(V\).

    \item \textbf{Choose a finite principal refinement.}
    Hence \Ref{prop:finite-principal-refinement} supplies a finite principal
    refinement
    \[
        \coprod_{i=1}^r
        V_i
        \longrightarrow
        V,
        \qquad
        V_i=D_V^{\mathrm{sp}}(f_i),
        \qquad
        (f_1,\ldots,f_r)=(1).
    \]

    \item \textbf{Descend to an \(A\)-module.}
    Restrict the given object of the \v{C}ech totalization to this finite
    affine cover. By \Ref{thm:affine-qcoh-zariski-descent}, it descends to an
    \(A\)-module \(M_V\).

    \item \textbf{Independence of principal refinement.}
    Suppose \(\{D_V^{\mathrm{sp}}(f_i)\}_i\) and
    \(\{D_V^{\mathrm{sp}}(g_j)\}_j\) are two finite principal refinements.
    Their pairwise intersections are
    \[
        D_V^{\mathrm{sp}}(f_i)
        \times_V
        D_V^{\mathrm{sp}}(g_j)
        \simeq
        D_V^{\mathrm{sp}}(f_ig_j).
    \]
    Since \((f_i)_i=(1)\) and \((g_j)_j=(1)\), the products
    \(\{f_ig_j\}_{i,j}\) also generate the unit ideal. Thus these
    intersections form a common finite principal refinement. Full
    faithfulness in affine module descent identifies the two descended
    modules by a canonical equivalence, and the space of such
    identifications is contractible.

    \item \textbf{Compatibility under affine base change.}
    For a morphism \(V'=\Spec(A')\to V=\Spec(A)\) over \(X\), base change of
    the chosen principal refinement is the finite principal family
    \[
        D_{V'}^{\mathrm{sp}}(f_i)
        \longrightarrow
        V'.
    \]
    The images of the \(f_i\) still generate the unit ideal, so this is again
    a cover. Both
    \[
        M_V\otimes_AA'
        \qquad\text{and}\qquad
        M_{V'}
    \]
    correspond to the same object of its affine descent totalization. Hence
    affine descent yields a canonical equivalence
    \[
        M_V\otimes_AA'\simeq M_{V'}.
    \]
    Full faithfulness of affine descent identifies these transition
    equivalences and all their higher homotopies with the corresponding
    morphisms in the totalization.

    \item \textbf{Assemble the compatible affine-point family.}
    The compatible family \(\{M_V\}\) defines an object of the affine-point
    limit
    \[
        \QCoh(X)
        =
        \varprojlim\nolimits_{(\Spec A\to X)\in\operatorname{AffPt}(X)^{\op}}\Mod_A.
    \]

    \item \textbf{Identify the inverse equivalence.}
    Affine descent shows that this construction and restriction to the \v{C}ech
    nerve are inverse.

    \item \textbf{Preserve the symmetric monoidal structure.}
    All restriction functors are strong symmetric
    monoidal and the defining limits are formed in
    \(\CAlg(\Pr^L_{\mathrm{st}})\), so the equivalence is symmetric
    monoidal.
\end{enumerate}
\end{proof}

\begin{corollary}[Zariski locality of quasi-coherent modules]
\label{cor:qcoh-zariski-locality}
A quasi-coherent module on a spectral scheme is equivalently a family of
modules on the members of an affine Zariski atlas, together with coherent
descent equivalences on all iterated intersections.
\end{corollary}

\begin{proof}
\leavevmode
\begin{enumerate}[label=\textup{(\arabic*)}]
    \item \textbf{Unfold the descent totalization.}
    Unfold the totalization in \Ref{thm:qcoh-atlas-computation}.
\end{enumerate}
\end{proof}

\begin{lemma}[Computation on the affine-open basis]
\label{lem:qcoh-affine-open-basis}
Let $X$ be a spectral scheme, and let
\[
    \operatorname{AffOp}(X)
\]
denote the $\infty$-category of affine open spectral subschemes of $X$.
Restriction induces a symmetric monoidal equivalence
\[
    \QCoh(X)
    \xrightarrow{\simeq}
    \varprojlim\nolimits_
    {V\in\operatorname{AffOp}(X)^{\op}}
    \QCoh(V).
\]
\end{lemma}

\begin{proof}
\leavevmode
\begin{enumerate}[label=\textup{(\arabic*)}]
    \item \textbf{Restriction to affine opens.}
    Restriction gives a symmetric monoidal functor
    \[
        \rho:\QCoh(X)
        \longrightarrow
        \varprojlim\nolimits_{V\in\operatorname{AffOp}(X)^{\op}}\QCoh(V).
    \]

    \item \textbf{Choose an affine atlas and an affine cover of the intersection.}
    Let \(\{M_V\}_V\) be an object of the limit on the right and choose an
    affine open atlas \(\{U_i\to X\}\). The intersection
    \[
        U_{ij}:=U_i\times_XU_j
    \]
    need not be affine. Choose an affine open cover
    \(\{W_\alpha\to U_{ij}\}\).

    \item \textbf{Compare the prescribed modules on the affine cover.}
    Every \(W_\alpha\) is also an affine open
    of \(X\), so the compatibility morphisms of the limit give
    \[
        M_{U_i}|_{W_\alpha}
        \simeq
        M_{W_\alpha}
        \simeq
        M_{U_j}|_{W_\alpha}.
    \]

    \item \textbf{Check compatibility on refinements.}
    On further affine refinements these equivalences agree because they are
    restrictions of one object of the categorical limit.

    \item \textbf{Glue the overlap equivalence.}
    By
    \Ref{thm:qcoh-atlas-computation} applied to \(U_{ij}\), they glue to a
    canonical equivalence
    \[
        M_{U_i}|_{U_{ij}}
        \xrightarrow{\simeq}
        M_{U_j}|_{U_{ij}}.
    \]

    \item \textbf{Higher compatibility.}
    On triple and higher intersections, every required equality between
    composites may be checked after an affine open cover. There it is exactly
    the corresponding equality among the transition maps of the object
    \(\{M_V\}_V\) in the limit. Full faithfulness of affine Zariski descent
    therefore supplies the global equality and all higher coherences. Thus
    the \(M_{U_i}\) define an object of
    \[
        \operatorname*{Tot}_{[n]\in\Delta}\QCoh(U_n).
    \]

    \item \textbf{Glue the descent object.}
    By \Ref{thm:qcoh-atlas-computation}, the resulting descent object glues to
    \(M\in\QCoh(X)\).

    \item \textbf{Compare it with a prescribed affine-open value.}
    Let \(V\subseteq X\) be any affine open. The opens
    \(V\times_XU_i\) cover \(V\); after refining them by affine opens, both
    \(M|_V\) and the prescribed object \(M_V\) restrict to the same members
    of the original compatible family.

    \item \textbf{Apply affine descent on the affine open.}
    Affine descent on \(V\) gives a
    canonical equivalence
    \[
        M|_V\simeq M_V.
    \]

    \item \textbf{Conclude that the constructions are inverse.}
    Hence this construction is inverse to \(\rho\).

    \item \textbf{Monoidality.}
    All restriction and gluing functors used above arise from limits in
    \(\CAlg(\Pr^L_{\mathrm{st}})\). The equivalence is therefore symmetric
    monoidal.
\end{enumerate}
\end{proof}

\subsection{Global flatness and local detection}

\begin{definition}[Flat morphisms and flat quasi-coherent modules]
\label{def:flat-spectral-scheme-morphism}
A morphism
\[
    f
    \colon
    Y
    \longrightarrow
    X
\]
of spectral schemes is \textbf{flat} if, for every affine spectral scheme
\[
    U
    =
    \Spec(A)
    \longrightarrow
    X
\]
and every affine open
\[
    V
    =
    \Spec(B)
    \subseteq
    Y\times_XU,
\]
the induced morphism $A\to B$ is flat.

A quasi-coherent module $\mathcal F\in\QCoh(X)$ is \textbf{flat} if, for
every affine open $U=\Spec(A)\subseteq X$, the corresponding $A$-module
$\mathcal F|_U$ is flat.
\end{definition}

\begin{proposition}[Zariski locality of flatness]
\label{prop:global-flatness-zariski-local}
Flatness of morphisms of spectral schemes is local on the source and the
target for the spectral Zariski topology. Consequently, in
\Ref{def:flat-spectral-scheme-morphism} it is enough to verify flatness on
an affine Zariski atlas of the target and an affine Zariski atlas of each
pullback to the source.
\end{proposition}

\begin{proof}
\leavevmode
\begin{enumerate}[label=\textup{(\arabic*)}]
    \item \textbf{Source locality in the affine case.}
    Let \(A\to B\) be a map of connective commutative ring spectra and let
    \(\{\Spec(B_i)\to\Spec(B)\}\) be a Zariski cover such that every
    composite \(A\to B_i\) is flat. For \(N\in(\Mod_{A})_{\leq0}\), put
    \[
        M:=N\otimes_AB.
    \]
    Then
    \[
        M\otimes_BB_i
        \simeq
        N\otimes_AB_i
        \in
        (\Mod_{B_i})_{\leq0}.
    \]

    \item \textbf{Detect coconnectivity after the cover.}
    The maps \(B\to B_i\) are open and hence flat, so extension of scalars is
    \(t\)-exact. Therefore
    \[
        \tau_{\geq1}^{B}(M)\otimes_BB_i
        \simeq
        \tau_{\geq1}^{B_i}(M\otimes_BB_i)
        \simeq0.
    \]
    Joint conservativity from
    \Ref{prop:zariski-cover-joint-conservativity} gives
    \(\tau_{\geq1}^{B}(M)\simeq0\). Thus \(-\otimes_AB\) preserves
    coconnective modules; right \(t\)-exactness is automatic, so \(A\to B\)
    is flat.

    \item \textbf{Establish target locality.}
    Locality on the target follows from base-change stability of flatness and
    the source-locality just proved.

    \item \textbf{Globalize by affine atlases.}
    Refining source and target by affine open
    atlases yields the stated global locality criterion.
\end{enumerate}
\end{proof}

\begin{proposition}[Zariski locality of flat modules]
\label{prop:flat-module-zariski-local}
Let \(A\) be a connective commutative ring spectrum, let
\[
    \{\Spec(A_i)\to\Spec(A)\}_{i\in I}
\]
be an affine spectral Zariski covering family, and let
\[
    M
    \in
    \Mod_A.
\]
Then \(M\) is flat over \(A\) if and only if
\[
    M\otimes_AA_i
\]
is flat over \(A_i\) for every \(i\).
\end{proposition}

\begin{proof}
\leavevmode
\begin{enumerate}[label=\textup{(\arabic*)}]
    \item \textbf{Base change of a flat module.}
    If \(M\) is flat over \(A\), then each
    \(M\otimes_AA_i\) is flat over \(A_i\) by
    \Ref{prop:flat-module-permanence}.

    \item \textbf{Recover connectivity from the cover.}
    Conversely, put \(M_i:=M\otimes_AA_i\) and assume every \(M_i\) is flat.
    The maps \(A\to A_i\) are open and hence flat, so extension of scalars is
    \(t\)-exact. Since every \(M_i\) is connective,
    \[
        \tau_{\leq-1}^{A}(M)\otimes_AA_i
        \simeq
        \tau_{\leq-1}^{A_i}(M_i)
        \simeq0.
    \]
    Joint conservativity gives \(\tau_{\leq-1}^{A}(M)\simeq0\), so \(M\) is
    connective.

    \item \textbf{Recover left $t$-exactness.}
    Let \(N\in(\Mod_{A})_{\leq0}\). For every \(i\),
    \[
    \begin{aligned}
        (N\otimes_AM)\otimes_AA_i
        &\simeq
        (N\otimes_AA_i)\otimes_{A_i}(M\otimes_AA_i)
        \\
        &\simeq
        (N\otimes_AA_i)\otimes_{A_i}M_i.
    \end{aligned}
    \]
    Flatness of \(A\to A_i\) makes \(N\otimes_AA_i\) coconnective, and
    flatness of \(M_i\) keeps the displayed tensor product coconnective.
    Hence
    \[
        \tau_{\geq1}^{A}(N\otimes_AM)\otimes_AA_i\simeq0.
    \]
    Joint conservativity gives
    \(\tau_{\geq1}^{A}(N\otimes_AM)\simeq0\).

    \item \textbf{Conclude flatness.}
    Tensoring with the connective module \(M\) is right \(t\)-exact, and the
    preceding step proves left \(t\)-exactness. Therefore \(M\) is flat.
\end{enumerate}
\end{proof}

\subsection{The global Postnikov structure and affine direct image}

\begin{theorem}[Canonical $t$-structure on global quasi-coherent modules]
\label{thm:qcoh-spectral-scheme-t-structure}
Let \(X\) be a spectral scheme. The full subcategories
\[
    \QCoh(X)_{\geq0}
    \qquad\text{and}\qquad
    \QCoh(X)_{\leq0}
\]
of quasi-coherent modules which are respectively connective and coconnective
on every affine open define an accessible, left complete, and right complete
\(t\)-structure on \(\QCoh(X)\). Pullback along every morphism of spectral
schemes is right \(t\)-exact, and pullback along a flat morphism is
\(t\)-exact.
\end{theorem}

\begin{proof}
\leavevmode
\begin{enumerate}[label=\textup{(\arabic*)}]
    \item \textbf{Limit presentation.}
    For each affine open write \(V=\Spec(A_V)\). By
    \Ref{lem:qcoh-affine-open-basis},
    \[
        \QCoh(X)
        \simeq
        \varprojlim\nolimits_{V=\Spec(A_V)\in\operatorname{AffOp}(X)^{\op}}
        \Mod_{A_V}.
    \]
    Every transition functor is pullback along an open immersion and is
    \(t\)-exact because open immersions are flat.

    \item \textbf{Construct global truncations componentwise.}
    For a compatible family \(\mathcal F=\{\mathcal F_V\}_V\), apply affine
    Postnikov truncation componentwise.

    \item \textbf{Assemble the truncation triangle.}
    The \(t\)-exact transition functors
    preserve the compatibility maps, so the affine truncation triangles
    assemble to
    \[
    \begin{tikzcd}[column sep=large]
        \tau_{\geq0}^{X}\mathcal F
            \arrow[r]
        &
        \mathcal F
            \arrow[r]
        &
        \tau_{\leq-1}^{X}\mathcal F.
    \end{tikzcd}
    \]

    \item \textbf{Establish global orthogonality.}
    Mapping spectra in the categorical limit are limits of the affine mapping
    spectra, so affine orthogonality implies global orthogonality.

    \item \textbf{Accessibility.}
    Choose a regular cardinal \(\kappa\) for which the small affine-open
    diagram, its transition functors, and the affine truncation functors are
    \(\kappa\)-accessible. The componentwise truncation functors on the limit
    are then \(\kappa\)-accessible.

    \item \textbf{Left completeness.}
    For each affine open \(V\), consider the tower
    \[
        \cdots
        \longrightarrow
        (\Mod_{A_V})_{\leq n}
        \xrightarrow{\tau_{\leq n-1}^{A_V}}
        (\Mod_{A_V})_{\leq n-1}
        \longrightarrow\cdots
    \]
    Affine left completeness gives an equivalence in \(\Cat_\infty\)
    \[
        \Mod_{A_V}
        \simeq
        \varprojlim\nolimits_n
        (\Mod_{A_V})_{\leq n}.
    \]
    Naturality in \(V\) and commutation of limits yield
    \[
    \begin{aligned}
        \QCoh(X)
        &\simeq
        \varprojlim\nolimits_V\varprojlim\nolimits_n
        (\Mod_{A_V})_{\leq n}
        \\
        &\simeq
        \varprojlim\nolimits_n\QCoh(X)_{\leq n}.
    \end{aligned}
    \]

    \item \textbf{Right completeness.}
    The canonical morphism
    \[
        \varinjlim_n\tau_{\geq-n}^{X}\mathcal F
        \longrightarrow
        \mathcal F
    \]
    becomes an equivalence after restriction to every affine open by affine
    right completeness. Affine-open restrictions jointly detect
    equivalences, so it is an equivalence globally.

    \item \textbf{Choose an affine open in the source and an affine atlas in the target.}
    Let \(f:Y\to X\) and let
    \(\mathcal F\in\QCoh(X)_{\geq0}\). Choose an affine open
    \(V=\Spec(B)\subseteq Y\) and an affine open atlas
    \(\{U_i=\Spec(A_i)\to X\}\). The opens
    \[
        V_i:=V\times_XU_i
    \]
    cover \(V\).

    \item \textbf{Refine the pullback cover by affine opens.}
    Refine each \(V_i\) by affine opens
    \(W_{ij}=\Spec(B_{ij})\).

    \item \textbf{Check local connectivity after extension of scalars.}
    On \(W_{ij}\), the restriction of
    \(f^*\mathcal F\) is extension of scalars along
    \(A_i\to B_{ij}\), hence is connective by
    \Ref{prop:base-change-exact-right-t-exact}.

    \item \textbf{Recognize the refinement as a Zariski cover.}
    Let \(M\in\Mod_B\) represent \((f^*\mathcal F)|_V\). The coproduct
    of the affine opens \(W_{ij}\to V\) is an effective epimorphism, so
    \Ref{prop:affine-open-cover-effective-epi} makes them an affine spectral
    Zariski covering family of \(V\). Their structural maps
    \(B\to B_{ij}\) are flat. Therefore
    \[
        \tau_{\leq-1}^{B}(M)\otimes_BB_{ij}
        \simeq
        \tau_{\leq-1}^{B_{ij}}(M\otimes_BB_{ij})
        \simeq0.
    \]

    \item \textbf{Detect global connectivity by joint conservativity.}
    Joint conservativity gives \(\tau_{\leq-1}^{B}(M)\simeq0\). Thus
    \(f^*\mathcal F\) is connective on every affine open of \(Y\), proving
    right \(t\)-exactness.

    \item \textbf{$t$-exactness for flat pullback.}
    If \(f\) is flat, every affine map \(A_i\to B_{ij}\) above is flat, so
    affine extension of scalars is \(t\)-exact. For
    \(\mathcal F\in\QCoh(X)_{\leq0}\), the same argument gives
    \[
        \tau_{\geq1}^{B}(M)\otimes_BB_{ij}
        \simeq
        \tau_{\geq1}^{B_{ij}}(M\otimes_BB_{ij})
        \simeq0,
    \]
    and joint conservativity gives \(\tau_{\geq1}^{B}(M)\simeq0\). Hence flat
    pullback preserves both halves of the global Postnikov \(t\)-structure.
\end{enumerate}
\end{proof}

\begin{definition}[Affine morphism of spectral schemes]
\label{def:affine-morphism-spectral-schemes}
A morphism $p:Y\to X$ of spectral schemes is \textbf{affine} if, for every
affine open $U\to X$, the pullback
\[
    Y\times_XU
\]
is an affine spectral scheme.
\end{definition}

\begin{proposition}[Affine direct image is restriction of scalars]
\label{prop:affine-direct-image-restriction-scalars}
Let
\[
    p:Y\longrightarrow X
\]
be an affine morphism of spectral schemes. For every affine open
\(U=\Spec(A)\subseteq X\), write
\[
    Y_U:=Y\times_XU\simeq\Spec(B).
\]
Then for every \(\mathcal F\in\QCoh(Y)\) there is a canonical equivalence
of \(A\)-modules
\[
    (p_*\mathcal F)|_U
    \simeq
    \operatorname{Res}_A^B
    \bigl(\mathcal F|_{Y_U}\bigr).
\]
These equivalences are compatible with restriction to smaller affine opens.
\end{proposition}

\begin{proof}
\leavevmode
\begin{enumerate}[label=\textup{(\arabic*)}]
    \item \textbf{Compatibility under restriction.}
    Let \(U'=\Spec(A')\to U=\Spec(A)\) be an affine open. Then
    \[
        Y_{U'}
        \simeq
        \Spec(B'),
        \qquad
        B':=B\otimes_AA'.
    \]
    For every \(B\)-module \(M\), associativity of relative tensor products
    gives a canonical equivalence of \(A'\)-modules
    \[
    \begin{aligned}
        \operatorname{Res}_A^B(M)\otimes_AA'
        &\simeq
        M\otimes_AA'
        \\
        &\simeq
        M\otimes_B(B\otimes_AA')
        \\
        &\simeq
        \operatorname{Res}_{A'}^{B'}(M\otimes_BB').
    \end{aligned}
    \]
    The associativity coherence supplies compatibility for chains of affine
    open restrictions.

    \item \textbf{Construct the candidate direct image.}
    The preceding compatibilities show that the modules
    \[
        U=\Spec(A)
        \longmapsto
        \operatorname{Res}_A^B(\mathcal F|_{Y_U})
    \]
    define an object \(G\in\QCoh(X)\) by
    \Ref{lem:qcoh-affine-open-basis}.

    \item \textbf{Apply the affine adjunction chartwise.}
    Let \(\mathcal E\in\QCoh(X)\). Choose an affine open atlas of \(X\)
    and affine refinements of all iterated intersections. On every affine
    chart the affine extension--restriction adjunction gives
    \[
        \Map_{\Mod_A}
        \bigl(\mathcal E|_U,
              \operatorname{Res}_A^B(\mathcal F|_{Y_U})\bigr)
        \simeq
        \Map_{\Mod_B}
        \bigl(\mathcal E|_U\otimes_AB,
              \mathcal F|_{Y_U}\bigr).
    \]

    \item \textbf{Check compatibility with restriction.}
    These equivalences commute with restriction by Step \textup{(1)}.

    \item \textbf{Pass to descent totalizations.}
    Taking the descent totalizations of
    \Ref{thm:qcoh-atlas-computation} yields
    \[
        \Map_{\QCoh(X)}(\mathcal E,G)
        \simeq
        \Map_{\QCoh(Y)}(p^*\mathcal E,\mathcal F).
    \]

    \item \textbf{Identify the right adjoint.}
    Hence \(G\) is right adjoint to \(p^*\). Uniqueness of right adjoints
    identifies \(G\simeq p_*\mathcal F\), proving the formula.
\end{enumerate}
\[
\begin{tikzcd}[column sep=huge, row sep=large]
    \Mod_B
        \arrow[r, "{-\otimes_BB'}"]
        \arrow[d, "{\operatorname{Res}_A^B}"']
        \arrow[dr, phantom, "\scriptstyle\simeq" description]
    &
    \Mod_{B'}
        \arrow[d, "{\operatorname{Res}_{A'}^{B'}}"]
    \\
    \Mod_A
        \arrow[r, "{-\otimes_AA'}"']
    &
    \Mod_{A'}.
\end{tikzcd}
\]
\end{proof}




\subsection{Quasi-coherent commutative algebras}

\begin{definition}[Quasi-coherent commutative algebra]
\label{def:quasi-coherent-commutative-algebra}
For a spectral scheme \(X\), define
\[
    \operatorname{QCAlg}(X)
    :=
    \CAlg\bigl(\QCoh(X)\bigr).
\]
Let
\[
    \operatorname{oblv}^{\mathrm{alg}}_{X}
    \colon
    \operatorname{QCAlg}(X)
    \longrightarrow
    \QCoh(X)
\]
denote the forgetful functor. The full subcategory of algebra objects
whose underlying quasi-coherent modules are connective is denoted
\[
    \operatorname{QCAlg}(X)^{\mathrm{cn}}
    :=
    \left\{
        \mathcal A
        \in
        \operatorname{QCAlg}(X)
        \,\middle|\,
        \operatorname{oblv}^{\mathrm{alg}}_{X}(\mathcal A)
        \in
        \QCoh(X)_{\geq0}
    \right\}.
\]
\end{definition}

\begin{remark}[Connectivity of quasi-coherent algebras]
\label{rem:qcalg-connectivity-not-t-structure}
Connectivity of a quasi-coherent commutative algebra is measured in the
underlying stable coefficient category. Thus the superscript
\({}^{\mathrm{cn}}\) in
\(\operatorname{QCAlg}(X)^{\mathrm{cn}}\) means that the underlying object
of \(\QCoh(X)\) is connective. It does not assert a $t$-structure on the
generally nonstable $\infty$-category \(\operatorname{QCAlg}(X)\); compare
\Ref{rem:postnikov-not-t-structure-on-algebras}.
\end{remark}

\begin{proposition}[Pullback of quasi-coherent algebras]
\label{prop:pullback-qcalg}
Every morphism $f:X\to Y$ of spectral schemes induces a functor
\[
    f^*
    \colon
    \operatorname{QCAlg}(Y)
    \longrightarrow
    \operatorname{QCAlg}(X).
\]
It preserves connective algebra objects, and for composable morphisms
there are canonical coherent equivalences
\[
    (g\circ f)^*
    \simeq
    f^*g^*.
\]
\end{proposition}

\begin{proof}
\leavevmode
\begin{enumerate}[label=\textup{(\arabic*)}]
    \item \textbf{Pass to commutative algebra objects.}
    The quasi-coherent pullback functor is strong symmetric monoidal by
    \Ref{prop:qcoh-pullback}, hence carries commutative algebra objects to
    commutative algebra objects.

    \item \textbf{Preserve connectivity.}
    Right $t$-exactness gives preservation of
    connectivity.

    \item \textbf{Inherit coherence.}
    Coherence is inherited from the functoriality of
    quasi-coherent pullback.
\end{enumerate}
\end{proof}

\begin{proposition}[Zariski descent for quasi-coherent algebras]
\label{prop:qcalg-zariski-descent}
Let \(X\) be a spectral scheme and let \(U\to X\) be an affine Zariski
atlas with \v{C}ech nerve \(U_\bullet\). Then restriction induces an
equivalence
\[
    \operatorname{QCAlg}(X)
    \xrightarrow{\simeq}
    \operatorname*{Tot}_{[n]\in\Delta}
    \operatorname{QCAlg}(U_n).
\]
It restricts to an equivalence on connective algebra objects:
\[
    \operatorname{QCAlg}(X)^{\mathrm{cn}}
    \xrightarrow{\simeq}
    \operatorname*{Tot}_{[n]\in\Delta}
    \operatorname{QCAlg}(U_n)^{\mathrm{cn}}.
\]
\end{proposition}

\[
\begin{tikzcd}[column sep=large, row sep=large]
    \QCoh(X)
        \arrow[rr, "\simeq"]
        \arrow[d, "\CAlg"']
    &&
    \displaystyle
    \operatorname*{Tot}_{[n]\in\Delta}\QCoh(U_n)
        \arrow[d, "\CAlg"]
    \\
    \operatorname{QCAlg}(X)
        \arrow[rr, "\simeq"']
    &&
    \displaystyle
    \CAlg\!\left(
        \operatorname*{Tot}_{[n]\in\Delta}\QCoh(U_n)
    \right)
        \arrow[r, "\simeq"']
    &
    \displaystyle
    \operatorname*{Tot}_{[n]\in\Delta}\operatorname{QCAlg}(U_n).
\end{tikzcd}
\]


\begin{proof}
\leavevmode
\begin{enumerate}[label=\textup{(\arabic*)}]
    \item \textbf{Apply commutative algebra objects to QCoh descent.}
    By \Ref{thm:qcoh-atlas-computation},
    \[
        \QCoh(X)
        \xrightarrow{\simeq}
        \operatorname*{Tot}_{[n]\in\Delta}\QCoh(U_n)
    \]
    is an equivalence of symmetric monoidal \(\infty\)-categories. The
    functor assigning commutative algebra objects preserves limits of
    symmetric monoidal \(\infty\)-categories; see
    \cite{LurieHigherAlgebra}. Hence
    \[
        \CAlg(\QCoh(X))
        \xrightarrow{\simeq}
        \operatorname*{Tot}_{[n]\in\Delta}
        \CAlg(\QCoh(U_n)).
    \]
    This is the first assertion.

    \item \textbf{Restrict to connective objects.}
    Connectivity of a quasi-coherent module is detected on affine opens by
    \Ref{thm:qcoh-spectral-scheme-t-structure}. Therefore an algebra object
    is connective exactly when all its restrictions to the affine atlas are
    connective, and the equivalence restricts as asserted.
\end{enumerate}
\end{proof}

\begin{proposition}[Lax monoidality of direct image]
\label{prop:qcoh-direct-image-lax-monoidal}
For every morphism \(f:X\to Y\) of spectral schemes, the right adjoint
\[
    f_*:\QCoh(X)\longrightarrow\QCoh(Y)
\]
carries a canonical lax symmetric monoidal structure. Consequently there is
an adjunction on commutative algebra objects
\[
\begin{tikzcd}[column sep=huge]
    \operatorname{QCAlg}(Y)
        \arrow[r, shift left=1.1ex, "f^*"]
    &
    \operatorname{QCAlg}(X)
        \arrow[l, swap, shift left=1.1ex, "f_*"'].
\end{tikzcd}
\]
If \(f=p\) is affine and \(U=\Spec(A)\subseteq Y\) has pullback
\(X\times_YU\simeq\Spec(B)\), then the commutative algebra
\(p_*\mathcal O_X\) restricts to \(B\) as a commutative \(A\)-algebra.
\end{proposition}

\begin{proof}
\leavevmode
\begin{enumerate}[label=\textup{(\arabic*)}]
    \item \textbf{Construct the lax unit.}
    The morphism
    \[
        \mathcal O_Y\longrightarrow f_*\mathcal O_X
    \]
    is the adjoint of the canonical equivalence
    \[
        f^*\mathcal O_Y\simeq\mathcal O_X.
    \]

    \item \textbf{Construct the lax tensor product.}
    For \(\mathcal F,\mathcal G\in\QCoh(X)\), define
    \[
        f_*\mathcal F\otimes f_*\mathcal G
        \longrightarrow
        f_*(\mathcal F\otimes\mathcal G)
    \]
    to be adjoint to the composite
    \[
    \begin{aligned}
        f^*(f_*\mathcal F\otimes f_*\mathcal G)
        &\simeq
        f^*f_*\mathcal F\otimes f^*f_*\mathcal G
        \\
        &\longrightarrow
        \mathcal F\otimes\mathcal G,
    \end{aligned}
    \]
    where the final morphism is the tensor product of the two adjunction
    counits.

    \item \textbf{Obtain coherence by adjoint mates.}
    Associativity, commutativity, and unitality are the mates under the
    adjunction of the corresponding identities for the strong symmetric
    monoidal functor \(f^*\).

    \item \textbf{Identify the canonical lax monoidal structure.}
    This is the canonical lax symmetric monoidal
    structure on the right adjoint of a strong symmetric monoidal left
    adjoint; see \cite{LurieHigherAlgebra}. Thus no independent coherence
    datum is chosen.

    \item \textbf{Lift to commutative algebra objects.}
    A lax symmetric monoidal functor carries commutative algebra objects to
    commutative algebra objects, and the original adjunction lifts to the
    displayed adjunction on \(\operatorname{QCAlg}\).

    \item \textbf{Affine identification on the structure sheaf.}
    Suppose \(f=p\) is affine and, over \(U=\Spec(A)\subseteq Y\), write
    \(X\times_YU\simeq\Spec(B)\). By
    \Ref{prop:affine-direct-image-restriction-scalars}, the underlying
    \(A\)-module of \((p_*\mathcal O_X)|_U\) is \(B\). Under the affine
    extension--restriction adjunction, the lax multiplication constructed in
    Step \textup{(2)} is the ordinary multiplication
    \[
        B\otimes_AB\longrightarrow B,
    \]
    and the lax unit is \(A\to B\). Hence the induced commutative
    \(A\)-algebra is canonically the original algebra \(B\).
\end{enumerate}
\end{proof}

The global coefficient system can now be computed from any affine atlas and
carries the structures needed for geometric re-entry: a global Postnikov
structure, local flatness, affine direct image, and quasi-coherent
commutative algebras satisfying descent. Relative spectrum can therefore be
introduced as the construction that represents connective algebra in a fibre
by affine geometry over the base.

\section{Affine comprehension}
\label{sec:affine-comprehension}

The indexed geometry now has all three ingredients needed for
representability: a geometric base of spectral schemes, stable coefficient
fibres varying by substitution, and Zariski descent. What remains is the
return from algebra in a coefficient fibre to geometry over its context.
Relative spectrum provides that return. For a spectral scheme \(X\),
connective objects of \(\operatorname{QCAlg}(X)\) determine affine
extensions of \(X\), and their algebra terms are represented after arbitrary
substitution.

\begin{remark}[Relative algebraic geometry in context]
\label{rem:relative-algebraic-geometry-in-context}
Affine comprehension repeats the absolute algebra-to-geometry construction
inside each context. The absolute passage
\[
    \Sp_{\geq0}^\otimes
    \longrightarrow
    \CAlg(\Sp)^{\mathrm{cn}}
    \longrightarrow
    \Aff_{\Sp}
\]
becomes fibrewise
\[
    \QCoh(X)^\otimes
    \longrightarrow
    \operatorname{QCAlg}(X)^{\mathrm{cn}}
    \longrightarrow
    \operatorname{QCAlg}(X)^{\mathrm{cn},\op}
    \simeq
    \bigl(\operatorname{Sch}^{\mathrm{sp}}_{/X}\bigr)_{\mathrm{aff}}.
\]
The relative spectrum theorem identifies the final term with the actual
affine slice geometry over \(X\). The displayed parallel is therefore both an
analogy and a construction.
\end{remark}



\begin{theorem}[Relative spectrum construction]
\label{thm:relative-spectrum-construction}
Let \(X\) be a spectral scheme. There is a functor
\[
    \Spec_X
    \colon
    \operatorname{QCAlg}(X)^{\mathrm{cn},\op}
    \longrightarrow
    \operatorname{Sch}^{\mathrm{sp}}_{/X}
\]
characterized by the following properties.
\begin{enumerate}[label=(\roman*)]
    \item If \(X=\Spec(A)\) and \(\mathcal A\) corresponds to a connective
    commutative \(A\)-algebra \(B\), then
    \[
        \Spec_X(\mathcal A)\simeq\Spec(B).
    \]

    \item For every morphism \(f:Y\to X\), there is a natural equivalence
    \[
        \Map_X\bigl(Y,\Spec_X(\mathcal A)\bigr)
        \simeq
        \Map_{\operatorname{QCAlg}(Y)}
        \bigl(f^*\mathcal A,\mathcal O_Y\bigr).
    \]
\end{enumerate}
\end{theorem}

\begin{proof}
\leavevmode
\begin{enumerate}[label=\textup{(\arabic*)}]
    \item \textbf{Define pullback on algebra objects.}
    Let \(T\) be an affine spectral scheme. Pullback defines a functor
    \[
        X(T)
        \longrightarrow
        \operatorname{QCAlg}(T)^\simeq,
        \qquad
        f\longmapsto f^*\mathcal A.
    \]

    \item \textbf{Define the contravariant mapping-space functor.}
    The functor
    \[
        \operatorname{QCAlg}(T)^{\simeq,\op}
        \longrightarrow
        \mathcal S,
        \qquad
        \mathcal B
        \longmapsto
        \Map_{\operatorname{QCAlg}(T)}(\mathcal B,\mathcal O_T),
    \]
    is contravariant in the algebra variable.

    \item \textbf{Use inversion to obtain the fibre functor.}
    Since \(X(T)\) is an
    \(\infty\)-groupoid, every path \(e:f\to f'\) induces an equivalence
    \[
        e^*:f^*\mathcal A\xrightarrow{\simeq}f'^*\mathcal A,
    \]
    and transport from the fibre over \(f\) to the fibre over \(f'\) is
    precomposition with \((e^*)^{-1}\). Equivalently, compose the
    contravariant mapping-space functor with the inversion equivalence
    \(X(T)\simeq X(T)^{\op}\). This gives a functor
    \[
        \Psi_{\mathcal A,T}
        \colon
        X(T)\longrightarrow\mathcal S,
        \qquad
        f\longmapsto
        \Map_{\operatorname{QCAlg}(T)}
        (f^*\mathcal A,\mathcal O_T).
    \]

    \item \textbf{Unstraighten the fibre functor.}
    Unstraighten \(\Psi_{\mathcal A,T}\). Let
    \[
        p_T:F_{\mathcal A}(T)\longrightarrow X(T)
    \]
    be the left fibration classified by this functor.

    \item \textbf{Identify the total space as a colimit.}
    Its total space is
    canonically the colimit
    \[
        F_{\mathcal A}(T)
        \simeq
        \operatorname*{colim}_{f\in X(T)}
        \Psi_{\mathcal A,T}(f).
    \]

    \item \textbf{Describe its objects and coherence.}
    Thus an object of the total space is precisely the Grothendieck-construction
    object represented by a pair
    \[
        (f,\alpha),
        \qquad
        f:T\to X,
        \quad
        \alpha:f^*\mathcal A\to\mathcal O_T,
    \]
    with all path and higher-coherence information supplied by
    unstraightening rather than imposed separately.

    \item \textbf{Define pullback on pairs.}
    For \(g:T'\to T\), pullback sends
    \[
        (f,\alpha)
        \longmapsto
        (f\circ g,g^*\alpha).
    \]

    \item \textbf{Use the coherent pullback equivalences.}
    This uses the coherent equivalences
    \[
        g^*f^*\mathcal A\simeq(f\circ g)^*\mathcal A,
        \qquad
        g^*\mathcal O_T\simeq\mathcal O_{T'}.
    \]

    \item \textbf{Assemble a prestack over \(X\).}
    Naturality of unstraightening therefore assembles the spaces
    \(F_{\mathcal A}(T)\) into a spectral prestack
    \[
        F_{\mathcal A}\longrightarrow X.
    \]

    \item \textbf{Establish contravariant functoriality in the algebra.}
    A morphism \(\mathcal A\to\mathcal B\) acts by precomposition on the
    mapping-space fibres, so \(\mathcal A\mapsto F_{\mathcal A}\) is
    contravariantly functorial.

    \item \textbf{Specify descent data on a finite affine cover.}
    Let \(T_\bullet\to T\) be the \v{C}ech nerve of a finite affine spectral
    Zariski cover. An object of
    \[
        \operatorname*{Tot}_{[n]\in\Delta}F_{\mathcal A}(T_n)
    \]
    consists of a coherent family of maps \(f_n:T_n\to X\) and compatible
    algebra maps
    \[
        \alpha_n:f_n^*\mathcal A\longrightarrow\mathcal O_{T_n}.
    \]

    \item \textbf{Glue the maps to \(X\).}
    Since \(X\) is a spectral Zariski sheaf, the \(f_n\) glue uniquely to
    \(f:T\to X\).

    \item \textbf{Apply algebra descent.}
    By
    \Ref{prop:qcalg-zariski-descent},
    \[
        \operatorname{QCAlg}(T)
        \simeq
        \operatorname*{Tot}_{[n]\in\Delta}
        \operatorname{QCAlg}(T_n).
    \]

    \item \textbf{Glue the algebra maps.}
    Mapping spaces in a limit of \(\infty\)-categories are limits of the
    corresponding mapping spaces, so the \(\alpha_n\) glue uniquely to
    \[
        \alpha:f^*\mathcal A\longrightarrow\mathcal O_T.
    \]

    \item \textbf{Obtain descent for the total space.}
    Hence
    \[
        F_{\mathcal A}(T)
        \xrightarrow{\simeq}
        \operatorname*{Tot}_{[n]\in\Delta}F_{\mathcal A}(T_n).
    \]

    \item \textbf{Conclude the prestack is a Zariski stack.}
    Finite principal covers generate the topology by
    \Ref{prop:finite-principal-refinement}; therefore
    \(F_{\mathcal A}\) is a spectral Zariski stack.

    \item \textbf{Choose an affine base and algebra.}
    Let \(U=\Spec(A)\to X\) be affine and let
    \(\mathcal A|_U\) correspond to a connective commutative
    \(A\)-algebra \(B\).

    \item \textbf{Compute the fibre on an affine test.}
    For \(T=\Spec(C)\to U\), the fibre over the
    chosen map \(T\to U\to X\) is
    \[
    \begin{aligned}
        \Map_{\operatorname{QCAlg}(T)}
        \bigl((\mathcal A|_U)|_T,\mathcal O_T\bigr)
        &\simeq
        \Map_{\CAlg(\Sp)_{C/}}(B\otimes_AC,C)
        \\
        &\simeq
        \Map_{\CAlg(\Sp)_{A/}}(B,C)
        \\
        &\simeq
        \Map_U(T,\Spec B).
    \end{aligned}
    \]

    \item \textbf{Apply Yoneda over the affine base.}
    Naturality in \(T\to U\) and Yoneda give a canonical equivalence
    \[
    \begin{tikzcd}[column sep=large, row sep=large]
        F_{\mathcal A}\times_XU
            \arrow[r, "\simeq"]
            \arrow[d]
        &
        \Spec(B)
            \arrow[d]
        \\
        U
            \arrow[r, equal]
        &
        U.
    \end{tikzcd}
    \]

    \item \textbf{Pull back an affine atlas.}
    Choose an affine spectral Zariski atlas
    \(\coprod_iU_i\to X\). Pulling it back gives an effective epimorphism
    \[
        \coprod_i(F_{\mathcal A}\times_XU_i)
        \longrightarrow
        F_{\mathcal A}.
    \]

    \item \textbf{Identify the pullback charts as affine opens.}
    Every source component is affine by Step \textup{(19)}, and its map to
    \(F_{\mathcal A}\) is an open immersion by base change. Hence
    \(F_{\mathcal A}\) is a spectral scheme over \(X\).

    \item \textbf{Define the relative spectrum.}
    Define
    \[
        \Spec_X(\mathcal A):=F_{\mathcal A}.
    \]

    \item \textbf{Define the two presheaves on open subschemes.}
    Fix \(f:Y\to X\). On the Zariski site of open spectral subschemes of
    \(Y\), define
    \[
        \mathscr P(V)
        :=
        \Map_X(V,\Spec_X\mathcal A)
    \]
    and
    \[
        \mathscr Q(V)
        :=
        \Map_{\operatorname{QCAlg}(V)}
        \bigl((f|_V)^*\mathcal A,\mathcal O_V\bigr).
    \]

    \item \textbf{Verify that \(\mathscr P\) is a sheaf.}
    The functor \(\mathscr P\) is a Zariski sheaf because
    \(\Spec_X\mathcal A\) is a spectral Zariski stack.

    \item \textbf{Verify that \(\mathscr Q\) is a sheaf.}
    To verify that
    \(\mathscr Q\) is a Zariski sheaf, it is enough to
    check descent on affine open covers, since affine opens form a basis.
    For such a cover, \Ref{prop:qcalg-zariski-descent} identifies the algebra
    category on \(V\) with the corresponding totalization, and mapping spaces
    in that totalization are the totalizations of the local mapping spaces.
    Hence \(\mathscr Q\) satisfies Zariski descent.

    \item \textbf{Compare the two sheaves on affine opens.}
    If \(V\subseteq Y\) is affine, Step \textup{(19)} gives a natural
    equivalence
    \[
        \mathscr P(V)\simeq\mathscr Q(V).
    \]

    \item \textbf{Extend the comparison from the affine basis.}
    Affine opens form a basis of the Zariski topology, so these affine
    equivalences determine an equivalence of sheaves
    \(\mathscr P\simeq\mathscr Q\).

    \item \textbf{Evaluate on \(Y\).}
    Evaluating on \(Y\) gives
    \[
        \Map_X\bigl(Y,\Spec_X(\mathcal A)\bigr)
        \simeq
        \Map_{\operatorname{QCAlg}(Y)}
        \bigl(f^*\mathcal A,\mathcal O_Y\bigr).
    \]

    \item \textbf{Record naturality.}
    Naturality in \(Y\), \(f\), and \(\mathcal A\) is inherited from the
    construction.

    \item \textbf{Record uniqueness.}
    Yoneda gives uniqueness of the represented object.
\end{enumerate}
\end{proof}

\begin{proposition}[Universal variable and term abstraction]
\label{prop:relative-spectrum-universal-variable}
Let \(\mathcal A\in\operatorname{QCAlg}(X)^{\mathrm{cn}}\) and write
\[
    \pi_{\mathcal A}:
    X.\mathcal A:=\Spec_X(\mathcal A)
    \longrightarrow
    X.
\]
There is a canonical universal algebra term
\[
    \mathsf{var}_{\mathcal A}:
    \pi_{\mathcal A}^*\mathcal A
    \longrightarrow
    \mathcal O_{X.\mathcal A}.
\]
For every \(f:T\to X\), composition with the relative-spectrum universal
property gives a natural equivalence between
\begin{enumerate}[label=(\roman*)]
    \item algebra terms \(a:f^*\mathcal A\to\mathcal O_T\), and
    \item morphisms \(\langle f,a\rangle:T\to X.\mathcal A\) over \(X\),
\end{enumerate}
characterized by
\[
    \langle f,a\rangle^*\mathsf{var}_{\mathcal A}
    \simeq
    a.
\]
Under substitution \(g:T'\to T\), the term \(a\) is carried to \(g^*a\),
and there is a canonical coherent equivalence
\[
    \langle f\circ g,g^*a\rangle
    \simeq
    \langle f,a\rangle\circ g.
\]
\end{proposition}

\begin{proof}
\leavevmode
\begin{enumerate}[label=\textup{(\arabic*)}]
    \item \textbf{Construct the universal variable.}
    Apply the equivalence of
    \Ref{thm:relative-spectrum-construction} with
    \(Y=X.\mathcal A\) and \(f=\pi_{\mathcal A}\). The identity map
    \[
        \id_{X.\mathcal A}
        \in
        \Map_X(X.\mathcal A,X.\mathcal A)
    \]
    corresponds to a canonical map
    \[
        \pi_{\mathcal A}^*\mathcal A
        \longrightarrow
        \mathcal O_{X.\mathcal A},
    \]
    which is \(\mathsf{var}_{\mathcal A}\).

    \item \textbf{Abstract a term.}
    For \(a:f^*\mathcal A\to\mathcal O_T\), let
    \(\langle f,a\rangle\) be its image under the inverse of the same
    representing equivalence. Naturality of the representing equivalence
    identifies pullback of \(\mathsf{var}_{\mathcal A}\) with \(a\), and
    also gives the displayed compatibility under \(g:T'\to T\). All higher
    coherence is the naturality coherence of the relative-spectrum
    equivalence.
\end{enumerate}
\end{proof}

\begin{remark}[Context extension and terms]
\label{rem:relative-spectrum-context-extension-terms}
The universal property of relative spectrum has the expected comprehension
form. For \(\mathcal A\in\operatorname{QCAlg}(X)^{\mathrm{cn}}\), write
\[
    \pi_{\mathcal A}
    :
    \Spec_X(\mathcal A)
    \longrightarrow
    X.
\]
After substitution along \(f:Y\to X\), maps into this context extension are
precisely algebra terms:
\[
    Y\longrightarrow\Spec_X(\mathcal A)
    \quad\text{over }X
    \qquad\Longleftrightarrow\qquad
    f^*\mathcal A\longrightarrow\mathcal O_Y.
\]
For a term \(a:f^*\mathcal A\to\mathcal O_T\), the corresponding
term-abstraction map is represented by
\[
\begin{tikzcd}[column sep=huge, row sep=large]
    T
        \arrow[rr, "{\langle f,a\rangle}"]
        \arrow[dr, "f"']
    &&
    X.\mathcal A
        \arrow[dl, "{\pi_{\mathcal A}}"]
    \\
    & X &
\end{tikzcd}
\]
together with
\(\langle f,a\rangle^*\mathsf{var}_{\mathcal A}\simeq a\).
In particular, sections of \(\pi_{\mathcal A}\) correspond to algebra maps
\(\mathcal A\to\mathcal O_X\), and the identity of the comprehending
object corresponds to the universal variable of
\Ref{prop:relative-spectrum-universal-variable}. The word
\textbf{comprehension} refers to this representability statement; it does not
identify \(\operatorname{QCAlg}(X)\) with a literal Lawvere predicate fibre.
\end{remark}

\begin{proposition}[Base change for relative spectra]
\label{prop:relative-spectrum-base-change}
Let $f:Y\to X$ be a morphism of spectral schemes and let
$\mathcal A\in\operatorname{QCAlg}(X)^{\mathrm{cn}}$. There is a canonical
equivalence over $Y$
\[
    \Spec_X(\mathcal A)
    \times_X
    Y
    \simeq
    \Spec_Y(f^*\mathcal A).
\]
Equivalently, the canonical square
\[
\begin{tikzcd}[column sep=large, row sep=large]
    \Spec_Y(f^*\mathcal A)
        \arrow[r]
        \arrow[d]
    &
    \Spec_X(\mathcal A)
        \arrow[d, "\pi_{\mathcal A}"]
    \\
    Y
        \arrow[r, "f"']
    &
    X
\end{tikzcd}
\]
is Cartesian.
\end{proposition}

\begin{proof}
\leavevmode
\begin{enumerate}[label=\textup{(\arabic*)}]
    \item \textbf{Compare the two functors of points.}
    For every \(T\to Y\),
    \[
    \begin{aligned}
        \Map_Y\bigl(T,\Spec_X(\mathcal A)\times_XY\bigr)
        &\simeq
        \Map_X\bigl(T,\Spec_X(\mathcal A)\bigr)
        \\
        &\simeq
        \Map_{\operatorname{QCAlg}(T)}
        \bigl(\mathcal A|_T,\mathcal O_T\bigr)
        \\
        &\simeq
        \Map_Y\bigl(T,\Spec_Y(f^*\mathcal A)\bigr).
    \end{aligned}
    \]
    The middle equivalence is the relative-spectrum universal property.

    \item \textbf{Apply Yoneda.}
    The comparison is natural in \(T\), so Yoneda gives the canonical
    equivalence over \(Y\).
\end{enumerate}
\end{proof}

\begin{theorem}[Affine comprehension and classification of affine morphisms]
\label{thm:classification-affine-morphisms}
For every spectral scheme \(X\), relative spectrum induces an equivalence
\[
    \operatorname{QCAlg}(X)^{\mathrm{cn},\op}
    \xrightarrow{\simeq}
    \bigl(\operatorname{Sch}^{\mathrm{sp}}_{/X}\bigr)_{\mathrm{aff}}
\]
onto the full subcategory of affine morphisms over \(X\). Under this
equivalence, for every morphism \(f:Y\to X\), pullback of affine morphisms
corresponds to reindexing of quasi-coherent algebras:
\[
    \Spec_X(\mathcal A)\times_XY
    \simeq
    \Spec_Y(f^*\mathcal A).
\]
Thus affine comprehension is stable under substitution in the Cartesian square
\[
\begin{tikzcd}[column sep=large, row sep=large]
    \Spec_Y(f^*\mathcal A)
        \arrow[r]
        \arrow[d]
    &
    \Spec_X(\mathcal A)
        \arrow[d]
    \\
    Y
        \arrow[r, "f"']
    &
    X.
\end{tikzcd}
\]
\end{theorem}

\[
\begin{tikzcd}[column sep=huge, row sep=large]
    \operatorname{QCAlg}(X)^{\mathrm{cn},\op}
        \arrow[r, "{\Spec_X}"]
        \arrow[d, "{(f^*)^{\op}}"']
    &
    \bigl(\operatorname{Sch}^{\mathrm{sp}}_{/X}\bigr)_{\mathrm{aff}}
        \arrow[d, "{-\times_XY}"]
    \\
    \operatorname{QCAlg}(Y)^{\mathrm{cn},\op}
        \arrow[r, "{\Spec_Y}"']
    &
    \bigl(\operatorname{Sch}^{\mathrm{sp}}_{/Y}\bigr)_{\mathrm{aff}}.
\end{tikzcd}
\]

\begin{proof}
\leavevmode
\begin{enumerate}[label=\textup{(\arabic*)}]
    \item \textbf{Attach the direct-image algebra.}
    Let \(p:Y\to X\) be affine. By
    \Ref{prop:qcoh-direct-image-lax-monoidal}, the structure sheaf determines
    a commutative algebra
    \[
        \mathcal A_p
        :=
        p_*\mathcal O_Y
        \in
        \operatorname{QCAlg}(X).
    \]

    \item \textbf{Identify it on an affine open.}
    For an affine open \(U=\Spec(A)\subseteq X\), write
    \[
        Y_U\simeq\Spec(B).
    \]
    By \Ref{prop:qcoh-direct-image-lax-monoidal},
    \[
        \mathcal A_p|_U\simeq B
    \]
    as a commutative \(A\)-algebra.

    \item \textbf{Detect connectivity locally.}
    Since \(B\) is connective,
    connectivity of \(\mathcal A_p\) is detected on affine opens, and
    \[
        \mathcal A_p\in\operatorname{QCAlg}(X)^{\mathrm{cn}}.
    \]

    \item \textbf{Apply the counit on algebra objects.}
    The counit of the lifted algebra adjunction gives
    \[
        p^*p_*\mathcal O_Y
        \longrightarrow
        \mathcal O_Y
    \]
    in \(\operatorname{QCAlg}(Y)\).

    \item \textbf{Use the relative-spectrum universal property.}
    The universal property of relative
    spectrum therefore determines a morphism
    \[
        \eta_p:
        Y\longrightarrow\Spec_X(\mathcal A_p)
    \]
    over \(X\).

    \item \textbf{Identify the comparison after affine base change.}
    After base change to \(U=\Spec(A)\subseteq X\), the preceding morphism is
    the map
    \[
        \Spec(B)\longrightarrow\Spec(B)
    \]
    corresponding contravariantly to the identity of \(B\).

    \item \textbf{Conclude local equivalence.}
    Hence it is an
    equivalence after base change to an affine open atlas of \(X\).

    \item \textbf{Match the local inverses on overlaps.}
    The
    local inverses are uniquely characterized as inverses to the restrictions
    of \(\eta_p\), so their restrictions to all iterated intersections agree.

    \item \textbf{Glue the local inverses.}
    Mapping descent,
    \Ref{lem:mapping-descent-effective-epimorphism}, glues them to a global
    inverse. Thus:
    \[
        Y\xrightarrow{\simeq}\Spec_X(p_*\mathcal O_Y).
    \]

    \item \textbf{Conclude essential surjectivity.}
    This proves essential surjectivity.

    \item \textbf{Identify the direct-image algebra locally.}
    Let
    \[
        q_{\mathcal B}:\Spec_X(\mathcal B)\longrightarrow X
    \]
    with \(\mathcal B\in\operatorname{QCAlg}(X)^{\mathrm{cn}}\). On every
    affine open \(U=\Spec(A)\), if \(\mathcal B|_U\) corresponds to the
    connective \(A\)-algebra \(B\), then
    \Ref{prop:qcoh-direct-image-lax-monoidal} gives an equivalence of
    commutative \(A\)-algebras
    \[
        \bigl(q_{\mathcal B,*}\mathcal O_{\Spec_X\mathcal B}\bigr)|_U
        \simeq
        B
        \simeq
        \mathcal B|_U.
    \]

    \item \textbf{Check compatibility under restriction.}
    The equivalences are compatible under restriction.

    \item \textbf{Glue the local algebra equivalences.}
    The affine-open
    basis theorem therefore gives a canonical algebra equivalence
    \[
        q_{\mathcal B,*}\mathcal O_{\Spec_X\mathcal B}
        \simeq
        \mathcal B.
    \]

    \item \textbf{Compute the mapping-space equivalence.}
    For connective algebras \(\mathcal A,\mathcal B\), the relative-spectrum
    universal property and the algebra adjunction give
    \[
    \begin{aligned}
        \Map_X
        \bigl(
            \Spec_X\mathcal B,
            \Spec_X\mathcal A
        \bigr)
        &\simeq
        \Map_{\operatorname{QCAlg}(\Spec_X\mathcal B)}
        \bigl(
            q_{\mathcal B}^*\mathcal A,
            \mathcal O_{\Spec_X\mathcal B}
        \bigr)
        \\
        &\simeq
        \Map_{\operatorname{QCAlg}(X)}
        \bigl(
            \mathcal A,
            q_{\mathcal B,*}\mathcal O_{\Spec_X\mathcal B}
        \bigr)
        \\
        &\simeq
        \Map_{\operatorname{QCAlg}(X)}(\mathcal A,\mathcal B).
    \end{aligned}
    \]

    \item \textbf{Conclude full faithfulness.}
    Thus relative spectrum is fully faithful.

    \item \textbf{Combine with essential surjectivity.}
    Together with
    Step \textup{(10)}, this proves the asserted anti-equivalence, with
    explicit inverse
    \[
        p:Y\to X
        \longmapsto
        p_*\mathcal O_Y.
    \]
\end{enumerate}
\end{proof}

Affine comprehension is the main structural return from coefficients to
geometry. Connective quasi-coherent commutative algebras are exactly affine
extensions of the base, and this identification commutes with substitution.
The remaining sections exploit this representability rather than adding a new
layer to the core construction.

\section{Closed context geometry, complements, and gluing}
\label{sec:closed-immersions-complements}

Affine comprehension supplies geometric extensions of a context. To manipulate
those extensions internally, we next record the complementary geometric
operations used later: closed subcontexts, their open complements, and gluing
along closed immersions. Each is constructed from the spectral geometry
already developed, with \(\pi_0\) used locally to detect the relevant closed
condition but without introducing a global classical truncation.

\subsection{Closed immersions}

\begin{definition}[Affine closed immersion]
\label{def:affine-closed-immersion-spectral}
A morphism
\[
    \Spec(B)
    \longrightarrow
    \Spec(A)
\]
of affine spectral schemes is a \textbf{closed immersion} if the induced
ordinary ring homomorphism
\[
    \pi_0(A)
    \longrightarrow
    \pi_0(B)
\]
is surjective.
\end{definition}

\begin{definition}[Closed immersion of spectral schemes]
\label{def:closed-immersion-spectral-schemes}
A morphism
\[
    i
    \colon
    Z
    \longrightarrow
    X
\]
of spectral schemes is a \textbf{closed immersion} if, after base change
along every affine open $U\to X$, the pullback
\[
    Z\times_XU
    \longrightarrow
    U
\]
is an affine closed immersion.
\end{definition}

\begin{proposition}[Permanence of closed immersions]
\label{prop:closed-immersion-permanence}
Closed immersions of spectral schemes are stable under equivalence,
composition, and base change.
\end{proposition}

\begin{proof}
\leavevmode
\begin{enumerate}[label=\textup{(\arabic*)}]
    \item \textbf{Reduce to the affine target.}
    The assertion is local on the target.

    \item \textbf{Handle equivalences.}
    On affines, equivalences preserve
    surjectivity on \(\pi_0\).

    \item \textbf{Handle composition.}
    A composite of surjective ring maps is
    surjective.

    \item \textbf{Compute degree zero after base change.}
    Let \(A\to B\) be surjective on \(\pi_0\) and let \(A\to A'\) be any
    morphism of connective commutative ring spectra. By
    \Ref{prop:classical-truncation-affine-base-change},
    \[
        \pi_0(B\otimes_AA')
        \simeq
        \pi_0(B)\otimes_{\pi_0(A)}\pi_0(A').
    \]

    \item \textbf{Preserve surjectivity under base change.}
    The induced map
    \(\pi_0(A')\to\pi_0(B\otimes_AA')\) is the base change of the
    surjection \(\pi_0(A)\to\pi_0(B)\), hence is surjective.

    \item \textbf{Globalize from the affine-local definition.}
    Thus affine
    closed immersions are stable under base change, and the global assertion
    follows from the affine-local definition.
\end{enumerate}
\end{proof}

\subsection{Open complements}

\begin{lemma}[Vanishing criterion for spectral schemes]
\label{lem:empty-spectral-scheme-classical-detection}
Let \(A\) be a connective commutative ring spectrum. If
\(\pi_0(A)=0\), then \(A\simeq0\) and \(\Spec(A)\simeq\varnothing\).
More generally, a spectral scheme is empty if an affine spectral Zariski
atlas consists entirely of empty affine schemes.
\end{lemma}

\begin{proof}
\leavevmode
\begin{enumerate}[label=\textup{(\arabic*)}]
    \item \textbf{Affine vanishing.}
    Every \(\pi_n(A)\) is naturally a module over \(\pi_0(A)\). If the
    latter is the zero ring, all homotopy groups of \(A\) vanish. By
    \Ref{prop:postnikov-completeness-detection}, \(A\simeq0\), and its
    representable functor of points is initial.

    \item \textbf{Global vanishing.}
    If an affine atlas of \(X\) has only empty members, its coproduct is the
    initial object. The atlas map is an effective epimorphism. In an
    \(\infty\)-topos an effective epimorphism from the initial object has
    initial target, so \(X\simeq\varnothing\).
\end{enumerate}
\end{proof}

\begin{theorem}[Complement of a closed immersion]
\label{thm:complement-closed-immersion}
Let
\[
    i:Z\longrightarrow X
\]
be a closed immersion of spectral schemes. There exists an open immersion
\[
    j:X\setminus Z\hookrightarrow X
\]
characterized by the universal property that a morphism \(T\to X\) factors
through \(X\setminus Z\) if and only if
\[
    T\times_XZ\simeq\varnothing.
\]
\end{theorem}

\begin{proof}
\leavevmode
\begin{enumerate}[label=\textup{(\arabic*)}]
    \item \textbf{Construct the complement over an affine context.}
    Suppose first that \(X=\Spec(A)\) and
    \(Z=\Spec(B)\), with \(\pi_0(A)\to\pi_0(B)\) surjective. Put
    \[
        I:=\ker\bigl(\pi_0(A)\to\pi_0(B)\bigr).
    \]
    In the spectral Zariski \(\infty\)-topos, form the image
    \[
    \begin{tikzcd}[column sep=huge]
        \displaystyle\coprod_{f\in I}D_A^{\mathrm{sp}}(f)
            \arrow[r, two heads]
        &
        U
            \arrow[r, hook]
        &
        X.
    \end{tikzcd}
    \]
    By \Ref{prop:unions-open-substacks}, this image is the open union of the
    displayed principal spectral opens and is stable under arbitrary base
    change. Since \(X\) is a spectral scheme, \(U\) is a spectral scheme.

    \item \textbf{The constructed open is disjoint from the closed part.}
    For \(f\in I\), the affine intersection
    \[
        D_A^{\mathrm{sp}}(f)\times_XZ
        \simeq
        \Spec\bigl(B\otimes_AA[f^{-1}]\bigr)
    \]
    has degree-zero ring
    \[
        \pi_0(B)\otimes_{\pi_0(A)}\pi_0(A)[f^{-1}].
    \]
    The image of \(f\) is zero in \(\pi_0(B)\) and invertible in the
    localization, so this ring is zero. By
    \Ref{lem:empty-spectral-scheme-classical-detection}, the intersection is
    empty. Since the principal opens cover \(U\), mapping descent gives
    \[
        U\times_XZ\simeq\varnothing.
    \]

    \item \textbf{Translate emptiness to degree zero.}
    Let \(T=\Spec(C)\to X\) satisfy
    \(T\times_XZ\simeq\varnothing\). Then
    \[
        0
        =
        \pi_0(C\otimes_AB)
        \simeq
        \pi_0(C)\otimes_{\pi_0(A)}\pi_0(B)
        \simeq
        \pi_0(C)/I\pi_0(C).
    \]

    \item \textbf{Generate the unit ideal.}
    Thus the image of \(I\) generates the unit ideal in \(\pi_0(C)\).

    \item \textbf{Choose finitely many generators.}
    Choose finitely many \(f_1,\ldots,f_r\in I\) whose images generate the
    unit ideal.

    \item \textbf{Construct the principal cover.}
    The principal opens
    \[
        \{D_C^{\mathrm{sp}}(f_j)\to T\}_{j=1}^r
    \]
    form a Zariski cover,

    \item \textbf{Produce the local factorizations.}
    Each composite to \(X\) factors through the
    corresponding \(D_A^{\mathrm{sp}}(f_j)\subseteq U\).

    \item \textbf{Check agreement on intersections.}
    Since \(U\to X\)
    is a monomorphism, these local factorizations agree on all intersections.

    \item \textbf{Glue the local factorizations.}
    Mapping descent along the effective epimorphic cover glues them to a
    unique factorization \(T\to U\to X\).

    \item \textbf{Verify the converse.}
    The converse follows from
    Step \textup{(2)}.

    \item \textbf{Choose an affine atlas.}
    For general \(X\), choose an affine open atlas
    \(\{X_a\to X\}\).

    \item \textbf{Construct the local complements.}
    The pullbacks \(Z_a:=Z\times_XX_a\) are affine
    closed immersions, and the preceding affine steps construct their
    open complements \(U_a\subseteq X_a\).

    \item \textbf{Compare the local complements on overlaps.}
    The affine universal property
    is preserved by arbitrary base change, so the \(U_a\) agree canonically
    on all iterated intersections.

    \item \textbf{Glue the global open subscheme.}
    They therefore glue to an open spectral
    subscheme \(U\hookrightarrow X\).

    \item \textbf{Verify the global universal property.}
    The same mapping-descent argument,
    now on an affine cover of an arbitrary \(T\to X\), shows that
    \[
        T\to U
        \quad\Longleftrightarrow\quad
        T\times_XZ\simeq\varnothing.
    \]

    \item \textbf{Define the complement.}
    Define \(X\setminus Z:=U\).

    \item \textbf{Record uniqueness.}
    The universal property gives uniqueness.
\end{enumerate}
\end{proof}

\begin{proposition}[Base change of complements]
\label{prop:complement-base-change}
Let \(Z\to X\) be closed and let \(Y\to X\) be any morphism. Then there is
a canonical equivalence
\[
    Y\times_X(X\setminus Z)
    \simeq
    Y\setminus(Y\times_XZ).
\]
\end{proposition}

\[
\begin{tikzcd}[column sep=huge, row sep=large]
    Y\setminus(Y\times_XZ)
        \arrow[r]
        \arrow[d, hook]
    &
    X\setminus Z
        \arrow[d, hook]
    \\
    Y
        \arrow[r]
    &
    X.
\end{tikzcd}
\]


\begin{proof}
\leavevmode
\begin{enumerate}[label=\textup{(\arabic*)}]
    \item \textbf{Translate maps into the pullback.}
    For every \(T\to Y\), a map
    \(T\to Y\times_X(X\setminus Z)\) over \(Y\) is equivalently a map
    \(T\to Y\) whose composite with \(X\) factors through
    \(X\setminus Z\).

    \item \textbf{Apply the complement criterion.}
    By \Ref{thm:complement-closed-immersion}, this is
    equivalent to
    \[
        T\times_XZ\simeq\varnothing.
    \]

    \item \textbf{Compare the two intersections.}
    Since
    \[
        T\times_Y(Y\times_XZ)
        \simeq
        T\times_XZ,
    \]
    the same condition represents maps into the complement of
    \(Y\times_XZ\) in \(Y\).

    \item \textbf{Apply Yoneda.}
    Yoneda gives the asserted equivalence.
\end{enumerate}
\end{proof}

\subsection{Pushouts along closed immersions}

\begin{theorem}[Closed pushouts]
\label{thm:closed-pushouts-spectral-schemes}
Suppose
\[
    Z
    \longrightarrow
    X
    \qquad\text{and}\qquad
    Z
    \longrightarrow
    Y
\]
are closed immersions of spectral schemes. Then the pushout
\[
    X\coprod_ZY
\]
exists in spectral schemes. The canonical maps
\[
    X
    \longrightarrow
    X\coprod_ZY
    \qquad\text{and}\qquad
    Y
    \longrightarrow
    X\coprod_ZY
\]
are closed immersions. The construction is compatible with arbitrary
base change. The resulting cocartesian square is displayed as
\[
\begin{tikzcd}[column sep=large, row sep=large]
    Z
        \arrow[r]
        \arrow[d]
    &
    X
        \arrow[d]
    \\
    Y
        \arrow[r]
    &
    X\coprod_ZY
        \arrow[ul, phantom, very near start, "\ulcorner"]
\end{tikzcd}
\]
\end{theorem}

\begin{proof}
\leavevmode
\begin{enumerate}[label=\textup{(\arabic*)}]
    \item \textbf{Affine pullback algebra.}
    First suppose
    \[
        X=\Spec(A),
        \qquad
        Y=\Spec(B),
        \qquad
        Z=\Spec(C),
    \]
    with \(A\to C\leftarrow B\) surjective on \(\pi_0\). Put
    \[
        D:=A\times_CB
    \]
    in \(\CAlg(\Sp)\). Since the forgetful functor creates limits, the
    underlying spectra fit into
    \[
    \begin{tikzcd}[column sep=large]
        D
            \arrow[r]
        &
        A\oplus B
            \arrow[r, "{\mathrm{diff}}"]
        &
        C.
    \end{tikzcd}
    \]

    \item \textbf{Compute \(\pi_{-1}(D)\).}
    The long exact homotopy sequence gives
    \[
        \pi_{-1}(D)
        \simeq
        \operatorname{coker}
        \bigl(\pi_0(A)\oplus\pi_0(B)\to\pi_0(C)\bigr)=0,
    \]
    because either structural map onto \(\pi_0(C)\) is surjective.

    \item \textbf{Deduce connectivity of \(D\).}
    Lower
    negative homotopy groups vanish automatically, so \(D\) is connective.

    \item \textbf{Compute \(\pi_0(D)\).}
    The same sequence identifies
    \[
        \pi_0(D)
        \simeq
        \pi_0(A)\times_{\pi_0(C)}\pi_0(B),
    \]
    whose projections to \(\pi_0(A)\) and \(\pi_0(B)\) are surjective.

    \item \textbf{Deduce closedness of the affine maps.}
    Hence
    \[
    \begin{tikzcd}[column sep=large, row sep=large]
        Z \arrow[r] \arrow[d]
            & X \arrow[d]
        \\
        Y \arrow[r]
            & \Spec(D)
    \end{tikzcd}
    \]
    is a pushout square of affine spectral schemes with closed horizontal and
    vertical maps.

    \item \textbf{Affine base change.}
    For any map \(D\to D'\), exactness of extension of scalars sends the
    fibre sequence above to
    \[
        D'
        \longrightarrow
        (A\otimes_DD')\oplus(B\otimes_DD')
        \longrightarrow
        C\otimes_DD'.
    \]
    The forgetful functor from commutative ring spectra creates limits, so the
    induced algebra map
    \[
        D'
        \longrightarrow
        (A\otimes_DD')
        \times_{C\otimes_DD'}
        (B\otimes_DD')
    \]
    is an equivalence. Thus the affine pushout is preserved by arbitrary base
    change.

    \item \textbf{Pass to the structured spectral model.}
    By \Ref{thm:spectral-scheme-model-comparison}, the functor-of-points
    spectral schemes used here identify with connective \(0\)-localic
    structured spectral schemes.

    \item \textbf{Match the affine closed-immersion condition.}
    On affines, the closed-immersion condition
    in the structured closed-gluing theorem is exactly surjectivity on
    \(\pi_0\), agreeing with
    \Ref{def:affine-closed-immersion-spectral}; see
    \cite[Theorem~4.4 and Remark~4.5]{LurieDAGIX}.

    \item \textbf{Apply the closed-gluing theorem.}
    Hence the
    spectral-scheme closed-gluing theorem applies and gives a pushout
    \[
    \begin{tikzcd}[column sep=large, row sep=large]
        Z \arrow[r] \arrow[d]
            & X \arrow[d]
        \\
        Y \arrow[r]
            & P
    \end{tikzcd}
    \]
    with the two canonical maps into \(P\) closed; see
    \cite[Theorem~6.1 and Remark~6.2]{LurieDAGIX}.

    \item \textbf{Identify the pushout.}
    Hence
    \(P\simeq X\coprod_ZY\).

    \item \textbf{Global base change.}
    The closed-gluing base-change theorem identifies, for every \(T\to P\),
    the square
    \[
    \begin{tikzcd}[column sep=large, row sep=large]
        T\times_PZ \arrow[r] \arrow[d]
            & T\times_PX \arrow[d]
        \\
        T\times_PY \arrow[r]
            & T
    \end{tikzcd}
    \]
    as a pushout. Its proof is affine-local and applies to spectral schemes;
    see \cite[Proposition~9.2]{LurieDAGIX}. This proves compatibility with
    arbitrary base change.
\end{enumerate}
\end{proof}
\begin{remark}[Closed gluing carries no presentation data]
\label{rem:closed-gluing-no-presentation-data}
The construction of the closed pushout uses several proof witnesses: the
affine pullback algebra \(A\times_CB\), affine charts used to verify
connectivity, and the structured-space comparison used for the global gluing
theorem. None is retained as structure on the result. The spectral scheme
\(X\coprod_ZY\) is determined by its pushout universal property, with no
chosen affine presentation, gluing presentation, or admissibility datum.
\end{remark}

Closed immersions therefore admit the two operations that will remain active
below: intrinsic open complements and base-change-stable pushouts. Together
with affine comprehension, this gives a geometric calculus in which the
linear constructions of the next section can be represented and manipulated.

\section{Linear comprehension and distinguished affine contexts}
\label{sec:vector-bundles-brave-new-affine-spaces}

Affine comprehension represents connective commutative algebra in a
coefficient fibre. Applying it to free commutative algebras produces the
linear geometric objects needed here. There are two related representation
problems: relative linear spaces represent linear functionals, whereas section
stacks represent sections. Finite locally free coefficients are the sector in
which duality identifies the latter with affine geometry. The constructions
below organize these two variances and then recover affine spaces, the brave
new affine line, and the spectral multiplicative group as distinguished
examples.

\subsection{Affine and global local freeness}


\begin{definition}[Locally free module]
\label{def:locally-free-affine-module}
Let $A$ be connective. An $A$-module $M$ is \textbf{locally free of rank
$n$} if there exists a Zariski covering family
\[
    \{A\to A_i\}_{i\in I}
\]
as defined in \Ref{def:affine-zariski-cover}, such that
\[
    M\otimes_AA_i
    \simeq
    A_i^{\oplus n}
\]
for every $i$.
\end{definition}


\begin{definition}[Locally free quasi-coherent module]
\label{def:locally-free-qcoh}
Let $X$ be a spectral scheme. A quasi-coherent module
$\mathcal E\in\QCoh(X)$ is \textbf{locally free of rank $n$} if there is a
Zariski open cover $\{U_i\to X\}$ such that
\[
    \mathcal E|_{U_i}
    \simeq
    \mathcal O_{U_i}^{\oplus n}
\]
for every $i$.
\end{definition}

A \textbf{vector bundle of rank $n$} on $X$ means a locally free
quasi-coherent module of rank $n$.

\begin{proposition}[Locally free modules are finite projective]
\label{prop:locally-free-finite-projective}
Let
\[
    \mathcal E
    \in
    \QCoh(X)
\]
be a quasi-coherent module. Suppose there is a Zariski open cover \(\{U_i\to X\}\) and nonnegative
integers \(n_i\) such that
\[
    \mathcal E|_{U_i}
    \simeq
    \mathcal O_{U_i}^{\oplus n_i}
\]
for every \(i\). Then \(\mathcal E\) is dualizable, connective, and flat. On
every affine open it corresponds to a finite projective module.
\end{proposition}

\begin{proof}
\leavevmode
\begin{enumerate}[label=\textup{(\arabic*)}]
    \item \textbf{Choose an affine trivializing atlas.}
    Refine the given Zariski cover by affine open subschemes and choose an
    affine Zariski atlas \(\{U_i\to X\}\) with
    \[
        E_i:=\mathcal E|_{U_i}
        \simeq
        \mathcal O_{U_i}^{\oplus n_i}.
    \]

    \item \textbf{Dualize the transition equivalences.}
    Each \(E_i\) is dualizable. On an overlap \(U_{ij}\), the descent
    equivalence
    \[
        g_{ij}:E_i|_{U_{ij}}\xrightarrow{\simeq}E_j|_{U_{ij}}
    \]
    induces functorially under dualization an equivalence
    \[
        (g_{ij}^{-1})^\vee:
        E_i^\vee|_{U_{ij}}
        \xrightarrow{\simeq}
        E_j^\vee|_{U_{ij}}.
    \]

    \item \textbf{Transfer the cocycle coherences.}
    Dualization is a functor on the maximal \(\infty\)-groupoid of
    dualizable objects, so the cocycle and all higher compatibilities are the
    images of those already carried by the descent object \(\mathcal E\).

    \item \textbf{Glue the local duals.}
    By \Ref{thm:qcoh-atlas-computation}, the \(E_i^\vee\) therefore glue to
    \(\mathcal E^\vee\in\QCoh(X)\).

    \item \textbf{Glue evaluation and coevaluation.}
    The local evaluation and coevaluation maps are natural under the
    transition equivalences and hence glue to
    \[
        \mathcal E^\vee\otimes\mathcal E\longrightarrow\mathcal O_X,
        \qquad
        \mathcal O_X\longrightarrow\mathcal E\otimes\mathcal E^\vee.
    \]

    \item \textbf{Check the triangle identities.}
    Their triangle identities hold after restriction to the covering family,
    and affine-open restrictions jointly detect equivalences and homotopies by
    descent.

    \item \textbf{Conclude global dualizability.}
    Thus \(\mathcal E\) is globally dualizable.

    \item \textbf{Connectivity.}
    Connectivity is detected on affine opens by
    \Ref{thm:qcoh-spectral-scheme-t-structure}. Every local trivialization is
    connective, so \(\mathcal E\) is connective.

    \item \textbf{Choose an affine open.}
    Let \(U=\Spec(A)\subseteq X\) be affine and let \(M\in\Mod_A\) represent
    \(\mathcal E|_U\).

    \item \textbf{Pull back the trivializing cover.}
    Pull back the trivializing cover to \(U\) and refine it
    by affine opens.

    \item \textbf{Recognize the affine Zariski cover.}
    The resulting affine-open coproduct is effective
    epimorphic, hence is an affine spectral Zariski cover by
    \Ref{prop:affine-open-cover-effective-epi}.

    \item \textbf{Record the local finite-free forms.}
    Thus there are affine opens
    \(\{\Spec(A_j)\to\Spec(A)\}\) and nonnegative integers \(r_j\) with
    \[
        M\otimes_AA_j\simeq A_j^{\oplus r_j}.
    \]

    \item \textbf{Apply locality of flatness.}
    The local modules are flat, so \Ref{prop:flat-module-zariski-local} makes
    \(M\) flat.

    \item \textbf{Conclude global flatness.}
    Thus \(\mathcal E\) is flat on every affine open.

    \item \textbf{Finite projectivity on affine opens.}
    On an affine open \(U=\Spec(A)\), dualizability gives perfectness of the
    corresponding module by \Ref{prop:perfect-equals-dualizable}. Together
    with connectivity and flatness,
    \Ref{prop:finite-projective-characterization} implies that the module is
    finite projective.
\end{enumerate}
\end{proof}

\begin{theorem}[Finite projective modules and vector bundles]
\label{thm:finite-projective-locally-free}
Let \(A\) be a connective commutative ring spectrum and let \(P\in\Mod_A\).
The following are equivalent.
\begin{enumerate}[label=(\roman*)]
    \item \(P\) is finite projective.
    \item \(P\) is Zariski-locally finite free of finite rank.
    \item \(P\) is perfect, connective, and flat.
\end{enumerate}
Consequently, a quasi-coherent module on a spectral scheme is
Zariski-locally free of locally constant finite rank exactly when its
restriction to every affine open is finite projective. The local ranks then
define a locally constant function
\[
    |X|
    \longrightarrow
    \mathbb N.
\]
It is a vector bundle of rank \(n\) in the sense of
\Ref{def:locally-free-qcoh} exactly when this rank function is identically
equal to \(n\).
\end{theorem}

\begin{proof}
\leavevmode
\begin{enumerate}[label=\textup{(\arabic*)}]
    \item \textbf{Algebraic characterization.}
    The equivalence \textup{(i)}\(\Leftrightarrow\)\textup{(iii)} is
    \Ref{prop:finite-projective-characterization}. The implication
    \textup{(ii)}\(\Rightarrow\)\textup{(iii)} follows from
    \Ref{prop:locally-free-finite-projective}.

    \item \textbf{Degree-zero local triviality.}
    Assume \(P\) finite projective. Then \(P\) is flat, and
    \[
        \pi_n(P)
        \simeq
        \pi_n(A)\otimes_{\pi_0(A)}\pi_0(P).
    \]
    The module \(\pi_0(P)\) is finitely generated projective over
    \(\pi_0(A)\). Hence there are finitely many elements
    \(f_i\in\pi_0(A)\) generating the unit ideal such that
    \[
        \pi_0(P)[f_i^{-1}]
        \simeq
        \pi_0(A)[f_i^{-1}]^{\oplus r_i}.
    \]

    \item \textbf{Lift the local bases.}
    Put
    \[
        A_i:=A[f_i^{-1}],
        \qquad
        P_i:=P\otimes_AA_i.
    \]
    A chosen basis of \(\pi_0(P_i)\) defines a homomorphism
    \[
        \pi_0(A_i)^{\oplus r_i}
        \xrightarrow{\simeq}
        \pi_0(P_i).
    \]
    Since maps from a finite free \(A_i\)-module are finite tuples of
    elements, this basis lifts to an \(A_i\)-linear morphism
    \[
        u_i:A_i^{\oplus r_i}\longrightarrow P_i.
    \]

    \item \textbf{Detect the lifted basis on homotopy groups.}
    Both sides are flat. Hence, for every integer \(n\),
    \[
        \pi_n(u_i)
        \simeq
        \pi_n(A_i)\otimes_{\pi_0(A_i)}\pi_0(u_i),
    \]
    which is an isomorphism. Homotopy groups detect equivalences, so
    \[
        A_i^{\oplus r_i}\xrightarrow{\simeq}P_i.
    \]
    Thus \(P\) is Zariski-locally finite free.

    \item \textbf{Globalize affine local freeness.}
    Apply the affine equivalence on an affine Zariski atlas of a spectral
    scheme. If a quasi-coherent module restricts to a finite projective
    module on every affine open, then each affine restriction is
    Zariski-locally finite free by the preceding steps.

    \item \textbf{Glue the local trivializations.}
    These local
    trivializations together give a Zariski cover on which the
    quasi-coherent module is finite free.

    \item \textbf{Construct the locally constant rank function.}
    The ranks agree after restriction to overlaps, since the rank of a
    finitely generated projective module is preserved by base change.
    They therefore determine a locally constant function
    \[
        |X|\longrightarrow\mathbb N.
    \]

    \item \textbf{Prove the converse local criterion.}
    Conversely, a Zariski-locally finite free quasi-coherent module restricts
    on every affine open to a locally free module, hence to a finite
    projective module by
    \Ref{prop:locally-free-finite-projective}.

    \item \textbf{Characterize fixed rank.}
    Finally, the module is locally free of the fixed rank \(n\) precisely
    when the resulting locally constant rank function is identically \(n\).
    This is exactly the notion of a vector bundle of rank \(n\) from
    \Ref{def:locally-free-qcoh}.
\end{enumerate}
\end{proof}

Local freeness is preserved by substitution; a trivializing cover pulls back
through
\[
\begin{tikzcd}[column sep=large, row sep=large]
    Y\times_XU_i \arrow[r] \arrow[d]
        & U_i \arrow[d] \\
    Y \arrow[r, "f"'] & X.
\end{tikzcd}
\]

\begin{proposition}[Base change of locally free modules]
\label{prop:base-change-locally-free}
If $f:Y\to X$ is any morphism and $\mathcal E$ is locally free of rank
$n$, then $f^*\mathcal E$ is locally free of rank $n$.
\end{proposition}

\begin{proof}
\leavevmode
\begin{enumerate}[label=\textup{(\arabic*)}]
    \item \textbf{Pull back the trivializing cover.}
    Pull back a trivializing Zariski cover of $X$.

    \item \textbf{Preserve the local data.}
    Open covers and finite free
    modules are preserved under base change.
\end{enumerate}
\end{proof}

\subsection{The relative free commutative algebra}

\begin{proposition}[Relative free commutative algebra]
\label{prop:relative-free-commutative-algebra}
For every spectral scheme $X$, the forgetful functor
\[
    \operatorname{oblv}^{\mathrm{alg}}_{X}
    \colon
    \operatorname{QCAlg}(X)
    \longrightarrow
    \QCoh(X)
\]
admits a left adjoint
\[
    \operatorname{Sym}_{\mathcal O_X}
    \colon
    \QCoh(X)
    \longrightarrow
    \operatorname{QCAlg}(X).
\]
The construction is compatible with pullback: for every $f:Y\to X$ there
is a canonical equivalence
\[
    f^*
    \operatorname{Sym}_{\mathcal O_X}(\mathcal E)
    \simeq
    \operatorname{Sym}_{\mathcal O_Y}(f^*\mathcal E).
\]
\end{proposition}

\[
\begin{tikzcd}[column sep=huge, row sep=large]
    \QCoh(X)
        \arrow[r, "{\operatorname{Sym}_{\mathcal O_X}}"]
        \arrow[d, "{f^*}"']
    &
    \operatorname{QCAlg}(X)
        \arrow[d, "{f^*}"]
    \\
    \QCoh(Y)
        \arrow[r, "{\operatorname{Sym}_{\mathcal O_Y}}"']
    &
    \operatorname{QCAlg}(Y).
\end{tikzcd}
\]


\begin{proof}
\leavevmode
\begin{enumerate}[label=\textup{(\arabic*)}]
    \item \textbf{Construct the free-algebra adjunction.}
    The category $\QCoh(X)$ is presentably symmetric monoidal, so the
    operadic free-algebra theorem supplies the adjunction.

    \item \textbf{Commute the construction with pullback.}
    Since pullback is
    colimit-preserving and strong symmetric monoidal, it commutes with the
    operadic free commutative algebra construction.
\end{enumerate}
\end{proof}

\begin{proposition}[Connectivity of the relative free algebra]
\label{prop:relative-free-algebra-connective}
If
\[
    \mathcal E
    \in
    \QCoh(X)_{\geq0},
\]
then
\[
    \operatorname{Sym}_{\mathcal O_X}(\mathcal E)
    \in
    \operatorname{QCAlg}(X)^{\mathrm{cn}}.
\]
\end{proposition}

\begin{proof}
\leavevmode
\begin{enumerate}[label=\textup{(\arabic*)}]
    \item \textbf{Reduce to an affine open.}
    The assertion is affine-local. Let $U=\Spec(A)\subseteq X$, and let
    $M\in(\Mod_{A})_{\geq0}$ correspond to $\mathcal E|_U$.

    \item \textbf{Write the underlying free algebra.}
    The underlying
    $A$-module of the free commutative $A$-algebra is
    \[
        \bigoplus_{k\geq0}
        \left(
            M^{\otimes_Ak}
        \right)_{h\Sigma_k}.
    \]

    \item \textbf{Check tensor-power connectivity.}
    Every tensor power $M^{\otimes_Ak}$ is connective.

    \item \textbf{Check colimit connectivity.}
    Homotopy orbits and
    small coproducts preserve connective objects, so the displayed
    $A$-module is connective.

    \item \textbf{Detect connectivity globally.}
    Connectivity of a quasi-coherent algebra is
    detected on affine opens, and the result follows.
\end{enumerate}
\end{proof}

\subsection{Total spaces and affine spaces}

\begin{definition}[Relative linear space]
\label{def:linear-spectral-space}
For a connective quasi-coherent module
$\mathcal E\in\QCoh(X)_{\geq0}$, define its \textbf{relative linear space} by
\[
    \mathbb L_X(\mathcal E)
    :=
    \Spec_X
    \operatorname{Sym}_{\mathcal O_X}(\mathcal E).
\]
\end{definition}

\[
\begin{tikzcd}[column sep=huge, row sep=large]
    \QCoh(X)_{\geq0}^{\op}
        \arrow[r, "{(\operatorname{Sym}_{\mathcal O_X})^{\op}}"]
        \arrow[dr, "{\mathbb L_X}"']
    &
    \operatorname{QCAlg}(X)^{\mathrm{cn},\op}
        \arrow[d, "{\Spec_X}"]
    \\
    &
    \bigl(\operatorname{Sch}^{\mathrm{sp}}_{/X}\bigr)_{\mathrm{aff}}.
\end{tikzcd}
\]


\begin{proposition}[Functor of points of a relative linear space]
\label{prop:linear-space-functor-points}
For every morphism $f:T\to X$, there is a natural equivalence
\[
    \Map_X
    \bigl(
        T,
        \mathbb L_X(\mathcal E)
    \bigr)
    \simeq
    \Map_{\QCoh(T)}
    \bigl(
        f^*\mathcal E,
        \mathcal O_T
    \bigr).
\]
The variance is summarized by
\[
\begin{tikzcd}[column sep=huge, row sep=large]
    T
        \arrow[rr, "\ell"]
        \arrow[dr, "f"']
    &&
    \mathbb L_X(\mathcal E)
        \arrow[dl]
    \\
    & X &
\end{tikzcd}
\qquad\longleftrightarrow\qquad
f^*\mathcal E\xrightarrow{\ \ell^\sharp\ }\mathcal O_T.
\]
\end{proposition}

\begin{proof}
\leavevmode
\begin{enumerate}[label=\textup{(\arabic*)}]
    \item \textbf{Apply relative spectrum.}
    Apply the relative spectrum universal property of
    \Ref{thm:relative-spectrum-construction}.

    \item \textbf{Apply the free-algebra adjunction.}
    Apply the free-algebra adjunction of
    \Ref{prop:relative-free-commutative-algebra}.
\end{enumerate}
\end{proof}

\begin{construction}[Section stack of a quasi-coherent coefficient]
\label{con:section-stack-qcoh}
Let \(X\) be a spectral scheme and let \(\mathcal E\in\QCoh(X)\). Define a
spectral prestack over \(X\), denoted
\[
    \operatorname{Sect}_X(\mathcal E)
    \longrightarrow
    X,
\]
by the following functor of affine points. For an affine spectral scheme
\(T\), the map to \(X(T)\) has fibre over \(f:T\to X\)
\[
    \Map_{\QCoh(T)}
    \bigl(\mathcal O_T,f^*\mathcal E\bigr).
\]
Equivalently, its total space is the Grothendieck construction of the functor
\[
    X(T)
    \longrightarrow
    \mathcal S,
    \qquad
    f\longmapsto
    \Map_{\QCoh(T)}
    \bigl(\mathcal O_T,f^*\mathcal E\bigr).
\]
\end{construction}

\begin{proposition}[Descent and base change for section stacks]
\label{prop:section-stack-descent-base-change}
For every \(\mathcal E\in\QCoh(X)\), the prestack
\(\operatorname{Sect}_X(\mathcal E)\) is a spectral Zariski stack over
\(X\). For every morphism \(g:Y\to X\), there is a canonical equivalence
over \(Y\)
\[
    \operatorname{Sect}_X(\mathcal E)\times_XY
    \simeq
    \operatorname{Sect}_Y(g^*\mathcal E).
\]
Thus the square
\[
\begin{tikzcd}[column sep=large, row sep=large]
    \operatorname{Sect}_Y(g^*\mathcal E)
        \arrow[r]
        \arrow[d]
    &
    \operatorname{Sect}_X(\mathcal E)
        \arrow[d]
    \\
    Y
        \arrow[r, "g"']
    &
    X
\end{tikzcd}
\]
is Cartesian.
\end{proposition}

\begin{proof}
\leavevmode
\begin{enumerate}[label=\textup{(\arabic*)}]
    \item \textbf{Specify descent data.}
    Let \(T_\bullet\to T\) be the \v{C}ech nerve of an affine spectral
    Zariski cover. An object of
    \[
        \operatorname*{Tot}_{[n]\in\Delta}
        \operatorname{Sect}_X(\mathcal E)(T_n)
    \]
    consists of a coherent family of maps \(f_n:T_n\to X\) and compatible
    sections
    \[
        s_n:\mathcal O_{T_n}\longrightarrow f_n^*\mathcal E.
    \]

    \item \textbf{Glue the maps to \(X\).}
    Since \(X\) is a spectral Zariski sheaf, the \(f_n\) glue uniquely to a
    map \(f:T\to X\).

    \item \textbf{Apply quasi-coherent descent.}
    By \Ref{thm:qcoh-atlas-computation},
    \[
        \QCoh(T)
        \simeq
        \operatorname*{Tot}_{[n]\in\Delta}\QCoh(T_n),
    \]
    and under this equivalence \(\mathcal O_T\) and \(f^*\mathcal E\)
    correspond to their restrictions.

    \item \textbf{Identify the mapping-space totalization.}
    Therefore
    \[
        \Map_{\QCoh(T)}(\mathcal O_T,f^*\mathcal E)
        \simeq
        \operatorname*{Tot}_{[n]\in\Delta}
        \Map_{\QCoh(T_n)}
        (\mathcal O_{T_n},f_n^*\mathcal E).
    \]

    \item \textbf{Glue the local sections.}
    The \(s_n\) therefore glue uniquely, together with all higher homotopies, to a
    section \(s:\mathcal O_T\to f^*\mathcal E\).

    \item \textbf{Conclude descent of the total prestack.}
    This constructs the
    inverse to the restriction map and proves descent for the total prestack.

    \item \textbf{Base change.}
    A \(T\)-point of the left-hand side over a map \(h:T\to Y\) is a section
    \[
        \mathcal O_T
        \longrightarrow
        (g\circ h)^*\mathcal E
        \simeq
        h^*g^*\mathcal E.
    \]
    This is exactly a \(T\)-point of
    \(\operatorname{Sect}_Y(g^*\mathcal E)\) over \(h\). Naturality in
    \(T\) and Yoneda give the claimed equivalence.
\end{enumerate}
\end{proof}

\begin{definition}[Total space representing sections]
\label{def:vector-bundle-dual-convention}
Let \(\mathcal E\) be a finite locally free quasi-coherent module on \(X\).
Define its \textbf{section total space} by
\[
    \mathbf V_X(\mathcal E)
    :=
    \mathbb L_X(\mathcal E^\vee).
\]
\end{definition}

\begin{proposition}[Affine representability and base change for section total spaces]
\label{prop:vector-bundle-total-space-sections}
Let \(\mathcal E\) be finite locally free on \(X\). Then the section stack
of \Ref{con:section-stack-qcoh} is represented by the affine spectral scheme
\(\mathbf V_X(\mathcal E)\):
\[
    \operatorname{Sect}_X(\mathcal E)
    \simeq
    \mathbf V_X(\mathcal E)
\]
over \(X\). Equivalently, for every morphism \(f:T\to X\), there is a
natural equivalence
\[
    \Map_X\bigl(T,\mathbf V_X(\mathcal E)\bigr)
    \simeq
    \Map_{\QCoh(T)}
    \bigl(\mathcal O_T,f^*\mathcal E\bigr).
\]
For every morphism \(g:Y\to X\), there is a canonical equivalence
\[
    \mathbf V_X(\mathcal E)\times_XY
    \simeq
    \mathbf V_Y(g^*\mathcal E).
\]
For the finite free module \(\mathcal O_X^{\oplus n}\), the canonical
self-duality identifies
\[
    \mathbf V_X(\mathcal O_X^{\oplus n})
    \simeq
    \mathbb L_X(\mathcal O_X^{\oplus n}).
\]
\end{proposition}

\[
\begin{tikzcd}[column sep=large, row sep=large]
    \mathbf V_Y(g^*\mathcal E)
        \arrow[r]
        \arrow[d]
    &
    \mathbf V_X(\mathcal E)
        \arrow[d]
    \\
    Y
        \arrow[r, "g"']
    &
    X.
\end{tikzcd}
\]


\begin{proof}
\leavevmode
\begin{enumerate}[label=\textup{(\arabic*)}]
    \item \textbf{Represent sections by duality.}
    Strong symmetric monoidal pullback preserves duals, so
    \[
        f^*(\mathcal E^\vee)
        \simeq
        (f^*\mathcal E)^\vee.
    \]
    By \Ref{prop:linear-space-functor-points} and the tensor--Hom adjunction
    for the dualizable module \(f^*\mathcal E\),
    \[
    \begin{aligned}
        \Map_X\bigl(T,\mathbf V_X(\mathcal E)\bigr)
        &\simeq
        \Map_{\QCoh(T)}
        \bigl((f^*\mathcal E)^\vee,\mathcal O_T\bigr)
        \\
        &\simeq
        \Map_{\QCoh(T)}
        \bigl(\mathcal O_T,f^*\mathcal E\bigr).
    \end{aligned}
    \]

    \item \textbf{Base change.}
    Relative linear spaces commute with base change by
    \Ref{prop:relative-spectrum-base-change} and
    \Ref{prop:relative-free-commutative-algebra}. Together with preservation
    of duals by pullback this gives
    \[
        \mathbf V_X(\mathcal E)\times_XY
        \simeq
        \mathbb L_Y((g^*\mathcal E)^\vee)
        =
        \mathbf V_Y(g^*\mathcal E).
    \]

    \item \textbf{Finite free self-duality.}
    The direct-sum decomposition gives the canonical duality
    \[
        (\mathcal O_X^{\oplus n})^\vee
        \simeq
        \mathcal O_X^{\oplus n},
    \]
    yielding the final equivalence.
\end{enumerate}
\end{proof}

\begin{remark}[Linear comprehension]
\label{rem:linear-comprehension}
The two linear representation problems now fit into a single picture. For
any coefficient \(\mathcal E\in\QCoh(X)\),
\Ref{con:section-stack-qcoh} constructs the stack representing sections
\[
    \mathcal O_T\longrightarrow f^*\mathcal E.
\]
When \(\mathcal E\) is finite locally free,
\Ref{prop:vector-bundle-total-space-sections} identifies this stack with the
affine context \(\mathbf V_X(\mathcal E)\). Thus section comprehension is
available as a stack for arbitrary quasi-coherent coefficients and becomes
affine representable in the finite locally free sector. In that sector,
duality gives the comparison
\[
    \mathbb L_X(\mathcal E)
    \simeq
    \mathbf V_X(\mathcal E^\vee),
\]
so linear-functional and section representations differ exactly by duality.
\end{remark}

\begin{remark}[Coordinate-module and section conventions]
\label{rem:linear-coordinate-section-conventions}
The relative linear space and the section total space represent opposite
variances, and duality is required to identify them. A map
\[
    T\longrightarrow\mathbb L_X(\mathcal E)
\]
over \(X\) represents a linear functional
\[
    f^*\mathcal E\longrightarrow\mathcal O_T,
\]
whereas, for finite locally free \(\mathcal F\), a map
\[
    T\longrightarrow\mathbf V_X(\mathcal F)
\]
represents a section
\[
    \mathcal O_T\longrightarrow f^*\mathcal F.
\]
For finite locally free \(\mathcal E\), duality relates the two conventions
by
\[
    \mathbb L_X(\mathcal E)
    \simeq
    \mathbf V_X(\mathcal E^\vee).
\]
The distinction is therefore one of variance, not an implicit choice of a
self-duality \(\mathcal E\simeq\mathcal E^\vee\).
\end{remark}

\begin{definition}[Brave new affine space]
\label{def:brave-new-affine-space}
For an integer $n\geq0$, define
\[
    \mathbb A_X^n
    :=
    \mathbb L_X
    \bigl(
        \mathcal O_X^{\oplus n}
    \bigr)
    =
    \Spec_X
    \operatorname{Sym}_{\mathcal O_X}
    \bigl(
        \mathcal O_X^{\oplus n}
    \bigr).
\]
The case $n=1$ is the \textbf{brave new affine line}
\[
    \mathbb A_X^1
    =
    \Spec_X
    \operatorname{Sym}_{\mathcal O_X}(\mathcal O_X).
\]
\end{definition}

\begin{proposition}[Base change of affine space]
\label{prop:affine-space-base-change}
For every morphism $f:Y\to X$, there is a canonical equivalence
\[
    \mathbb A_X^n
    \times_X
    Y
    \simeq
    \mathbb A_Y^n.
\]
\end{proposition}

\[
\begin{tikzcd}[column sep=large, row sep=large]
    \mathbb A_Y^n
        \arrow[r]
        \arrow[d]
    &
    \mathbb A_X^n
        \arrow[d]
    \\
    Y
        \arrow[r, "f"']
    &
    X.
\end{tikzcd}
\]


\begin{proof}
\leavevmode
\begin{enumerate}[label=\textup{(\arabic*)}]
    \item \textbf{Commute the free algebra with pullback.}
    Use the pullback compatibility of the free algebra from
    \Ref{prop:relative-free-commutative-algebra}.

    \item \textbf{Apply relative-spectrum base change.}
    Use relative spectrum base
    change from \Ref{prop:relative-spectrum-base-change}.
\end{enumerate}
\end{proof}

\begin{construction}[Zero and unit sections]
\label{con:zero-unit-sections-affine-line}
The zero endomorphism and identity endomorphism of $\mathcal O_X$ determine
maps of quasi-coherent modules
\[
    \mathcal O_X
    \rightrightarrows
    \mathcal O_X.
\]
By the free-algebra and relative-spectrum adjunctions, they induce sections
\[
    0,1
    \colon
    X
    \rightrightarrows
    \mathbb A_X^1.
\]
\end{construction}

\begin{proposition}[Sections of brave new affine space are closed]
\label{prop:affine-space-sections-closed}
Let \(n\geq0\). Every section
\[
\begin{tikzcd}[column sep=huge, row sep=large]
    X
        \arrow[rr, "s"]
        \arrow[dr, equal]
    &&
    \mathbb A_X^n
        \arrow[dl]
    \\
    & X &
\end{tikzcd}
\]
is a closed immersion.
\end{proposition}

\begin{proof}
\leavevmode
\begin{enumerate}[label=\textup{(\arabic*)}]
    \item \textbf{Reduce to an affine context.}
    Closed immersions are local on the target, so take \(X=\Spec(A)\). The
    section \(s\) corresponds contravariantly to an \(A\)-algebra morphism
    \[
        \operatorname{Sym}_A(A^{\oplus n})
        \longrightarrow
        A.
    \]
    By the free-algebra adjunction it is determined by an \(n\)-tuple of
    elements of \(\Omega^\infty A\).

    \item \textbf{Pass to degree zero.}
    By \Ref{prop:pi-zero-relative-free-algebra}, the induced ordinary ring
    morphism is an evaluation map
    \[
        \pi_0(A)[t_1,\ldots,t_n]
        \longrightarrow
        \pi_0(A),
        \qquad
        t_i\longmapsto a_i.
    \]
    It is surjective because it restricts to the identity on the constant
    subring \(\pi_0(A)\). Therefore the corresponding affine morphism is a
    closed immersion by \Ref{def:affine-closed-immersion-spectral}.
\end{enumerate}
\end{proof}

\begin{proposition}[The zero section is closed]
\label{prop:zero-section-affine-line-closed}
The zero section
\[
    0
    \colon
    X
    \longrightarrow
    \mathbb A_X^1
\]
is a closed immersion.
\end{proposition}

\begin{proof}
\leavevmode
\begin{enumerate}[label=\textup{(\arabic*)}]
    \item \textbf{Recognize the zero map as a section.}
    The zero section is a section of \(\mathbb A_X^1\to X\).

    \item \textbf{Apply closedness of sections.}
    Apply
    \Ref{prop:affine-space-sections-closed}.
\end{enumerate}
\end{proof}

\begin{definition}[Spectral multiplicative group]
\label{def:spectral-multiplicative-group}
The \textbf{spectral multiplicative group} over $X$ is the open complement
\[
    \mathbb G_{m,X}
    :=
    \mathbb A_X^1
    \setminus
    0(X).
\]
\end{definition}

\begin{proposition}[Group structure on the spectral multiplicative group]
\label{prop:spectral-multiplicative-group-structure}
The spectral scheme
\[
    \mathbb G_{m,X}
\]
is canonically a commutative group object in spectral schemes over \(X\).
For every morphism \(T\to X\), there is a natural equivalence
\[
    \Map_X(T,\mathbb G_{m,X})
    \simeq
    \operatorname{Aut}_{\QCoh(T)}(\mathcal O_T).
\]
\end{proposition}

\begin{proof}
\leavevmode
\begin{enumerate}[label=\textup{(\arabic*)}]
    \item \textbf{Affine coordinate algebra.}
    Work first over \(X=\Spec(A)\) and write
    \[
        A\{t\}:=\operatorname{Sym}_A(A).
    \]
    By the complement construction and
    \Ref{prop:principal-spectral-open},
    \[
        \mathbb G_{m,X}
        \simeq
        \Spec\bigl(A\{t\}[t^{-1}]\bigr).
    \]

    \item \textbf{Coordinate algebra of the product.}
    There is a canonical equivalence
    \[
        A\{t_1,t_2\}[(t_1t_2)^{-1}]
        \simeq
        A\{t_1,t_2\}[t_1^{-1},t_2^{-1}],
    \]
    because an inverse of \(t_1t_2\) gives
    \[
        t_1^{-1}=t_2(t_1t_2)^{-1},
        \qquad
        t_2^{-1}=t_1(t_1t_2)^{-1}.
    \]
    The right side is the coordinate algebra of
    \(\mathbb G_{m,X}\times_X\mathbb G_{m,X}\).

    \item \textbf{Define the Hopf algebra maps.}
    The universal property of localization gives
    \[
    \begin{aligned}
        \Delta(t)&=t_1t_2,\\
        \epsilon(t)&=1,\\
        \operatorname{inv}(t)&=t^{-1}.
    \end{aligned}
    \]
    Thus
    \[
    \begin{tikzcd}[column sep=huge, row sep=large]
        &
        A\{t\}[t^{-1}]
            \arrow[dl, "\epsilon"']
            \arrow[d, "\operatorname{inv}"]
            \arrow[dr, "\Delta"]
        &
        \\
        A
        &
        A\{t\}[t^{-1}]
        &
        A\{t_1,t_2\}[t_1^{-1},t_2^{-1}].
    \end{tikzcd}
    \]

    \item \textbf{Verify the group identities.}
    The commutative group identities are the identities satisfied by the
    universal invertible element.

    \item \textbf{Pass to relative spectra.}
    Passing contravariantly to relative spectra
    produces multiplication, unit, and inversion.

    \item \textbf{Identify points of affine line with endomorphisms.}
    A map \(T\to\mathbb A_X^1\) corresponds by
    \Ref{prop:linear-space-functor-points} to an
    \(\mathcal O_T\)-linear endomorphism of \(\mathcal O_T\).

    \item \textbf{Reduce invertibility to an affine test.}
    The assertion
    that such an endomorphism is invertible is Zariski-local on \(T\), so
    take \(T=\Spec(C)\).

    \item \textbf{Identify an endomorphism with multiplication by a point.}
    Then
    \[
        \Map_{\Mod_C}(C,C)
        \simeq
        \Omega^\infty C,
    \]
    and a point \(c\in\Omega^\infty C\) acts by multiplication
    \(m_c:C\to C\).

    \item \textbf{Show an equivalence gives a unit in degree zero.}
    If \(m_c\) is an equivalence, multiplication by \([c]\) on
    \(\pi_0(C)\) is an isomorphism, so \([c]\in\pi_0(C)^\times\).

    \item \textbf{Show a unit in degree zero gives an equivalence.}
    Conversely, if \([c]\) is a unit, choose the inverse component
    \([d]\). The products \(cd\) and \(dc\) lie in the component of the
    unit, so multiplication by \(c\) and multiplication by \(d\) are inverse
    equivalences of \(C\)-modules. Thus \(m_c\) is invertible exactly when
    \([c]\) is a unit.

    \item \textbf{Characterize factorization through the localization.}
    By the universal property of localization, the corresponding map
    \(T\to\mathbb A_X^1\) factors through the principal open
    \(D_{A\{t\}}^{\mathrm{sp}}(t)=\mathbb G_{m,X}\) exactly when the image of \(t\) is a unit in
    \(\pi_0(C)\).

    \item \textbf{Identify the functor of points.}
    Therefore
    \[
        \Map_X(T,\mathbb G_{m,X})
        \simeq
        \operatorname{Aut}_{\QCoh(T)}(\mathcal O_T).
    \]

    \item \textbf{Globalize by Zariski descent.}
    Zariski descent extends this equivalence to arbitrary spectral schemes
    \(T\) over \(X\).

    \item \textbf{Globalization.}
    The affine constructions commute with restriction to affine opens and
    therefore glue over an arbitrary spectral scheme \(X\).
\end{enumerate}
\end{proof}

\begin{proposition}[Base change of the multiplicative group]
\label{prop:multiplicative-group-base-change}
For every morphism $Y\to X$, there is a canonical equivalence
\[
    \mathbb G_{m,X}
    \times_X
    Y
    \simeq
    \mathbb G_{m,Y}.
\]
\end{proposition}

\[
\begin{tikzcd}[column sep=large, row sep=large]
    \mathbb G_{m,Y}
        \arrow[r]
        \arrow[d]
    &
    \mathbb G_{m,X}
        \arrow[d]
    \\
    Y
        \arrow[r]
    &
    X.
\end{tikzcd}
\]


\begin{proof}
\leavevmode
\begin{enumerate}[label=\textup{(\arabic*)}]
    \item \textbf{Apply affine-space base change.}
    Use affine-space base change.

    \item \textbf{Apply base change of open complements.}
    Use base change of open complements from
    \Ref{prop:complement-base-change}.
\end{enumerate}
\end{proof}

\begin{proposition}[Homotopy of the brave new polynomial algebra]
\label{prop:affine-line-genuinely-spectral}
Let \(A\) be a connective commutative ring spectrum and put
\[
    A\{t\}
    :=
    \operatorname{Sym}_A(A).
\]
Then, as an underlying \(A\)-module,
\[
    A\{t\}
    \simeq
    \bigoplus_{n\geq0}
    A\otimes\Sigma^\infty_+B\Sigma_n.
\]
Consequently, if
\[
    A_m(K)
    :=
    \pi_m\bigl(A\otimes\Sigma^\infty_+K\bigr),
\]
then
\[
    \pi_m(A\{t\})
    \simeq
    \bigoplus_{n\geq0}A_m(B\Sigma_n)
\]
for every \(m\geq0\), while
\[
    \pi_0(A\{t\})
    \simeq
    \pi_0(A)[t].
\]
In particular, for \(A=HR\) with \(R\) a commutative ring,
\[
    \pi_m(HR\{t\})
    \simeq
    \bigoplus_{n\geq0}H_m(B\Sigma_n;R),
\]
recovering the degreewise calculation over a discrete spectral base. These
positive-degree groups can be nonzero: for example,
\[
    H_1(B\Sigma_2;\mathbb Z)
    \simeq
    \mathbb Z/2,
\]
so \(\pi_1(H\mathbb Z\{t\})\) contains a \(\mathbb Z/2\)-summand. They need
not be nonzero for every base: if \(R\) is a \(\mathbb Q\)-algebra, transfer
for the finite groups \(\Sigma_n\) gives
\[
    H_m(B\Sigma_n;R)=0
    \qquad(m>0),
\]
and the rank-one free commutative algebra is discrete. The formula therefore
displays exactly when the brave new affine line carries additional positive
spectral homotopy.
\end{proposition}

\begin{proof}
\leavevmode
\begin{enumerate}[label=\textup{(\arabic*)}]
    \item \textbf{Apply the relative free-algebra formula.}
    By the relative version of
    \Ref{prop:underlying-free-spectral-algebra},
    \[
        A\{t\}
        \simeq
        \bigoplus_{n\geq0}
        (A^{\otimes_A n})_{h\Sigma_n}
        \simeq
        \bigoplus_{n\geq0}(A)_{h\Sigma_n}.
    \]
    The \(\Sigma_n\)-action is trivial on the tensor power of the rank-one
    free \(A\)-module, so
    \[
        (A)_{h\Sigma_n}
        \simeq
        A\otimes\Sigma^\infty_+B\Sigma_n.
    \]

    \item \textbf{Take homotopy groups.}
    Homotopy groups commute with the displayed coproduct, giving
    \[
        \pi_m(A\{t\})
        \simeq
        \bigoplus_{n\geq0}A_m(B\Sigma_n).
    \]
    The degree-zero algebra identification is already
    \Ref{prop:pi-zero-relative-free-algebra} applied to the connective
    \(A\)-module \(A\):
    \[
        \pi_0(A\{t\})\simeq\pi_0(A)[t].
    \]
    When \(A=HR\), the generalized homology groups above are ordinary
    homology groups \(H_m(B\Sigma_n;R)\). Compare
    \cite{LurieHigherAlgebra,KhanBraveNewMotivicI}.
\end{enumerate}
\end{proof}

The two linear representation conventions are now separated and related
precisely: relative linear spaces represent functionals, section stacks
represent sections, and finite locally free coefficients make the latter
affine by duality. The final substantive step is retrospective rather than
foundational: having built the spectral geometry natively, we recover its
ordinary \(0\)-truncated sector as an interface with classical algebraic
geometry.

\section{Reconstructing the \texorpdfstring{$0$}{0}-truncated comparison}
\label{sec:zero-truncated-reconstruction}

The spectral geometry above was defined natively. Ordinary schemes can now be
recovered from it through the degree-zero reflection already constructed for
connective commutative ring spectra. This section extracts that \(0\)-truncated
interface: first on affines, then globally by Zariski gluing. The completed
spectral construction therefore yields the comparison as a consequence.

Let
\[
    \delta
    :
    \operatorname{AffSch}
    \hookrightarrow
    \Aff_{\Sp},
    \qquad
    \Spec(R)\longmapsto\Spec(HR),
\]
be the fully faithful embedding of ordinary affine schemes obtained from the
Eilenberg--Mac Lane functor.

\begin{proposition}[Restriction to \texorpdfstring{$0$}{0}-truncated affine tests]
\label{prop:zero-truncated-affine-reconstruction}
For every connective commutative ring spectrum \(A\), restriction of the
spectral functor of points along \(\delta\) is represented by
\(\Spec\pi_0(A)\): for every commutative ring \(R\),
\[
\begin{aligned}
    h_A(\Spec(HR))
    &\simeq
    \Map_{\CAlg(\Sp)}(A,HR)
    \\
    &\simeq
    \Map_{\operatorname{CRing}}(\pi_0(A),R).
\end{aligned}
\]
Moreover, for \(r\in R\), there is a canonical equivalence
\[
    HR[r^{-1}]
    \simeq
    H(R[r^{-1}]).
\]
Consequently, restriction along \(\delta\) carries the spectral principal
Zariski basis to the usual principal Zariski basis.
\end{proposition}

\[
\begin{tikzcd}[column sep=huge, row sep=large]
    t_0D_A^{\mathrm{sp}}(f)
        \arrow[r]
        \arrow[d, "\simeq"']
    &
    t_0\Spec(A)
        \arrow[d, "\simeq"]
    \\
    D_{\pi_0(A)}(f)
        \arrow[r]
    &
    \Spec\pi_0(A).
\end{tikzcd}
\]


\begin{proof}
\leavevmode
\begin{enumerate}[label=\textup{(\arabic*)}]
    \item \textbf{Restrict the affine functor of points.}
    The mapping-space equivalence is exactly the adjunction
    \(\pi_0\dashv H_{\mathrm{CRing}}\) of
    \Ref{prop:pi-zero-h-adjunction}. Ordinary affine Yoneda therefore
    identifies the restricted functor with \(\Spec\pi_0(A)\).

    \item \textbf{Compare principal localization.}
    By \Ref{prop:homotopy-groups-localization}, the localization
    \(HR[r^{-1}]\) is connective, has degree-zero homotopy
    \(R[r^{-1}]\), and has no positive homotopy groups. Hence it is
    canonically \(H(R[r^{-1}])\). Thus the principal covering families are
    preserved under restriction.
\end{enumerate}
\end{proof}

\begin{theorem}[Global \texorpdfstring{$0$}{0}-truncated reconstruction and adjunction]
\label{thm:global-zero-truncated-reconstruction}
The affine enhancement \(\delta\) extends by Zariski gluing to a functor
\[
    (-)^{\mathrm{disc}}:
    \operatorname{Sch}
    \longrightarrow
    \operatorname{Sch}^{\mathrm{sp}},
\]
characterized on affines by
\[
    (\Spec R)^{\mathrm{disc}}
    \simeq
    \Spec(HR).
\]
Restriction of the functor of points along \(\delta\) carries every spectral
scheme \(X\) to a Zariski sheaf represented by an ordinary scheme \(t_0X\),
and the two constructions form an adjunction
\[
\begin{tikzcd}[column sep=huge]
    \operatorname{Sch}
        \arrow[r, shift left=1.1ex, "{(-)^{\mathrm{disc}}}"]
    &
    \operatorname{Sch}^{\mathrm{sp}}
        \arrow[l, swap, shift left=1.1ex, "t_0"'].
\end{tikzcd}
\qquad
(-)^{\mathrm{disc}}\dashv t_0.
\]
It has the following properties.
\begin{enumerate}[label=(\roman*)]
    \item For every connective commutative ring spectrum \(A\),
    \[
        t_0\Spec(A)
        \simeq
        \Spec\pi_0(A).
    \]

    \item The adjunction unit
    \[
        S\longrightarrow t_0(S^{\mathrm{disc}})
    \]
    is an equivalence for every ordinary scheme \(S\). Hence
    \((-)^{\mathrm{disc}}\) is fully faithful.

    \item The right adjoint \(t_0\) preserves finite limits.

    \item Principal opens correspond under \(t_0\):
    \[
        t_0D_A^{\mathrm{sp}}(f)
        \simeq
        D_{\pi_0(A)}(f).
    \]
    Consequently the posets of open subschemes of \(X\) and of \(t_0X\) are
    canonically equivalent.

    \item The adjunction counit is a natural morphism
    \[
        \iota_X:
        (t_0X)^{\mathrm{disc}}
        \longrightarrow
        X,
    \]
    whose restriction to \(\Spec(A)\subseteq X\) is the morphism represented
    by
    \[
        A\longrightarrow H\pi_0(A).
    \]

    \item Under the standard connective spectral-scheme recognition theorem,
    affineness, quasi-compactness, quasi-separatedness, and the qcqs condition
    are detected by \(t_0\).
\end{enumerate}
The construction in this theorem is an interface with the \(0\)-truncated
sector; none of its detection statements is used in the native construction
of spectral schemes, coefficients, locality, or comprehension above.
\end{theorem}

\[
\begin{tikzcd}[column sep=huge, row sep=large]
    \operatorname{AffSch}
        \arrow[r, hook]
        \arrow[d, "\delta"']
    &
    \operatorname{Sch}
        \arrow[d, "{(-)^{\mathrm{disc}}}"]
    \\
    \Aff_{\Sp}
        \arrow[r, hook]
    &
    \operatorname{Sch}^{\mathrm{sp}}.
\end{tikzcd}
\]


\begin{proof}
\leavevmode
\begin{enumerate}[label=\textup{(\arabic*)}]
    \item \textbf{Choose an affine atlas of the ordinary scheme.}
    Let an ordinary scheme \(S\) have an affine Zariski atlas
    \(\coprod_i\Spec(R_i)\to S\).

    \item \textbf{Enhance principal intersections spectrally.}
    By
    \Ref{prop:zero-truncated-affine-reconstruction}, Eilenberg--Mac Lane
    enhancement carries principal intersections and their restriction maps to
    the corresponding spectral principal intersections.

    \item \textbf{Glue the discrete enhancement.}
    Spectral Zariski
    descent therefore glues the affines \(\Spec(HR_i)\) to a spectral scheme
    \(S^{\mathrm{disc}}\).

    \item \textbf{Remove dependence on the atlas.}
    Common refinement and mapping descent make the
    construction independent of the chosen atlas and functorial in \(S\).

    \item \textbf{Restrict a spectral affine atlas to degree zero.}
    Conversely, let
    \[
        \coprod_i\Spec(A_i)\longrightarrow X
    \]
    be an affine spectral Zariski atlas. Restriction along \(\delta\) carries
    each chart to \(\Spec\pi_0(A_i)\) and every principal spectral cover to
    the corresponding ordinary principal cover.

    \item \textbf{Recognize the ordinary scheme.}
    Since principal covers
    generate both Zariski topologies, the restricted functor is an ordinary
    Zariski sheaf with an effective affine open atlas and is therefore
    represented by an ordinary scheme.

    \item \textbf{Define \(t_0X\).}
    Define that scheme to be \(t_0X\).

    \item \textbf{Record the affine formula.}
    On affines this gives
    \(t_0\Spec(A)\simeq\Spec\pi_0(A)\).

    \item \textbf{Establish the affine mapping equivalence.}
    Let \(T=\Spec R\) be an ordinary affine scheme and let \(X\) be an
    arbitrary spectral scheme. By the functor-of-points definition and the
    construction of \(t_0X\), there are natural equivalences
    \[
    \begin{aligned}
        \Map_{\operatorname{Sch}^{\mathrm{sp}}}
        \bigl(T^{\mathrm{disc}},X\bigr)
        &\simeq
        X(\Spec(HR))
        \\
        &\simeq
        (t_0X)(\Spec R)
        \\
        &\simeq
        \Map_{\operatorname{Sch}}(T,t_0X).
    \end{aligned}
    \]

    \item \textbf{Identify its first two components.}
    The first equivalence is spectral Yoneda and the second is the defining
    restriction of the functor of points along \(\delta\).

    \item \textbf{Choose an affine atlas of a general ordinary scheme.}
    For a general
    ordinary scheme \(S\), choose an affine Zariski atlas
    \(U\to S\).

    \item \textbf{Apply mapping descent.}
    Its discrete enhancement is an effective affine open atlas
    \(U^{\mathrm{disc}}\to S^{\mathrm{disc}}\), and mapping descent on the
    two sides gives
    \[
    \begin{tikzcd}[column sep=huge]
        \Map_{\operatorname{Sch}^{\mathrm{sp}}}(S^{\mathrm{disc}},X)
            \arrow[r, "\simeq"]
        &
        \Map_{\operatorname{Sch}}(S,t_0X).
    \end{tikzcd}
    \]

    \item \textbf{Establish naturality.}
    Naturality of the affine equivalence on the \v{C}ech nerve makes this
    equivalence natural in both variables.

    \item \textbf{Conclude the global adjunction.}
    This is the asserted adjunction.

    \item \textbf{Identify the unit on affines.}
    For an ordinary affine scheme \(\Spec R\), the unit is the identity after
    the canonical identification
    \[
        t_0\Spec(HR)\simeq\Spec R.
    \]

    \item \textbf{Globalize the unit.}
    Zariski descent therefore makes the unit an equivalence for every ordinary
    scheme.

    \item \textbf{Deduce full faithfulness.}
    Thus discrete enhancement is fully faithful.

    \item \textbf{Identify the counit on affines.}
    The counit on
    \(\Spec A\) is contravariantly represented by the unit of the algebraic
    reflection
    \[
        A\longrightarrow H\pi_0(A),
    \]

    \item \textbf{Glue the counit.}
    These affine maps glue by naturality.

    \item \textbf{Use the right-adjoint property to preserve finite limits.}
    Finally, as a right adjoint,
    \(t_0\) preserves every limit that exists in spectral schemes, in
    particular the finite limits constructed in
    \Ref{cor:spectral-schemes-finite-limits}.

    \item \textbf{Identify principal open subobjects.}
    The equivalence
    \[
        HR[r^{-1}]\simeq H(R[r^{-1}])
    \]
    of \Ref{prop:zero-truncated-affine-reconstruction} identifies principal
    spectral opens with ordinary principal opens.

    \item \textbf{Pass to small unions.}
    Small unions of open
    substacks are constructed by \Ref{prop:unions-open-substacks} and commute
    with base change.

    \item \textbf{Glue the open-subobject correspondence.}
    Since principal opens form a basis on both sides, the
    principal correspondence glues to a canonical equivalence of open-subobject
    posets.

    \item \textbf{Invoke only the remaining recognition properties.}
    The standard connective spectral-scheme comparison identifies the functor
    just constructed with the usual \(0\)-truncation of spectral algebraic
    geometry, under which affineness, quasi-compactness, and
    quasi-separatedness are detected; see \cite{LurieSAG}. This invocation is
    used only for these recognition properties and imports no additional
    structure on \(X\).
\end{enumerate}
\end{proof}

\begin{corollary}[The truncation counit is a closed immersion]
\label{cor:zero-truncation-counit-closed}
For every spectral scheme \(X\), the counit
\[
    \iota_X:
    (t_0X)^{\mathrm{disc}}
    \longrightarrow
    X
\]
of
\(\Ref{thm:global-zero-truncated-reconstruction}\)
is a closed immersion. It induces a homeomorphism on underlying topological
spaces
\[
    \bigl|(t_0X)^{\mathrm{disc}}\bigr|
    \xrightarrow{\ \cong\ }
    |X|,
\]
and its \(0\)-truncation
\[
    t_0(\iota_X):
    t_0\bigl((t_0X)^{\mathrm{disc}}\bigr)
    \longrightarrow
    t_0X
\]
is an equivalence.
\end{corollary}

\begin{proof}
\leavevmode
\begin{enumerate}[label=\textup{(\arabic*)}]
    \item \textbf{Reduce closedness to an affine target.}
    The assertion that \(\iota_X\) is a closed immersion is local on the
    target. Let
    \[
        X=\Spec(A)
    \]
    be affine.

    \item \textbf{Identify the affine form of the counit.}
    By
    \Ref{thm:global-zero-truncated-reconstruction},
    \[
        t_0X\simeq\Spec\pi_0(A),
    \]
    and the restriction of the counit is
    \[
        \Spec(H\pi_0(A))
        \longrightarrow
        \Spec(A),
    \]
    represented contravariantly by the canonical morphism
    \[
        A\longrightarrow H\pi_0(A).
    \]

    \item \textbf{Compute degree-zero homotopy.}
    On degree-zero homotopy this induces the canonical isomorphism
    \[
        \pi_0(A)
        \xrightarrow{\ \cong\ }
        \pi_0(H\pi_0(A))
        \simeq
        \pi_0(A).
    \]

    \item \textbf{Deduce affine closedness.}
    In particular, it is surjective. Hence the displayed affine morphism is
    a closed immersion by
    \Ref{def:affine-closed-immersion-spectral}.

    \item \textbf{Globalize closedness.}
    Since closed immersions are
    defined affine-locally on the target, \(\iota_X\) is a closed immersion
    for every spectral scheme \(X\).

    \item \textbf{Underlying topological support.}
    On the same affine open,
    \[
        \bigl|\Spec(H\pi_0(A))\bigr|
        =
        \bigl|\Spec\pi_0(A)\bigr|
        =
        \bigl|\Spec(A)\bigr|.
    \]
    Under these identifications the map induced by \(\iota_X\) is the
    identity. These affine identifications are compatible on intersections
    and therefore glue to a homeomorphism
    \[
        \bigl|(t_0X)^{\mathrm{disc}}\bigr|
        \xrightarrow{\ \cong\ }
        |X|.
    \]

    \item \textbf{Apply the \(0\)-truncation functor.}
    On an affine open \(X=\Spec(A)\), applying \(t_0\) to the counit gives
    \[
        \Spec\pi_0(H\pi_0(A))
        \longrightarrow
        \Spec\pi_0(A).
    \]
    Since
    \[
        \pi_0(H\pi_0(A))
        \simeq
        \pi_0(A),
    \]
    this morphism is an equivalence. The affine equivalences agree on
    overlaps, so they glue to a global equivalence
    \[
        t_0\bigl((t_0X)^{\mathrm{disc}}\bigr)
        \xrightarrow{\simeq}
        t_0X.
    \]
\end{enumerate}
\end{proof}

\begin{convention}[Classical comparison notation]
\label{conv:classical-comparison-notation}
For compatibility with later chapters, when the \(0\)-truncated comparison is
used we write
\[
    X^{\mathrm{cl}}:=t_0X,
    \qquad
    f^{\mathrm{cl}}:=t_0(f).
\]
The notation is an abbreviation for the right-adjoint reconstruction above;
it does not reintroduce classical geometry into the native construction.
\end{convention}

The \(0\)-truncated comparison therefore closes the circle without altering
the native construction: ordinary schemes embed discretely, \(t_0\) recovers
their degree-zero geometry, and the truncation counit identifies the ordinary
support inside the spectral scheme. We can now summarize the chapter in terms
of the structures that remain available for later use.

\section{Chapter summary}
\label{sec:spectral-geometric-output}

The chapter has produced a coherent spectral geometry whose persistent
structure can be compressed into four interacting outputs.

First, the geometric contexts are spectral schemes built from connective
coordinates. The connective symmetric-monoidal category
\(\Sp_{\geq0}^\otimes\) supplies the semantic universe for the commutative
operadic theory, giving
\[
    \Alg_{\Comm^\otimes}(\Sp_{\geq0}^\otimes)
    \simeq
    \CAlg(\Sp)^{\mathrm{cn}}.
\]
Variance reversal produces the affine contexts
\[
    \Aff_{\Sp}
    =
    \bigl(\CAlg(\Sp)^{\mathrm{cn}}\bigr)^{\op},
\]
principal spectral localizations generate their Zariski topology, and Zariski
gluing produces \(\operatorname{Sch}^{\mathrm{sp}}\). The resulting category
has the finite-limit and open-subobject operations used throughout the later
geometry and agrees with the standard connective \(0\)-localic structured
model.

Second, every spectral context carries a fully stable coefficient fibre. On
an affine context \(\Spec(A)\) the fibre is \(\Mod_A\), and substitution
along \(\Spec(B)\to\Spec(A)\) is extension of scalars. Right Kan extension
and Zariski descent globalize these fibres to the coefficient system
\[
    \QCoh
    :
    \bigl(\operatorname{Sch}^{\mathrm{sp}}\bigr)^{\op}
    \longrightarrow
    \CAlg(\Pr^L_{\mathrm{st}}).
\]
Its pullbacks are coherent and strong symmetric monoidal, its right adjoints
are functorial, and the global fibres inherit the Postnikov structure needed
to formulate flatness and connective algebra without truncating the stable
coefficient theory.

Third, locality is detected through those coefficients and reconstructs both
geometry and algebra. Affine open immersions are flat, locally finitely
presented monomorphisms, and an affine family covers exactly when the
associated substitutions are jointly conservative. Finite principal covers
are descendable, so Amitsur and comonadic descent reconstruct representables,
modules, and commutative algebras from their restrictions. The geometric site
and the coefficient system are therefore linked by one reconstruction
mechanism rather than by two independent locality theories.

Fourth, relative spectrum returns connective coefficient algebra to geometry.
For every spectral scheme \(X\),
\[
    \operatorname{QCAlg}(X)^{\mathrm{cn},\op}
    \xrightarrow{\simeq}
    \bigl(\operatorname{Sch}^{\mathrm{sp}}_{/X}\bigr)_{\mathrm{aff}},
\]
and this equivalence is compatible with arbitrary substitution. Maps into the
resulting affine extension represent algebra terms
\(f^*\mathcal A\to\mathcal O_Y\), with the identity of the relative spectrum
corresponding to the universal variable. This is the chapter's comprehension
principle: connective relative algebra is represented by affine geometry over
the base.

The subsequent constructions are consequences and specializations of this
core. Closed immersions admit intrinsic open complements and
base-change-stable pushouts. Free commutative algebras give relative linear
spaces representing linear functionals; arbitrary coefficients have section
stacks; and finite locally free coefficients have affine section total spaces
by duality. The brave new affine spaces and \(\mathbb G_m\) are distinguished
instances. Finally, the adjunction
\[
    (-)^{\mathrm{disc}}
    \dashv
    t_0
\]
recovers ordinary algebraic geometry as the \(0\)-truncated interface, with
\(t_0\Spec(A)\simeq\Spec\pi_0(A)\) and the same underlying topological
support.

The dependency structure of the chapter is therefore
\[
    \boxed{
    \text{operadic theory}
    \longrightarrow
    \text{connective models}
    \longrightarrow
    \text{spectral contexts}
    \longrightarrow
    \text{stable substitution}
    \longrightarrow
    \text{Zariski descent}
    \longrightarrow
    \text{affine and linear comprehension}.}
\]
What remains active is the coexistence of two levels. Geometry is built from
connective commutative ring spectra, while coefficients remain fully stable.
Substitution transports the stable fibres coherently, locality detects and
reconstructs their geometric base, and relative spectrum returns connective
coefficient algebras to geometry. The later closed, linear, and classical
comparison constructions all operate inside this same indexed structure.
